% Rédiger une dissertation — Français 1re Tech — Cours signé OnBuch+
% Programme officiel MINESEC. Compiler avec : tectonic N04-rediger-dissertation.tex
\newcommand{\DOCMATIERE}{Français}
\newcommand{\DOCNIVEAU}{1re Tech}
\newcommand{\DOCMODULE}{Français – classe de Première technique}
\newcommand{\DOCLECON}{N04}
\newcommand{\DOCTITRE}{Rédiger une dissertation}
% OnBuch+ — préambule « cours signés » ENRICHI (compilé avec tectonic / XeLaTeX)
% Système visuel « Pop » d'OnBuch+ : fond crème (jamais blanc), bordures encre
% épaisses, ombres « dures » décalées, titres Archivo Black en capitales.
% Version MATHÉMATIQUES : graphes (pgfplots/tikz), boîtes pédagogiques
% (définition, propriété, méthode, attention, le savais-tu, exercice résolu,
% exercice, TP, bilan), page de garde et signature OnBuch+ ancrée en bas.
%
% Macros attendues dans main.tex AVANT \input{preamble} :
%   \DOCMATIERE  \DOCNIVEAU  \DOCMODULE  \DOCLECON (numéro)  \DOCTITRE
\documentclass[11pt]{article}

\usepackage{fontspec}
\usepackage[french]{babel}
\usepackage[a4paper,margin=2cm,top=3.4cm,bottom=2.8cm,headheight=1.4cm,footskip=1.3cm]{geometry}
\usepackage{amsmath,amssymb,amsfonts}
\usepackage{mathtools} % \xmapsto, \coloneqq... souvent utilisés par les modèles
\usepackage{cancel} % simplifications barrées (bilans de forces)
% \abs{z} (module, valeur absolue) et \norm{u} (norme), utilisés par les modèles
\providecommand{\abs}{}\let\abs\relax\DeclarePairedDelimiter{\abs}{\lvert}{\rvert}
\providecommand{\norm}{}\let\norm\relax\DeclarePairedDelimiter{\norm}{\lVert}{\rVert}
\usepackage[table]{xcolor} % [table] : fournit aussi \rowcolors (alternance de lignes)
\usepackage{pagecolor}
\usepackage{tikz}
\usetikzlibrary{arrows.meta,positioning,calc,shapes.geometric,shapes.misc,shapes.arrows,decorations.pathmorphing,decorations.pathreplacing,decorations.markings,patterns,shadows,mindmap,trees,fit,backgrounds,matrix,intersections,angles,quotes}
\usepackage{pgfplots}
\usepackage[version=4]{mhchem} % équations nucléaires \ce{^{235}_{92}U}
\usepackage[european,siunitx]{circuitikz} % schémas électriques
\pgfplotsset{compat=1.18}
\usepgfplotslibrary{fillbetween,groupplots} % aires entre courbes (calcul intégral)
\usepackage{enumitem}
\usepackage{fancyhdr}
% [most] charge toutes les bibliothèques tcolorbox (listings, minted, raster,
% documentation...) et gonfle inutilement la mémoire TeX sur un document long
% (~300 boîtes) : on ne charge que ce qui est réellement utilisé.
\usepackage[skins,breakable]{tcolorbox}
\usepackage{graphicx}
\usepackage{colortbl}
\usepackage{booktabs}
\usepackage{tabularx}
\usepackage{multirow}
\usepackage{diagbox} % case d'en-tête diagonale des tableaux à double entrée
% Colonnes extensibles centrée / gauche / droite (C, L, R), souvent utilisées par les modèles
\newcolumntype{C}{>{\centering\arraybackslash}X}
\newcolumntype{L}{>{\raggedright\arraybackslash}X}
\newcolumntype{R}{>{\raggedleft\arraybackslash}X}
\usepackage{array}
\usepackage{multicol}
\usepackage{siunitx}
% parse-numbers=false : accepte aussi \SI{2,5 \times 10^{-3}}{...}
\sisetup{locale=FR,output-decimal-marker={,},group-minimum-digits=4,parse-numbers=false}
\DeclareSIUnit\ohm{\ensuremath{\symup{\Omega}}} % U+2126 absent de la police
\DeclareSIUnit\division{div} % réglages d'oscilloscope (ms/div, V/div)
\DeclareSIUnit\tr{tr} % tours (vitesse de rotation, ex. \SI{1500}{\tr\per\minute})
\DeclareSIUnit\molar{M} % concentration molaire (ex. \SI{3}{\molar}, \SI{1}{\micro\molar})
\DeclareSIUnit\an{an} % année (le modèle confond parfois avec \year, un compteur TeX) — ex. \SI{5}{\kilo\meter\per\an}
\DeclareSIUnit\FCFA{FCFA} % franc CFA (ex. \SI{500}{\FCFA\per\kilo\gram})
\DeclareSIUnit\degreeBrix{\textdegree Brix} % teneur en sucre d'un sirop/jus (réfractométrie)
\DeclareSIUnit\month{mois} % le modèle utilise parfois \month, un compteur TeX, comme unité de durée
\DeclareSIUnit\microliter{\micro\liter} % le modèle écrit parfois \microliter en un seul token
\DeclareSIUnit\million{millions} % dénombrement (ex. \SI{15}{\million\per\mL} spermatozoïdes)
\usepackage{qrcode}
% \degree : commande fréquemment utilisée par les modèles pour un angle ou une
% température en dehors de \SI{}{} (non fournie par les paquets chargés).
\providecommand{\degree}{^{\circ}}
% \qed (fin de démonstration), utilisé par les modèles sans amsthm
\providecommand{\qed}{\hfill$\square$}
% \barre : double barre des valeurs interdites dans un tableau de variations (tkz-tab)
\providecommand{\barre}{\ensuremath{\|}}
% environnement proof (amsthm non chargé) utilisé par les modèles
\ifdefined\proof\else\newenvironment{proof}[1][Démonstration]{\par\noindent\textit{#1.}\ }{\qed\par}\fi
\usepackage{lastpage}
\usepackage{hyperref}
\hypersetup{hidelinks}

% ============================================================
% Palette « Pop » OnBuch+ — mêmes hex que la classe P (plus_ui.dart)
% ============================================================
\definecolor{poporange}{HTML}{F59321}
\definecolor{popdark}{HTML}{E07A0C}
\definecolor{popcream}{HTML}{FFF6E8}
\definecolor{popink}{HTML}{1C1714}
\definecolor{poppurple}{HTML}{7A5AE0}
\definecolor{popgreen}{HTML}{2E9E6A}
\definecolor{poppink}{HTML}{E0457B}
\definecolor{popblue}{HTML}{2F7DE1}
\definecolor{popgold}{HTML}{C98A12}
\definecolor{popmuted}{HTML}{8A7F76}
\definecolor{popline}{HTML}{EADFD2}
\definecolor{popgray}{HTML}{8A7F76}
\colorlet{popgrey}{popgray}
% Alias tolérants (le modèle nomme parfois les couleurs différemment)
\colorlet{popwhite}{white}
\colorlet{popblack}{popink}
\colorlet{popred}{poppink}
\colorlet{popyellow}{popgold}
% Teintes claires (fonds de tableaux/encadrés) — le modèle nomme parfois
% les couleurs avec un suffixe L (ex. popblueL) sans les avoir définies.
\colorlet{poporangeL}{poporange!20}
\colorlet{popdarkL}{popdark!20}
\colorlet{poppurpleL}{poppurple!20}
\colorlet{popgreenL}{popgreen!20}
\colorlet{poppinkL}{poppink!20}
\colorlet{popblueL}{popblue!20}
\colorlet{popgoldL}{popgold!20}
\colorlet{popredL}{popred!20}
\colorlet{popgrayL}{popgray!20}
% Teintes claires pour fonds de tableaux / zones
\colorlet{popblueL}{popblue!10!white}
\colorlet{poporangeL}{poporange!14!white}
\colorlet{popgreenL}{popgreen!12!white}
\colorlet{poppurpleL}{poppurple!12!white}
\colorlet{poppinkL}{poppink!12!white}
\colorlet{popgoldL}{popgold!14!white}

\pagecolor{popcream}

% ============================================================
% Typographie
% ============================================================
\defaultfontfeatures{Ligatures=TeX}
\setmainfont{PlusJakartaSans-Medium.ttf}[Path=../../../fonts/,BoldFont=PlusJakartaSans-Bold.ttf]
% Maths en sans-serif (Fira Math) avec chiffres et lettres droites de Plus
% Jakarta, pour que les formules se fondent dans le texte.
\usepackage[math-style=ISO,bold-style=ISO]{unicode-math}
\setmathfont{FiraMath-Regular.otf}[Path=../../../fonts/]
\setmathfont{PlusJakartaSans-Medium.ttf}[Path=../../../fonts/,range={up/num,up/{Latin,latin},\mathup}]
\usepackage{icomma} % virgule décimale sans espace en mode math
% Caractères mathématiques Unicode absents des polices de texte (Plus Jakarta,
% Archivo Black) : rendus par leur équivalent mathématique, en texte comme en
% formule — sinon ils s'affichent en carré vide (ex. « Arithmétique dans ℤ »).
\usepackage{newunicodechar}
\newunicodechar{ℝ}{\ensuremath{\mathbb{R}}}
\newunicodechar{ℤ}{\ensuremath{\mathbb{Z}}}
\newunicodechar{ℂ}{\ensuremath{\mathbb{C}}}
\newunicodechar{ℕ}{\ensuremath{\mathbb{N}}}
\newunicodechar{ℚ}{\ensuremath{\mathbb{Q}}}
\newunicodechar{𝔻}{\ensuremath{\mathbb{D}}}
\newunicodechar{α}{\ensuremath{\alpha}}
\newunicodechar{β}{\ensuremath{\beta}}
\newunicodechar{γ}{\ensuremath{\gamma}}
\newunicodechar{δ}{\ensuremath{\delta}}
\newunicodechar{ε}{\ensuremath{\varepsilon}}
\newunicodechar{θ}{\ensuremath{\theta}}
\newunicodechar{λ}{\ensuremath{\lambda}}
\newunicodechar{μ}{\ensuremath{\mu}}
\newunicodechar{σ}{\ensuremath{\sigma}}
\newunicodechar{τ}{\ensuremath{\tau}}
\newunicodechar{φ}{\ensuremath{\varphi}}
\newunicodechar{ψ}{\ensuremath{\psi}}
\newunicodechar{ω}{\ensuremath{\omega}}
\newunicodechar{Σ}{\ensuremath{\Sigma}}
% majuscules produites par \MakeUppercase dans les titres (α-aminés -> Α)
\newunicodechar{Α}{\ensuremath{\alpha}}
\newunicodechar{Β}{\ensuremath{\beta}}
\newunicodechar{Γ}{\ensuremath{\Gamma}}
\newunicodechar{Δ}{\ensuremath{\Delta}}
\newunicodechar{Ε}{\ensuremath{\varepsilon}}
\newunicodechar{Θ}{\ensuremath{\Theta}}
\newunicodechar{Λ}{\ensuremath{\Lambda}}
\newunicodechar{Μ}{\ensuremath{\mu}}
\newunicodechar{Τ}{\ensuremath{\tau}}
\newunicodechar{Φ}{\ensuremath{\Phi}}
\newunicodechar{Ψ}{\ensuremath{\Psi}}
\newunicodechar{Ω}{\ensuremath{\Omega}}
\newunicodechar{↦}{\ensuremath{\mapsto}}
\newunicodechar{∈}{\ensuremath{\in}}
\newunicodechar{∉}{\ensuremath{\notin}}
\newunicodechar{∧}{\ensuremath{\wedge}}
\newunicodechar{⇔}{\ensuremath{\Leftrightarrow}}
\newunicodechar{⟺}{\ensuremath{\Longleftrightarrow}}
\newunicodechar{⇒}{\ensuremath{\Rightarrow}}
\newunicodechar{⟹}{\ensuremath{\Longrightarrow}}
\newunicodechar{∘}{\ensuremath{\circ}}
\newunicodechar{∪}{\ensuremath{\cup}}
\newunicodechar{∩}{\ensuremath{\cap}}
\newunicodechar{≡}{\ensuremath{\equiv}}
\newunicodechar{⊂}{\ensuremath{\subset}}
\newunicodechar{⊥}{\ensuremath{\perp}}
\newunicodechar{∃}{\ensuremath{\exists}}
\newunicodechar{∀}{\ensuremath{\forall}}
\newunicodechar{∛}{\ensuremath{\sqrt[3]{\,}}}
\newunicodechar{∜}{\ensuremath{\sqrt[4]{\,}}}
\newunicodechar{⃗}{\ensuremath{{}^{\to}}}
\newunicodechar{ⁿ}{\ifmmode^{n}\else\textsuperscript{\ensuremath{n}}\fi}
\newunicodechar{ˣ}{\ifmmode^{x}\else\textsuperscript{\ensuremath{x}}\fi}
\newunicodechar{ᵇ}{\ifmmode^{b}\else\textsuperscript{\ensuremath{b}}\fi}
\newunicodechar{ʳ}{\ifmmode^{r}\else\textsuperscript{\ensuremath{r}}\fi}
\newunicodechar{ᵘ}{\ifmmode^{u}\else\textsuperscript{\ensuremath{u}}\fi}
\newunicodechar{ʸ}{\ifmmode^{y}\else\textsuperscript{\ensuremath{y}}\fi}
\newunicodechar{ᵃ}{\ifmmode^{a}\else\textsuperscript{\ensuremath{a}}\fi}
\newunicodechar{ᵉ}{\ifmmode^{e}\else\textsuperscript{\ensuremath{e}}\fi}
\newunicodechar{ᵏ}{\ifmmode^{k}\else\textsuperscript{\ensuremath{k}}\fi}
\newunicodechar{ᵖ}{\ifmmode^{p}\else\textsuperscript{\ensuremath{p}}\fi}
\newunicodechar{ᴺ}{\ifmmode^{N}\else\textsuperscript{\ensuremath{N}}\fi}
\newunicodechar{ᶜ}{\ifmmode^{c}\else\textsuperscript{\ensuremath{c}}\fi}
\newunicodechar{ʰ}{\ifmmode^{h}\else\textsuperscript{\ensuremath{h}}\fi}
\newunicodechar{ᵠ}{\ifmmode^{q}\else\textsuperscript{\ensuremath{q}}\fi}
\newunicodechar{ₙ}{\ifmmode_{n}\else\textsubscript{\ensuremath{n}}\fi}
\newunicodechar{ₐ}{\ifmmode_{a}\else\textsubscript{\ensuremath{a}}\fi}
\newunicodechar{ᵢ}{\ifmmode_{i}\else\textsubscript{\ensuremath{i}}\fi}
\newunicodechar{ᵣ}{\ifmmode_{r}\else\textsubscript{\ensuremath{r}}\fi}
\newunicodechar{ₖ}{\ifmmode_{k}\else\textsubscript{\ensuremath{k}}\fi}
\newunicodechar{ᵦ}{\ifmmode_{\beta}\else\textsubscript{\ensuremath{\beta}}\fi}
\newunicodechar{ₓ}{\ifmmode_{x}\else\textsubscript{\ensuremath{x}}\fi}
\newfontfamily\popdisplay{ArchivoBlack-Regular.ttf}[Path=../../../fonts/]
\newfontfamily\poplogo{SpaceGrotesk-Bold.ttf}[Path=../../../fonts/]
\newfontfamily\popsemi{PlusJakartaSans-SemiBold.ttf}[Path=../../../fonts/]
\newfontfamily\popxbold{PlusJakartaSans-ExtraBold.ttf}[Path=../../../fonts/]
\newfontfamily\popxboldls{PlusJakartaSans-ExtraBold.ttf}[Path=../../../fonts/,LetterSpace=3.0]

\setlist[enumerate]{leftmargin=1.7em,itemsep=5pt,topsep=4pt}
\setlist[itemize]{leftmargin=1.7em,itemsep=4pt,topsep=4pt,label={\color{poporange}\textbullet}}
\setlength{\parindent}{0pt}
\setlength{\parskip}{5pt}
\renewcommand{\arraystretch}{1.25}

% pgfplots : style Pop par défaut
\pgfplotsset{
  every axis/.append style={axis line style={popink,line width=1pt},tick style={popink},
    grid style={popline,line width=0.5pt},label style={font=\small},tick label style={font=\footnotesize}},
  popcurve/.style={poporange,line width=1.8pt,no markers,smooth},
}
% Même style disponible dans un \draw TikZ simple (hors pgfplots)
\tikzset{popcurve/.style={poporange,line width=1.8pt}}
\tikzset{popgrid/.style={popline,very thin}}
\colorlet{popcurve}{poporange} % le nom du style est parfois utilisé comme couleur

% ============================================================
% Pill / squiggle
% ============================================================
\newcommand{\poppill}[3]{%
  \tikz[baseline=-0.55ex]\node[draw=popink,line width=1.2pt,rounded corners=2.6pt,%
    fill=#2,inner xsep=6.5pt,inner ysep=3pt]{%
    \popxboldls\fontsize{7}{7}\selectfont\color{#3}\MakeUppercase{#1}};%
}
\newcommand{\popsquiggle}[1]{%
  \tikz[baseline=-0.4ex]{
    \draw[#1,line width=2.6pt,line cap=round]
      plot [smooth] coordinates {(0,0.09) (0.34,0.01) (0.68,0.21) (1.02,0.01) (1.36,0.10)};
  }%
}
% Mot-clé surligné (marqueur orange sous le texte)
\newcommand{\cle}[1]{{\popxbold\color{popdark}#1}}

% ============================================================
% En-tête / pied
% ============================================================
\pagestyle{fancy}
\fancyhf{}
\renewcommand{\headrulewidth}{0pt}
\renewcommand{\footrulewidth}{0pt}
\lhead{\raisebox{-0.15ex}{\poplogo\fontsize{13}{13}\selectfont\color{popink}OnBuch\color{poporange}+}}
\rhead{\poppill{\DOCMATIERE\ \textbullet\ \DOCNIVEAU\ \textbullet\ Leçon \DOCLECON}{white}{popink}}
\lfoot{\popxbold\fontsize{7.6}{7.6}\selectfont\color{popmuted}\MakeUppercase{Cours signé OnBuch+}\ \textbullet\ {\popsemi\DOCTITRE}}
\rfoot{\tikz[baseline=-0.6ex]\node[rounded corners=8.5pt,fill=popink,minimum height=17pt,minimum width=17pt,inner xsep=5pt,inner sep=0]{\popxbold\fontsize{8}{8}\selectfont\color{popcream}\thepage\,/\,\pageref*{LastPage}};}

% ============================================================
% Titres
% ============================================================
\newcommand{\chaptitre}[2]{%
  \begin{tcolorbox}[enhanced,colback=poporange,colframe=popink,boxrule=2.6pt,arc=11pt,
      drop shadow=popink,left=15pt,right=15pt,top=12pt,bottom=11pt,before skip=3pt,after skip=18pt]
    \poppill{#2}{popink}{popcream}\par\vspace{9pt}
    {\raggedright\hyphenpenalty=10000\exhyphenpenalty=10000\popdisplay\fontsize{22}{24}\selectfont\color{popcream}\MakeUppercase{#1}\par}
  \end{tcolorbox}
}
% Section numérotée : pastille orange avec le numéro + titre Archivo
\newcounter{popsec}
\newcommand{\coursec}[1]{%
  \stepcounter{popsec}\setcounter{popsub}{0}%
  \par\vspace{16pt}\noindent
  \begin{minipage}{\linewidth}
  \tikz[baseline=-0.7ex]\node[draw=popink,line width=1.6pt,fill=poporange,rounded corners=4pt,minimum size=22pt,inner sep=0]{\popdisplay\fontsize{12}{12}\selectfont\color{popcream}\thepopsec};%
  \hspace{8pt}{\popdisplay\fontsize{15}{17}\selectfont\color{popink}\MakeUppercase{#1}}\par
  \vspace{4pt}\noindent{\color{poporange}\rule{36pt}{3.6pt}}
  \end{minipage}\par\nopagebreak\vspace{8pt}
}
\newcounter{popsub}
\newcommand{\courssub}[1]{%
  \stepcounter{popsub}%
  \par\vspace{9pt}\noindent{\popxbold\fontsize{11.5}{13}\selectfont\color{popink}{\color{poporange}\thepopsec.\thepopsub}\hspace{6pt}#1}\par\nopagebreak\vspace{4pt}
}

% ============================================================
% Boîtes pédagogiques — même recette : fond blanc, bordure épaisse colorée,
% ombre dure encre, badge plein. Titre optionnel [#1].
% ============================================================
\tcbset{popbase/.style={breakable,enhanced,colback=white,boxrule=2.3pt,arc=9pt,
  shadow={3pt}{-3pt}{0pt}{popink},left=14pt,right=14pt,top=11pt,bottom=11pt,before skip=11pt,after skip=15pt,
  fontupper=\popsemi\fontsize{10.6}{14.5}\selectfont\color{popink},
  fontlower=\popsemi\fontsize{10.6}{14.5}\selectfont\color{popink}}}
% Coche dessinée (le glyphe ✓ n'existe pas dans les polices du document)
\newcommand{\popcheck}[1]{\tikz[baseline=-0.5ex]\draw[#1,line width=1.6pt,line cap=round,line join=round](-0.13,0.0)--(-0.04,-0.09)--(0.13,0.1);}
\providecommand{\coche}{\popcheck{popgreen}}
\providecommand{\resultat}[1]{\par\smallskip\noindent\fcolorbox{popgreen}{popgreenL}{\parbox{\dimexpr\linewidth-2\fboxsep-2\fboxrule\relax}{\textbf{Résultat.} #1}}\par\smallskip}
\newcommand{\popboxhead}[3]{\poppill{#1}{#2}{white}\ifx\relax#3\relax\else\hspace{6pt}{\popxbold\fontsize{10.6}{12}\selectfont\color{popink}#3}\fi\par\vspace{7pt}}

\newtcolorbox{aretenir}[1][]{popbase,colframe=popgreen,before upper={\popboxhead{À retenir}{popgreen}{#1}}}
\newtcolorbox{exemplebox}[1][]{popbase,colframe=poporange,before upper={\popboxhead{Exemple}{poporange}{#1}}}
\newtcolorbox{definition}[1][]{popbase,colframe=popblue,before upper={\popboxhead{Définition}{popblue}{#1}}}
\newtcolorbox{propriete}[1][]{popbase,colframe=poppurple,before upper={\popboxhead{Propriété}{poppurple}{#1}}}
\newtcolorbox{methode}[1][]{popbase,colframe=popgold,before upper={\popboxhead{Méthode}{popgold}{#1}}}
\newtcolorbox{attention}[1][]{popbase,colframe=poppink,before upper={\popboxhead{Attention}{poppink}{#1}}}
\newtcolorbox{experience}[1][]{popbase,colframe=popblue,colback=popblueL,before upper={\popboxhead{Expérience}{popblue}{#1}}}
% « Le savais-tu ? » : carte violette pleine (contexte, culture, Cameroun)
\newtcolorbox{savaistu}[1][]{popbase,colback=poppurple,colframe=popink,
  fontupper=\popsemi\fontsize{10.4}{14.2}\selectfont\color{white},
  before upper={\poppill{Le savais-tu\,?}{white}{poppurple}\ifx\relax#1\relax\else\hspace{6pt}{\popxbold\color{white}#1}\fi\par\vspace{7pt}}}
% Exercice résolu : énoncé en haut, \tcblower puis la solution sur fond crème
\newcounter{popexo}\newcounter{popexr}
\newtcolorbox{exoresolu}[1][]{popbase,colframe=poporange,colbacklower=popcream,
  segmentation style={popink,line width=1.2pt,dashed},
  before upper={\stepcounter{popexr}\popboxhead{Exercice résolu \thepopexr}{poporange}{#1}},
  before lower={\poppill{Solution}{popink}{popcream}\par\vspace{6pt}}}
% Exercice (non résolu en place) — la correction est dans la partie Corrigés
\newtcolorbox{exercice}[1][]{popbase,colframe=popink,
  before upper={\stepcounter{popexo}\popboxhead{Exercice \thepopexo}{popink}{#1}}}
% Corrigé
\newtcolorbox{corrige}[1][]{popbase,colframe=popgreen,colback=popgreenL,
  before upper={\popboxhead{Corrigé}{popgreen}{#1}}}
% Objectifs / prérequis / bilan
\newtcolorbox{objectifs}{popbase,colframe=popink,colback=poporangeL,
  before upper={\popboxhead{Objectifs de la leçon}{popink}{}}}
\newtcolorbox{prerequis}{popbase,colframe=popink,colback=popblueL,
  before upper={\popboxhead{Prérequis}{popblue}{}}}
\newtcolorbox{bilan}{popbase,colframe=popink,colback=white,boxrule=2.8pt,
  before upper={\popboxhead{Fiche bilan}{poporange}{L'essentiel à connaître pour l'examen}}}
% Figure encadrée avec légende
\newtcolorbox{popfigure}[1][]{enhanced,breakable=false,colback=white,colframe=popline,boxrule=1.4pt,arc=8pt,
  left=10pt,right=10pt,top=10pt,bottom=8pt,before skip=10pt,after skip=12pt,halign=center,
  lower separated=false,halign lower=center,
  fontlower=\popsemi\fontsize{9}{11.5}\selectfont\color{popmuted},
  #1}
\newcommand{\legende}[1]{\par\vspace{6pt}{\popsemi\fontsize{9}{11.5}\selectfont\color{popmuted}#1}}

% ============================================================
% Signature OnBuch+
% ============================================================
\newcommand{\poplogoinline}[1]{{\poplogo\fontsize{#1}{#1}\selectfont\color{popink}OnBuch\color{poporange}+}}
\newcommand{\signatureonbuch}{%
  \begin{tcolorbox}[enhanced,colback=popink,colframe=popink,boxrule=2.2pt,arc=10pt,
      drop shadow=poporange,left=14pt,right=14pt,top=10pt,bottom=10pt,before skip=0pt,after skip=0pt]
    \tikz[baseline=-0.7ex]\node[circle,fill=popgreen,draw=popcream,line width=1.3pt,minimum size=22pt,inner sep=0]{\popcheck{white}};%
    \hspace{9pt}\begin{minipage}[c]{0.62\linewidth}
      {\popxbold\fontsize{12}{13}\selectfont\color{popcream}Signé\ \textbullet\ {\poplogo OnBuch\color{poporange}+}}\par\vspace{2pt}
      {\raggedright\popsemi\fontsize{8}{10.5}\selectfont\color{popline}\MakeUppercase{Équipe pédagogique OnBuch+}\\\MakeUppercase{Conforme au programme officiel MINESEC}\par}
    \end{minipage}\hfill
    {\poplogo\fontsize{22}{22}\selectfont\color{popcream}OnBuch\color{poporange}+}
  \end{tcolorbox}%
}

% ============================================================
% Page de garde
% #1 titre · #2 matière · #3 niveau · #4 module · #5 n° leçon · #6 durée ·
% #7 description / objectifs (texte ou itemize)
% ============================================================
\newcommand{\pagegarde}[7]{%
  \thispagestyle{empty}
  \newgeometry{a4paper,margin=1.7cm,top=1.6cm,bottom=1.4cm}%
  \begin{tikzpicture}[remember picture,overlay]
    \fill[poporange!18] ([xshift=-3.2cm,yshift=-3.6cm]current page.north east) circle (4.6cm);
    \fill[poppurple!14] ([xshift=2.4cm,yshift=7.6cm]current page.south west) circle (3.8cm);
  \end{tikzpicture}%
  {\poplogo\fontsize{30}{30}\selectfont\color{popink}OnBuch\color{poporange}+}\hfill
  \raisebox{0.6ex}{\poppill{Cours signé \textbullet\ Premium}{popink}{popcream}}\par
  \vspace{1pt}\popsquiggle{poppurple}\par
  \vspace{8pt}
  \begin{tcolorbox}[enhanced,colback=poporange,colframe=popink,boxrule=2.8pt,arc=13pt,
      drop shadow=popink,left=18pt,right=18pt,top=12pt,bottom=13pt,after skip=10pt]
    \poppill{#2 \textbullet\ #3}{popink}{popcream}\hspace{5pt}\poppill{#4}{white}{popink}\par\vspace{8pt}
    {\popdisplay\fontsize{14}{14}\selectfont\color{popink}LEÇON #5}\par\vspace{4pt}
    {\raggedright\hyphenpenalty=10000\exhyphenpenalty=10000\popdisplay\fontsize{25}{27}\selectfont\color{popcream}\MakeUppercase{#1}\par}
    \vspace{8pt}
    {\popxbold\fontsize{9.5}{9.5}\selectfont\color{popink}\MakeUppercase{Durée conseillée : #6}}
  \end{tcolorbox}
  \begin{tcolorbox}[enhanced,colback=white,colframe=popink,boxrule=2.2pt,arc=10pt,
      drop shadow=popink,left=16pt,right=16pt,top=10pt,bottom=10pt,after skip=8pt]
    \poppill{Dans cette leçon}{poppurple}{white}\par\vspace{5pt}
    \popsemi\fontsize{9.3}{12.6}\selectfont\color{popink}#7
  \end{tcolorbox}
  \noindent\begin{minipage}[t]{0.487\linewidth}
    \begin{tcolorbox}[enhanced,colback=popgreen,colframe=popink,boxrule=2.2pt,arc=10pt,
        drop shadow=popink,left=11pt,right=11pt,top=10pt,bottom=10pt,height=3.1cm,valign=center]
      \tcbox[colback=white,colframe=popink,boxrule=1.3pt,arc=5pt,boxsep=0pt,nobeforeafter,
        left=3pt,right=3pt,top=3pt,bottom=3pt,valign=center]{\qrcode[height=1.9cm]{https://play.google.com/store/apps/details?id=com.luvvix.onbuch&hl=fr}}%
      \hspace{8pt}\begin{minipage}[c]{0.5\linewidth}
        {\popdisplay\fontsize{11}{12.5}\selectfont\color{white}\MakeUppercase{Télécharge OnBuch}}\par\vspace{4pt}
        {\raggedright\popsemi\fontsize{8.4}{11}\selectfont\color{white}Gratuit \textbullet\ Google Play\\Tuteur IA, annales, concours, résultats\par}
      \end{minipage}
    \end{tcolorbox}
  \end{minipage}\hfill
  \begin{minipage}[t]{0.487\linewidth}
    \begin{tcolorbox}[enhanced,colback=poppurple,colframe=popink,boxrule=2.2pt,arc=10pt,
        drop shadow=popink,left=11pt,right=11pt,top=10pt,bottom=10pt,height=3.1cm,valign=center]
      \tikz[baseline=-0.7ex]\node[circle,fill=white,draw=popink,line width=1.3pt,minimum size=1.6cm,inner sep=0]{\popdisplay\fontsize{24}{24}\selectfont\color{poppurple}+};%
      \hspace{8pt}\begin{minipage}[c]{0.6\linewidth}
        {\popdisplay\fontsize{11}{12.5}\selectfont\color{white}\MakeUppercase{Passe à OnBuch+}}\par\vspace{4pt}
        {\raggedright\popsemi\fontsize{8.4}{11}\selectfont\color{white}Tous les cours signés, Tuteur IA illimité, fiches PDF — onglet OnBuch+ de l'app\par}
      \end{minipage}
    \end{tcolorbox}
  \end{minipage}
  \vspace{8pt}
  \popcontenu
  \vfill
  \signatureonbuch
  \restoregeometry
  \newpage
}

% Rangée « Ce cours contient » (page de garde)
\newcommand{\poptile}[3]{% #1 couleur #2 gros texte #3 légende
  \begin{tcolorbox}[enhanced,colback=#1,colframe=popink,boxrule=2pt,arc=9pt,shadow={2.5pt}{-2.5pt}{0pt}{popink},
      left=6pt,right=6pt,top=9pt,bottom=9pt,halign=center,valign=center,height=1.75cm,width=\linewidth,nobeforeafter]
    {\popdisplay\fontsize{12.5}{13}\selectfont\color{white}\MakeUppercase{#2}}\par\vspace{3pt}
    {\centering\popsemi\fontsize{7.8}{9.5}\selectfont\color{white}#3\par}
  \end{tcolorbox}}
\newcommand{\popcontenu}{%
  \par\noindent{\popxboldls\fontsize{7.5}{7.5}\selectfont\color{popmuted}CE COURS SIGNÉ CONTIENT}\par\vspace{7pt}
  \noindent\begin{minipage}[t]{0.235\linewidth}\poptile{popblue}{Cours}{complet, illustré, pas à pas}\end{minipage}\hfill
  \begin{minipage}[t]{0.235\linewidth}\poptile{poporange}{Méthodes}{et exercices résolus}\end{minipage}\hfill
  \begin{minipage}[t]{0.235\linewidth}\poptile{poppink}{Exercices}{type examen + corrigés}\end{minipage}\hfill
  \begin{minipage}[t]{0.235\linewidth}\poptile{popgreen}{Fiche bilan}{l'essentiel à retenir}\end{minipage}\par}

% Sommaire
\newtcolorbox{sommaire}{enhanced,colback=white,colframe=popink,boxrule=2.2pt,arc=10pt,
  drop shadow=popink,left=18pt,right=18pt,top=13pt,bottom=13pt,before skip=6pt,after skip=20pt,
  before upper={\poppill{Sommaire}{poppurple}{white}\par\vspace{9pt}\popsemi\fontsize{10.6}{14.5}\selectfont\color{popink}}}

% Signature de fin de cours — poussée tout en bas de la dernière page
\newcommand{\findecours}{%
  \par\vfill\nopagebreak
  \begin{center}\popsquiggle{poporange}\end{center}\vspace{4pt}
  \signatureonbuch
}
\AtBeginDocument{\def\llbracket{\mathopen{[\mkern-3.5mu[}}\def\rrbracket{\mathclose{]\mkern-3.5mu]}}} % intervalles d'entiers (glyphe ⟦ absent de la police)
% Glyphes absents de Fira Math : redéfinis à partir de glyphes disponibles
\AtBeginDocument{%
  \def\setminus{\mathbin{\backslash}}\let\smallsetminus\setminus
  \def\vdots{\mathord{\vbox{\baselineskip3.2pt\lineskiplimit0pt\kern4pt\hbox{.}\hbox{.}\hbox{.}}}}%
  \def\ddots{\mathinner{\mkern1mu\raise6.5pt\hbox{.}\mkern2mu\raise3.6pt\hbox{.}\mkern2mu\raise0.7pt\hbox{.}\mkern1mu}}%
  \def\triangle{\mathord{\scalebox{0.9}{$\symup{\Delta}$}}}%
  \def\nabla{\mathord{\raisebox{\depth}{\scalebox{1}[-1]{$\symup{\Delta}$}}}}%
  \def\checkmark{\popcheck{popdark}}%
  \def\star{\mathbin{\ast}}\def\bigstar{\mathord{\ast}}%
  \def\diamond{\mathbin{\scalebox{0.6}{$\Diamond$}}}%
  \def\bigcup{\mathop{\mathchoice{\raisebox{-0.3em}{\scalebox{1.7}{$\cup$}}}{\scalebox{1.25}{$\cup$}}{\cup}{\cup}}\displaylimits}%
  \def\bigcap{\mathop{\mathchoice{\raisebox{-0.3em}{\scalebox{1.7}{$\cap$}}}{\scalebox{1.25}{$\cap$}}{\cap}{\cap}}\displaylimits}%
  \def\lesseqqgtr{\mathrel{\lessgtr}}\def\bigoplus{\mathop{\oplus}\displaylimits}%
}
\providecommand{\ssi}{si et seulement si} % abréviation fréquente
\AtBeginDocument{\providecommand{\wideparen}{\overparen}} % arc de cercle

\begin{document}

\pagegarde{Rédiger une dissertation}{Français}{1re Tech}{Français – classe de Première technique}{N04}{22 séances (fiche de progression harmonisée)}{Cette leçon outille l'élève de Première technique pour rédiger une dissertation complète : analyse du sujet, introduction, développement avec citations et transitions, conclusion. Elle s'appuie sur la lecture méthodique du Lion et la perle et consolide les notions de langue (types de phrases, dénotation/connotation, énonciation, champ lexical et sémantique, texte descriptif).
\vspace{4pt}
\begin{multicols}{2}\small
\begin{itemize}
\item À la fin de cette leçon, je sais analyser un sujet de dissertation (dégager le thème, la problématique et le plan).
\item Je sais rédiger une introduction complète (amorce, présentation du sujet, problématique, annonce du plan) et l'améliorer.
\item Je sais rédiger un paragraphe argumentatif en insérant correctement des citations tirées du Lion et la perle.
\item Je sais rédiger des transitions qui articulent les parties du développement.
\item Je sais rédiger une conclusion (bilan, ouverture) et l'améliorer.
\item Je sais distinguer dénotation et connotation et exploiter les champs lexical et sémantique dans mes écrits.
\end{itemize}\end{multicols}}

\begin{sommaire}
\begin{multicols}{2}
\begin{enumerate}
\item Lecture méthodique du Lion et la perle (1) et les types de phrases
\item Analyser un sujet et rédiger l'introduction
\item Lecture méthodique (2), dénotation/connotation et le paragraphe argumentatif
\item Lecture méthodique (3), énonciation et conclusion
\item Champ lexical, champ sémantique et texte descriptif
\item Intégration : rédiger une dissertation complète
\item Activité d'intégration
\item Méthodes et exercices résolus
\item Exercices
\item Corrigés des exercices
\item Fiche bilan
\end{enumerate}
\end{multicols}
\end{sommaire}

\begin{objectifs}
\begin{itemize}
\item À la fin de cette leçon, je sais analyser un sujet de dissertation (dégager le thème, la problématique et le plan).
\item Je sais rédiger une introduction complète (amorce, présentation du sujet, problématique, annonce du plan) et l'améliorer.
\item Je sais rédiger un paragraphe argumentatif en insérant correctement des citations tirées du Lion et la perle.
\item Je sais rédiger des transitions qui articulent les parties du développement.
\item Je sais rédiger une conclusion (bilan, ouverture) et l'améliorer.
\item Je sais distinguer dénotation et connotation et exploiter les champs lexical et sémantique dans mes écrits.
\item Je sais identifier et employer les types de phrases obligatoires et les marques de l'énonciation.
\item Je sais reconnaître les éléments constitutifs du texte descriptif et les mobiliser dans une production écrite.
\item Je sais justifier ma démarche de rédaction et appliquer ces compétences à des situations concrètes de communication.
\end{itemize}
\end{objectifs}

\begin{prerequis}
\begin{itemize}
\item La lecture analytique d'un texte littéraire (classe de Seconde)
\item Les classes grammaticales et la fonction des mots dans la phrase
\item La notion de paragraphe et de texte argumentatif (Seconde)
\item La connaissance générale de l'œuvre Le Lion et la perle (début de lecture en classe)
\item Les connecteurs logiques de base (d'abord, ensuite, cependant, donc...)
\end{itemize}
\end{prerequis}

\newpage

\coursec{Lecture méthodique du Lion et la perle (1) et les types de phrases}

\begin{exemplebox}[Situation d'accroche : deux copies, deux destins]
À l'issue du concours d'entrée dans une grande école de Yaoundé, deux candidats ont traité le même sujet de dissertation. Le premier a obtenu 6/20, le second 17/20. Pourtant, tous deux connaissaient l'œuvre au programme. La différence ? Le second savait analyser un sujet, construire une introduction percutante, articuler ses paragraphes avec des citations et conclure avec élégance. Comment rédiger une dissertation qui convainc un correcteur ? C'est tout l'art que cette leçon va vous apprendre, en vous appuyant sur \emph{Le Lion et la perle} de Wole Soyinka.
\end{exemplebox}

\textbf{Problématique de la leçon :} comment lire méthodiquement une œuvre dramatique et comment maîtriser les types de phrase pour construire une argumentation claire et efficace ?

\courssub{Présentation de l'œuvre : Wole Soyinka et \emph{Le Lion et la perle}}

Avant de lire une œuvre, il faut savoir \emph{qui} l'a écrite, \emph{quand} et \emph{pourquoi}. C'est la première étape de toute lecture méthodique : situer le texte dans son contexte.

\textbf{Wole Soyinka} est un dramaturge, poète et essayiste nigérian, né en 1934 à Abeokuta, dans l'ouest du Nigeria. Formé au Nigeria puis en Angleterre, il est profondément marqué par la culture yoruba, dont il puise les mythes, les rites et la langue théâtrale. Engagé politiquement, il a connu la prison pour ses positions contre la guerre du Biafra. En 1986, il reçoit le prix Nobel de littérature.

\begin{savaistu}[Un Africain Nobel]
Wole Soyinka est le premier Africain à avoir reçu le prix Nobel de littérature, en 1986. Le jury a salué un écrivain « qui, dans une large perspective culturelle et avec des nuances poétiques, façonne le drame de l'existence ». Cette consécration mondiale montre que la littérature africaine, longtemps marginalisée, occupe désormais une place centrale dans les lettres universelles.
\end{savaistu}

\emph{Le Lion et la perle} (\emph{The Lion and the Jewel}, 1959, publié en 1963) est une comédie en trois actes intitulés « Matin », « Midi » et « Nuit ». L'action se déroule dans le village fictif d'\textbf{Ilujinle}, au Nigeria. Le conflit central de la pièce oppose deux visions du monde :

\begin{itemize}
\item la \cle{tradition}, incarnée par \textbf{Baroka}, le vieux chef du village, surnommé « le Lion », polygame, rusé, attaché aux coutumes ;
\item la \cle{modernité}, incarnée par \textbf{Lakunle}, le jeune instituteur, qui rêve de transformer le village à l'image de la ville et refuse les coutumes qu'il juge « barbares ».
\end{itemize}

Entre eux se trouve \textbf{Sidi}, la plus belle jeune fille du village, « la perle », dont les deux hommes veulent faire leur épouse. \textbf{Sadiku}, la première épouse de Baroka, joue le rôle d'entremetteuse. La pièce pose ainsi une question qui dépasse le simple triangle amoureux : que faire de la tradition face à la modernité ? Et quelle place la femme occupe-t-elle dans ce combat d'hommes ?

\begin{attention}[Ne pas caricaturer les personnages]
Soyinka ne fait pas de Lakunle un héros du progrès ni de Baroka un simple « vieux réactionnaire ». Lakunle est souvent ridicule : il récite des mots anglais sans les comprendre et refuse de payer la dot par pur calcul. Baroka, lui, est intelligent et fin stratège. Dans votre dissertation, évitez donc les jugements simplistes : analysez la complexité des personnages.
\end{attention}

\courssub{Lecture méthodique n° 1 : l'acte I, « Matin »}

La lecture méthodique d'une pièce de théâtre consiste à observer, dans l'ordre : la \cle{didascalie} (indications scéniques), les \cle{personnages} présents, la \cle{situation initiale} et les premiers \cle{conflits}.

\textbf{Ce que nous apprenons dans « Matin » :}

\begin{enumerate}
\item \textbf{L'exposition des personnages.} Sidi porte l'eau près de l'école ; Lakunle la courtise en la sermonnant sur ses « mauvaises » habitudes (porter des charges sur la tête, exposer ses épaules). Il lui promet un mariage moderne, sans dot.
\item \textbf{L'élément perturbateur.} Un photographe étranger est passé au village et a publié des photos de Sidi dans un magazine. Sa beauté est désormais connue de tous : elle se prend pour une « perle » précieuse.
\item \textbf{L'annonce du conflit.} Baroka, le chef, a vu les photos et veut épouser Sidi. Sadiku est chargée de porter la demande. La rivalité entre le Lion et l'instituteur est déclarée.
\end{enumerate}

\begin{popfigure}
\begin{center}
\begin{tikzpicture}[
  perso/.style={rectangle, rounded corners=3pt, draw=popdark, fill=popblueL, thick, minimum width=2.6cm, minimum height=0.9cm, align=center, font=\small},
  rel/.style={-{Stealth[length=2.5mm]}, thick, popdark},
  node distance=1.4cm]
\node[perso, fill=popgoldL, align=center] (baroka){\textbf{Baroka}\\ le chef, « le Lion »};
\node[perso, fill=poppinkL, below left=1.6cm and -0.5cm of baroka, align=center] (sidi){\textbf{Sidi}\\ « la perle »};
\node[perso, below right=1.6cm and -0.5cm of baroka, align=center] (lakunle){\textbf{Lakunle}\\ l'instituteur};
\node[perso, fill=popgreenL, above right=1.4cm and 0.2cm of baroka, align=center] (sadiku){\textbf{Sadiku}\\ première épouse};
\draw[rel, poppurple] (baroka) -- node[above left, font=\footnotesize, align=center]{veut épouser\\ (tradition)} (sidi);
\draw[rel, popblue] (lakunle) -- node[above right, font=\footnotesize, align=center]{courtise\\ (modernité)} (sidi);
\draw[rel, popgreen!60!black] (sadiku) -- node[right, font=\footnotesize, align=center]{entremetteuse\\ de Baroka} (baroka);
\draw[rel, dashed, poporange] (lakunle) to[bend left=15] node[below, font=\footnotesize]{rivalité} (baroka);
\end{tikzpicture}
\end{center}
\legende{Arbre des personnages du \emph{Lion et la perle} : Sidi est au centre du conflit qui oppose Baroka (tradition) et Lakunle (modernité) ; Sadiku sert d'intermédiaire.}
\end{popfigure}

\textbf{Relevé des thèmes majeurs.} Dès le premier acte, quatre thèmes se dessinent et pourront alimenter vos sujets de dissertation :

\begin{center}
\begin{tabularx}{\linewidth}{|l|X|}
\hline
\rowcolor{popblueL} \textbf{Thème} & \textbf{Comment il apparaît dans « Matin »} \\
\hline
La tradition & La dot, la polygamie, l'autorité du chef, les coutumes du village. \\
\hline
La modernité & L'école, la langue anglaise, le magazine, le refus de la dot par Lakunle. \\
\hline
Le statut de la femme & Sidi est convoitée comme un trophée ; peut-elle choisir librement ? \\
\hline
Le pouvoir & Baroka règne sur Ilujinle ; la beauté de Sidi est aussi un pouvoir. \\
\hline
\end{tabularx}
\end{center}

\courssub{Qu'est-ce qu'une dissertation ?}

\begin{definition}[La dissertation]
La \cle{dissertation} est un exercice écrit qui traite un sujet de manière \textbf{ordonnée} et \textbf{argumentée}. Elle comprend trois parties : une \textbf{introduction} (qui présente le sujet et pose la problématique), un \textbf{développement} (qui expose les arguments en plusieurs paragraphes) et une \textbf{conclusion} (qui répond à la problématique et ouvre éventuellement le débat).
\end{definition}

Au Probatoire et au Baccalauréat technique, la dissertation porte souvent sur l'œuvre au programme. Exemple de sujet : \emph{« Le conflit entre tradition et modernité dans Le Lion et la perle »}. Pour le traiter, il faut d'abord maîtriser l'outil de base de toute argumentation : la phrase. Or toutes les phrases ne font pas le même travail.

\courssub{Morphosyntaxe : les types de phrase}

Observez ces quatre phrases, qui parlent toutes de Sidi :

\begin{enumerate}
\item Sidi refuse la demande de Baroka.
\item Sidi refuse-t-elle la demande de Baroka ?
\item Refuse la demande de Baroka !
\item Comme Sidi refuse fièrement la demande de Baroka !
\end{enumerate}

Le contenu est le même, mais l'\textbf{intention du locuteur} change : affirmer, questionner, ordonner, s'émouvoir. C'est ce qu'on appelle le \cle{type de phrase}.

\begin{definition}[Les quatre types de phrase]
\begin{itemize}
\item La phrase \cle{déclarative} : elle donne une information, affirme un fait. Elle se termine par un point.
\item La phrase \cle{interrogative} : elle pose une question. Elle se termine par un point d'interrogation.
\item La phrase \cle{impérative} : elle exprime un ordre, un conseil, une prière. Le verbe est souvent à l'impératif ; elle se termine par un point ou un point d'exclamation.
\item La phrase \cle{exclamative} : elle exprime une émotion (admiration, colère, surprise). Elle se termine par un point d'exclamation.
\end{itemize}
\end{definition}

\begin{methode}[Repérer le type d'une phrase en trois observations]
\begin{enumerate}
\item \textbf{Observer la ponctuation finale} : point (déclarative), point d'interrogation (interrogative), point d'exclamation (impérative ou exclamative).
\item \textbf{Observer l'ordre des mots} : inversion du sujet ou mot interrogatif (\emph{que, qui, comment, pourquoi, est-ce que}) $\rightarrow$ interrogative ; verbe à l'impératif sans sujet exprimé $\rightarrow$ impérative.
\item \textbf{Observer les marques modales} : adverbes exclamatifs (\emph{comme, que, combien}), tournures emphatiques $\rightarrow$ exclamative. En cas de doute, demandez-vous : le locuteur affirme-t-il, questionne-t-il, ordonne-t-il ou s'émeut-il ?
\end{enumerate}
\end{methode}

\begin{exemplebox}[Deux phrases, deux types]
\begin{itemize}
\item \emph{« Sidi est-elle libre de choisir son époux ? »} : phrase \textbf{interrogative} (inversion du sujet, point d'interrogation) — elle pose un problème.
\item \emph{« Sidi choisit son époux. »} : phrase \textbf{déclarative} (ordre sujet-verbe, point final) — elle affirme un fait.
\end{itemize}
\end{exemplebox}

\begin{popfigure}
\begin{center}
\begin{tikzpicture}
\node[font=\small, align=center] at (0,0){
\begin{tabular}{|l|l|l|}
\hline
\rowcolor{popblueL} \textbf{Type de phrase} & \textbf{Marques formelles} & \textbf{Fonction dans la dissertation} \\
\hline
Déclarative & point final ; ordre sujet-verbe & affirmer une thèse, un argument, une analyse \\
\hline
Interrogative & « ? » ; inversion, \emph{est-ce que}, mot en \emph{qu-} & formuler la problématique dans l'introduction \\
\hline
Impérative & « ! » ou « . » ; verbe à l'impératif & rare : à éviter dans une dissertation \\
\hline
Exclamative & « ! » ; \emph{comme, que, combien} & à manier avec prudence (effet oratoire) \\
\hline
\end{tabular}};
\end{tikzpicture}
\end{center}
\legende{Tableau récapitulatif : chaque type de phrase a des marques formelles et une fonction précise dans le devoir.}
\end{popfigure}

\courssub{Les types de phrase au service de la dissertation}

Dans une dissertation, deux types de phrase sont \textbf{obligatoires} :

\begin{itemize}
\item La phrase \textbf{interrogative} sert à formuler la \cle{problématique} à la fin de l'introduction : \emph{« Comment Soyinka met-il en scène le conflit entre tradition et modernité ? »} C'est elle qui annonce le plan du devoir.
\item La phrase \textbf{déclarative} est la phrase de l'argumentation : elle énonce la thèse, les arguments et les analyses. Un développement est essentiellement une suite de phrases déclaratives bien articulées.
\end{itemize}

Il faut aussi distinguer le \textbf{type} de phrase de sa \textbf{forme} :

\begin{itemize}
\item \textbf{Forme affirmative / négative} : \emph{« Lakunle respecte les coutumes »} (affirmative) / \emph{« Lakunle ne respecte pas les coutumes »} (négative). La négation peut renforcer une opposition dans l'argumentation.
\item \textbf{Forme active / passive} : \emph{« Baroka épouse Sidi »} (active : le sujet fait l'action) / \emph{« Sidi est épousée par Baroka »} (passive : le sujet subit l'action). La voix passive met en avant celui qui \emph{subit} : écrire \emph{« Sidi est traitée comme un trophée »} souligne le statut de la femme dans la pièce.
\end{itemize}

\begin{attention}[Erreur fréquente : confondre type et forme]
Le \textbf{type} de phrase (déclarative, interrogative, impérative, exclamative) dépend de l'\textbf{énonciation}, c'est-à-dire de l'intention du locuteur. La \textbf{forme} (affirmative/négative, active/passive) dépend de la construction grammaticale. Une phrase déclarative peut être négative (\emph{« Sidi ne refuse pas. »}) et une phrase interrogative peut être affirmative (\emph{« Sidi refuse-t-elle ? »}). Ne dites jamais « phrase négative » quand on vous demande le \emph{type} de la phrase.
\end{attention}

\begin{exoresolu}[Identifier et transformer les types de phrase]
\textbf{Énoncé.} Voici trois phrases inspirées de l'acte I. 1) Indiquez le type de chaque phrase en justifiant. 2) Transformez la phrase a) en phrase interrogative, puis la phrase b) en phrase déclarative négative.
\begin{itemize}
\item[a)] Lakunle refuse de payer la dot.
\item[b)] Sidi acceptera-t-elle la demande de Baroka ?
\item[c)] Comme ce village est attaché à ses coutumes !
\end{itemize}
\tcblower
\textbf{Solution détaillée.}
\begin{enumerate}
\item \textbf{Identification.}
\begin{itemize}
\item a) Phrase \textbf{déclarative} : point final, ordre sujet-verbe-complément ; le locuteur affirme un fait.
\item b) Phrase \textbf{interrogative} : point d'interrogation et inversion du sujet (\emph{Sidi acceptera-t-elle}) ; le locuteur pose une question.
\item c) Phrase \textbf{exclamative} : point d'exclamation et adverbe exclamatif \emph{comme} ; le locuteur exprime son admiration.
\end{itemize}
\item \textbf{Transformations.}
\begin{itemize}
\item a) $\rightarrow$ interrogative : \emph{« Lakunle refuse-t-il de payer la dot ? »} (inversion du sujet avec le \emph{t} euphonique, point d'interrogation). Autre possibilité : \emph{« Est-ce que Lakunle refuse de payer la dot ? »}
\item b) $\rightarrow$ déclarative négative : \emph{« Sidi n'acceptera pas la demande de Baroka. »} (ordre sujet-verbe rétabli, négation \emph{ne\ldots pas}, point final).
\end{itemize}
\end{enumerate}
\textbf{Vérification :} la transformation change le type ou la forme, mais le sens de fond (le refus, l'acceptation) reste le même.
\end{exoresolu}

\begin{exercice}[À vous de jouer]
Relevez dans un extrait de l'acte I (« Matin ») deux phrases déclaratives et une phrase interrogative prononcées par Lakunle. Pour chacune : 1) justifiez le type par les marques formelles ; 2) transformez chaque phrase déclarative en phrase exclamative, puis à la forme négative ; 3) dites quel effet de sens produit la voix passive dans : \emph{« Sidi est admirée par tout le village »}.
\end{exercice}

\begin{corrige}[Exercice]
Exemple de réponse attendue. 1) \emph{« Vous portez des charges sur votre tête comme une femme du village. »} : déclarative (point final, ordre sujet-verbe). \emph{« Ne voyez-vous pas que ces coutumes sont dépassées ? »} : interrogative (point d'interrogation, inversion avec \emph{ne\ldots pas}). 2) Exclamative : \emph{« Comme vous portez fièrement ces charges sur votre tête ! »} ; négative : \emph{« Vous ne portez pas de charges sur votre tête. »} 3) La voix passive dans \emph{« Sidi est admirée par tout le village »} place Sidi en position de sujet qui \emph{subit} le regard des autres : elle devient un objet d'admiration, ce qui annonce son statut de « perle » convoitée, privée de libre choix.
\end{corrige}

\begin{aretenir}[L'essentiel de la section]
\begin{itemize}
\item \emph{Le Lion et la perle} (Wole Soyinka, Nobel 1986) met en scène à Ilujinle le conflit entre Baroka (tradition) et Lakunle (modernité) autour de Sidi ; l'acte I, « Matin », expose les personnages et déclenche la rivalité.
\item Les quatre thèmes majeurs : tradition, modernité, statut de la femme, pouvoir.
\item Les quatre types de phrase : déclarative, interrogative, impérative, exclamative ; on les repère par la ponctuation, l'ordre des mots et les marques modales.
\item Dans la dissertation : la phrase interrogative pour la problématique, la phrase déclarative pour l'argumentation ; ne pas confondre type (énonciation) et forme (affirmative/négative, active/passive).
\end{itemize}
\end{aretenir}

\coursec{Analyser un sujet et rédiger l'introduction}

Dans la section précédente, nous avons lu méthodiquement le début du \emph{Lion et la perle} de Wole Soyinka et appris à reconnaître les types de phrases (déclarative, interrogative, impérative, exclamative). Cette lecture nous a donné de la matière : des personnages, des thèmes, des citations. Mais au Probatoire et au Baccalauréat technique, avoir des connaissances ne suffit pas : il faut savoir \cle{répondre exactement à la question posée}. Or toute dissertation commence par deux gestes décisifs : \textbf{analyser le sujet} et \textbf{rédiger l'introduction}. C'est l'objet de cette section. Un correcteur, en lisant vos dix premières lignes, sait déjà si vous avez compris le sujet. Apprenons donc à ne rien laisser au hasard.

\courssub{Analyser le sujet : les quatre étapes}

\begin{attention}[Le piège du hors-sujet]
La faute la plus grave en dissertation n'est pas la faute d'orthographe : c'est le \cle{hors-sujet}. Un élève qui rédige trois pages brillantes... à côté de la question posée obtient une note très faible. L'analyse du sujet, au brouillon, prend 15 à 20 minutes : c'est un investissement, pas une perte de temps.
\end{attention}

\textbf{Étape 1 : lire le sujet plusieurs fois et définir les mots-clés.}\par
Lisez le sujet au moins trois fois, lentement. Soulignez ensuite les \cle{mots-clés}, c'est-à-dire les mots qui portent le sens : les noms, les verbes importants, les notions abstraites. Puis définissez chacun d'eux au brouillon, avec vos propres mots, comme si vous deviez l'expliquer à un camarade.

\begin{exemplebox}[Définir les mots-clés]
Sujet : \emph{Le conflit tradition/modernité dans Le Lion et la perle.}
\begin{itemize}
\item \textbf{Tradition} : ensemble des coutumes, des valeurs et des pratiques transmises par les ancêtres (ici, le village d'Ilujinle, le mariage coutumier, la dot).
\item \textbf{Modernité} : mode de vie et idées venus de l'Occident (école, progrès technique, nouvelles mœurs), incarnés par Lakunle, l'instituteur.
\item \textbf{Conflit} : opposition, affrontement entre deux forces. Le sujet demande donc de montrer \emph{comment} ces deux mondes s'affrontent dans la pièce.
\end{itemize}
\end{exemplebox}

\textbf{Étape 2 : dégager le thème et les limites du sujet.}\par
Le \cle{thème} est le grand sujet général dont parle le texte. Les \cle{limites} précisent ce qu'il faut traiter... et ce qu'il faut exclure. Demandez-vous : de quoi parle le sujet ? Sur quelle œuvre ? Sur quelle période ? Qu'est-ce qui n'est \emph{pas} demandé ?

\begin{exemplebox}[Thème et limites]
Pour notre sujet, le thème est \textbf{l'affrontement entre coutume africaine et modernité occidentale}. Les limites : on traite \emph{uniquement} de \emph{Le Lion et la perle} (pas d'autres pièces de Soyinka), et on s'intéresse au \emph{conflit}, pas à une simple description du village. Raconter toute l'histoire de Sidi serait donc hors de propos.
\end{exemplebox}

\textbf{Étape 3 : formuler la problématique.}\par
Une fois les mots définis et le thème cerné, il faut transformer le sujet en une \textbf{question} qui guidera toute la réflexion.

\begin{definition}[La problématique]
La \cle{problématique} est la \textbf{question centrale} que pose le sujet et à laquelle toute la dissertation devra répondre. Elle oriente la réflexion du début à la fin : chaque paragraphe du développement doit apporter un élément de réponse à cette question. Elle se formule avec un point d'interrogation, souvent introduit par \emph{dans quelle mesure}, \emph{comment}, \emph{pourquoi}.
\end{definition}

\begin{exemplebox}[Formuler une problématique]
Sujet : \emph{Le conflit tradition/modernité dans Le Lion et la perle.}
Problématiques possibles :
\begin{itemize}
\item Comment Soyinka met-il en scène l'affrontement entre la tradition et la modernité ?
\item Dans quelle mesure ce conflit révèle-t-il les tensions de l'Afrique contemporaine ?
\end{itemize}
\end{exemplebox}

\textbf{Étape 4 : mobiliser ses connaissances et choisir un plan.}\par
Faites un \cle{brainstorming} (remue-méninges) au brouillon : notez en vrac tout ce que vous savez en lien avec la problématique — personnages, scènes, citations, idées. Puis triez et regroupez ces idées en deux ou trois grandes parties : c'est le \cle{plan}. Trois types de plans sont au programme :

\begin{center}
\begin{tabularx}{\linewidth}{|l|X|X|}
\hline
\rowcolor{popblueL} \textbf{Type de plan} & \textbf{Principe} & \textbf{Quand l'utiliser} \\
\hline
Thématique & On examine les différents aspects d'un thème (aspect 1, aspect 2, aspect 3). & Sujets descriptifs ou d'explication. \\
\hline
Dialectique & Thèse / antithèse / synthèse : on confronte deux opinions opposées avant de dépasser l'opposition. & Sujets de débat, questions polémiques. \\
\hline
Analytique & Constat / causes / conséquences (ou solutions). & Sujets portant sur un phénomène ou un problème. \\
\hline
\end{tabularx}
\end{center}

\begin{methode}[Analyser un sujet en 4 étapes]
\begin{enumerate}
\item \textbf{Mots-clés} : lire le sujet trois fois, souligner et définir chaque terme important.
\item \textbf{Thème et limites} : dégager le grand thème, préciser ce qu'il faut traiter et ce qu'il faut exclure.
\item \textbf{Problématique} : transformer le sujet en une question centrale.
\item \textbf{Plan} : brainstorming des idées, puis regroupement en 2 ou 3 parties selon un plan thématique, dialectique ou analytique.
\end{enumerate}
\end{methode}

\begin{popfigure}
\begin{center}
\begin{tikzpicture}[node distance=0.9cm,
  etape/.style={rectangle, rounded corners, draw=popblue, fill=popblueL, thick, text width=9.5cm, align=center, minimum height=1cm},
  fleche/.style={-{Stealth[length=3mm]}, thick, popdark}]
\node[etape] (e1) {\textbf{1. Mots-clés} : lire, souligner, définir chaque terme};
\node[etape, below=of e1] (e2) {\textbf{2. Thème et limites} : quoi traiter, quoi exclure};
\node[etape, below=of e2] (e3) {\textbf{3. Problématique} : la question centrale};
\node[etape, below=of e3] (e4) {\textbf{4. Plan} : brainstorming puis 2 ou 3 parties};
\draw[fleche] (e1) -- (e2);
\draw[fleche] (e2) -- (e3);
\draw[fleche] (e3) -- (e4);
\end{tikzpicture}
\end{center}
\legende{Organigramme des quatre étapes de l'analyse d'un sujet de dissertation. Chaque étape prépare la suivante : on ne peut pas formuler la problématique sans avoir défini les mots-clés.}
\end{popfigure}

\courssub{Rédiger l'introduction : les quatre mouvements}

L'introduction est la \textbf{porte d'entrée} de votre devoir. Imaginez que vous recevez un visiteur : vous ne lui jetez pas le repas à la figure ; vous l'accueillez, vous lui présentez la maison, puis vous l'invitez à table. L'introduction fonctionne de même : elle conduit progressivement le lecteur du général vers le sujet précis. C'est un mouvement \textbf{en entonnoir}.

\begin{definition}[Les quatre mouvements de l'introduction]
Une introduction de dissertation comporte obligatoirement quatre mouvements, dans cet ordre :
\begin{enumerate}
\item L'\cle{amorce} (ou accroche) : une phrase générale qui met le lecteur dans le bain.
\item La \cle{présentation du sujet} : on situe l'œuvre, l'auteur, et on reformule le sujet avec ses propres mots.
\item La \cle{problématique} : la question centrale, annoncée clairement.
\item L'\cle{annonce du plan} : on présente les grandes parties du développement.
\end{enumerate}
\end{definition}

\begin{popfigure}
\begin{center}
\begin{tikzpicture}
\draw[thick, popblue, fill=popblueL] (-4.5,4) -- (4.5,4) -- (2.6,2) -- (-2.6,2) -- cycle;
\draw[thick, popblue, fill=popgreenL] (-2.6,2) -- (2.6,2) -- (1.5,0.2) -- (-1.5,0.2) -- cycle;
\draw[thick, popblue, fill=poporangeL] (-1.5,0.2) -- (1.5,0.2) -- (0.8,-1.4) -- (-0.8,-1.4) -- cycle;
\draw[thick, popblue, fill=poppinkL] (-0.8,-1.4) -- (0.8,-1.4) -- (0.3,-2.6) -- (-0.3,-2.6) -- cycle;
\node at (0,3) {\textbf{1. Amorce} (fait général)};
\node at (0,1.1) {\textbf{2. Présentation du sujet}};
\node at (0,-0.6) {\textbf{3. Problématique}};
\node at (0,-2) {\small \textbf{4. Plan}};
\end{tikzpicture}
\end{center}
\legende{Schéma en entonnoir de l'introduction : on part du plus général (amorce) pour aboutir au plus précis (annonce du plan).}
\end{popfigure}

\courssub{Les techniques d'amorce}

L'amorce est la première phrase : elle doit être juste, générale et en lien avec le sujet. Trois techniques sont à votre disposition :

\begin{itemize}
\item \textbf{Le fait général} : une vérité large, observable par tous. Exemple : « Toute société connaît, à un moment de son histoire, la confrontation entre ses coutumes et les idées nouvelles. »
\item \textbf{La citation} : une parole d'auteur célèbre, à condition de la connaître exactement et de nommer son auteur. Une citation fausse ou approximative dessert le candidat.
\item \textbf{Le contexte historique ou littéraire} : situer l'œuvre dans son époque ou son mouvement. Exemple : « Publiée en 1963, à l'époque des indépendances africaines, la pièce de Soyinka s'inscrit dans le débat sur l'avenir culturel du continent. »
\end{itemize}

\begin{attention}[Erreurs fréquentes dans l'introduction]
\begin{itemize}
\item Annoncer le plan avec « je vais parler de... » ou « dans ma première partie je vais dire... ». Le « je » est à bannir. Préférez : « Nous verrons d'abord que..., puis nous montrerons que... »
\item Commencer par une amorce trop éloignée (« Depuis la nuit des temps, l'homme... ») : restez proche du sujet.
\item Recopier le sujet mot pour mot : il faut le \emph{reformuler}.
\item Rédiger une introduction trop longue : \textbf{8 à 12 lignes maximum}.
\end{itemize}
\end{attention}

\begin{methode}[Rédiger l'introduction pas à pas]
\begin{enumerate}
\item Rédiger l'amorce (1 à 2 phrases) : fait général, citation ou contexte.
\item Présenter le sujet (2 à 3 phrases) : nommer l'auteur, l'œuvre, le genre, et reformuler le sujet.
\item Poser la problématique (1 phrase interrogative).
\item Annoncer le plan (1 à 2 phrases) avec des connecteurs : « nous verrons d'abord que..., ensuite que..., enfin que... »
\item Vérifier la longueur totale : 8 à 12 lignes.
\end{enumerate}
\end{methode}

\courssub{Production guidée : une introduction complète}

\begin{exemplebox}[Introduction rédigée]
Sujet : \emph{Le conflit tradition/modernité dans Le Lion et la perle.}\par\smallskip
« L'Afrique des indépendances est partagée entre la fidélité aux coutumes ancestrales et l'attrait du progrès occidental. C'est précisément ce débat que met en scène le dramaturge nigérian Wole Soyinka, prix Nobel de littérature, dans sa pièce \emph{Le Lion et la perle}, publiée en 1963. Dans le village d'Ilujinle, la belle Sidi est courtisée par deux hommes que tout oppose : Lakunle, le jeune instituteur moderniste qui méprise la dot, et Baroka, le vieux chef traditionnel, rusé et polygame. À travers cette rivalité amoureuse, c'est tout un monde qui s'affronte. Comment Soyinka met-il en scène l'affrontement entre la tradition et la modernité, et quelle vision de l'Afrique ce conflit révèle-t-il ? Nous verrons d'abord que la pièce oppose deux personnages et deux visions du monde ; nous montrerons ensuite que ce conflit, loin d'être manichéen, débouche sur une réflexion nuancée sur l'identité africaine. »
\end{exemplebox}

Observez ce modèle : l'amorce est un fait général ; la présentation nomme l'auteur, l'œuvre et les personnages ; la problématique est une vraie question ; le plan est annoncé avec « nous verrons d'abord... nous montrerons ensuite... »

\begin{savaistu}[Le poids de l'introduction à l'examen]
Au Probatoire et au Baccalauréat technique, l'introduction représente environ \textbf{un quart de la note} de dissertation. Les correcteurs y vérifient quatre points précis : la présence de l'amorce, la présentation du sujet, la problématique et l'annonce du plan. Une introduction absente ou bâclée pénalise l'ensemble de la copie, même si le développement est bon !
\end{savaistu}

\courssub{Améliorer son introduction : la grille de relecture}

Rédiger, c'est d'abord \textbf{réécrire}. Après avoir produit votre introduction au brouillon, relisez-la à l'aide de la grille suivante et corrigez chaque point faible avant de recopier.

\begin{center}
\begin{tabularx}{\linewidth}{|l|X|X|}
\hline
\rowcolor{popblueL} \textbf{Critère} & \textbf{Question à se poser} & \textbf{Correction à apporter} \\
\hline
Clarté & Mes phrases sont-elles simples et compréhensibles ? & Raccourcir les phrases trop longues, supprimer les mots vagues. \\
\hline
Correction & Orthographe, grammaire, conjugaisons sont-elles exactes ? & Vérifier les accords, les terminaisons, la ponctuation. \\
\hline
Cohérence & Les idées s'enchaînent-elles logiquement ? & Ajouter des connecteurs (en effet, ainsi, c'est pourquoi). \\
\hline
Respect des étapes & Les 4 mouvements sont-ils présents et ordonnés ? & Ajouter le mouvement manquant, remettre dans l'ordre. \\
\hline
\end{tabularx}
\end{center}

\begin{exoresolu}[Repérer et corriger les défauts d'une introduction]
Énoncé. Voici l'introduction d'un élève sur le sujet « Le conflit tradition/modernité dans \emph{Le Lion et la perle} » : « Dans cette dissertation, je vais parler du conflit entre la tradition et la modernité. Le Lion et la perle est une pièce de Soyinka. Lakunle veut épouser Sidi sans payer la dot. Baroka aussi veut l'épouser. Je vais d'abord parler de Lakunle, ensuite de Baroka. » Relevez quatre défauts et proposez une correction.
\tcblower
Solution détaillée. Quatre défauts sont identifiables :
\begin{enumerate}
\item \textbf{Absence d'amorce} : l'élève entre directement dans le sujet. Correction : ajouter un fait général, par exemple « Toute société africaine moderne vit la tension entre coutumes et progrès. »
\item \textbf{Emploi du « je »} : « je vais parler de » est interdit. Correction : « nous verrons d'abord que Lakunle incarne une modernité superficielle, puis que Baroka défend une tradition rusée. »
\item \textbf{Absence de problématique} : aucune question n'est posée. Correction : ajouter « Comment ce conflit entre deux prétendants reflète-t-il les tensions de l'Afrique contemporaine ? »
\item \textbf{Simple résumé de l'histoire} : raconter l'intrigue ne présente pas le sujet. Correction : reformuler le sujet en nommant les notions (tradition incarnée par Baroka, modernité incarnée par Lakunle).
\end{enumerate}
L'introduction corrigée comportera donc : amorce, présentation de l'œuvre et du sujet, problématique, annonce du plan au « nous ».
\end{exoresolu}

\begin{exercice}[À vous de rédiger]
Sujet : « La condition de la femme dans \emph{Le Lion et la perle}. »
\begin{enumerate}
\item Soulignez les mots-clés et définissez-les.
\item Dégagez le thème et les limites du sujet.
\item Formulez une problématique.
\item Rédigez une introduction complète de 8 à 12 lignes comportant les quatre mouvements, puis relisez-la avec la grille d'amélioration.
\end{enumerate}
\end{exercice}

\begin{corrige}[Exercice : La condition de la femme]
\textbf{1. Mots-clés} : \emph{condition} = situation sociale, place dans la société ; \emph{femme} = ici Sidi, Sadiku et les autres femmes du village.\par
\textbf{2. Thème et limites} : la place des femmes dans la société traditionnelle d'Ilujinle ; on se limite à la pièce de Soyinka, sans exposé général sur les femmes dans le monde.\par
\textbf{3. Problématique} : quelle image de la femme la pièce donne-t-elle à voir : objet des hommes ou force agissante ?\par
\textbf{4. Introduction possible} : « Dans les sociétés traditionnelles, la place de la femme est souvent débattue entre soumission apparente et influence réelle. C'est cette question que soulève Wole Soyinka dans \emph{Le Lion et la perle}, pièce publiée en 1963. La jeune Sidi, courtisée par Lakunle et Baroka, semble d'abord un prix à conquérir ; pourtant, son orgueil et ses choix font d'elle un personnage central. Quelle image de la femme la pièce donne-t-elle à voir ? Nous verrons d'abord que la femme apparaît comme un objet de désir et d'échange, puis nous montrerons qu'elle exerce en réalité un pouvoir décisif sur le dénouement. »
\end{corrige}

\begin{aretenir}[L'essentiel de la section]
\begin{itemize}
\item Analyser un sujet se fait en \textbf{4 étapes} : mots-clés, thème et limites, problématique, plan.
\item La \textbf{problématique} est la question centrale qui oriente toute la dissertation.
\item L'introduction comporte \textbf{4 mouvements} : amorce, présentation du sujet, problématique, annonce du plan ; elle fait 8 à 12 lignes.
\item On annonce le plan avec « nous verrons que... », jamais avec « je vais parler de... ».
\item Toute introduction se relit avec la grille : clarté, correction, cohérence, respect des étapes.
\end{itemize}
\end{aretenir}

Maintenant que nous savons analyser un sujet et construire une introduction solide, la section suivante nous apprendra, à partir de la deuxième lecture méthodique de la pièce, à distinguer dénotation et connotation et à rédiger un paragraphe argumentatif avec citations et transitions.

\coursec{Lecture méthodique (2), dénotation/connotation et le paragraphe argumentatif}

Dans la section précédente, nous avons appris à analyser un sujet de dissertation et à rédiger une introduction complète (amorce, présentation de l'œuvre, problématique, annonce du plan). Mais une introduction ne suffit pas : c'est dans le \cle{développement} que l'élève démontre réellement sa capacité à penser et à écrire. Or le développement est fait de \cle{paragraphes argumentatifs}. Pour rédiger de bons paragraphes, il faut deux outils : une connaissance précise de l'œuvre étudiée (nous poursuivons donc la lecture méthodique du \emph{Le Lion et la perle} de Wole Soyinka) et une maîtrise du sens des mots (dénotation et connotation), car argumenter, c'est d'abord comprendre ce que les mots disent au-delà d'eux-mêmes.

\courssub{Lecture méthodique n°2 : la deuxième partie, « Midi »}

\textbf{Rappel de la première partie (« Matin »).} Dans le village d'Ilujinlè, Sidi, la plus belle jeune fille du village, a vu ses photographies publiées dans un magazine. Elle se croit désormais plus célèbre que le Baalè (chef) lui-même, Baroka, et refuse sa demande en mariage, transmise par Sadiku, sa première épouse.

\textbf{Déroulement de la deuxième partie.} L'action se passe à midi, au palais de Baroka.

\begin{enumerate}
\item \textbf{La confidence de Baroka.} Baroka, âgé de soixante-deux ans, est un chef rusé et attaché aux traditions. Il révèle à Sadiku un prétendu secret : il serait devenu \emph{impuissant} depuis peu. Sadiku, ravie de tenir une arme contre son mari tyrannique, court annoncer la nouvelle à Sidi et l'invite à venir narguer le vieux chef lors d'un souper au palais.
\item \textbf{Le stratagème du « Lion ».} En réalité, l'impuissance de Baroka est un \cle{mensonge calculé} : c'est un piège pour attirer Sidi au palais. Sidi, trop sûre d'elle-même et flattée par sa beauté photographiée, accepte de venir se moquer du « lion » affaibli.
\item \textbf{La séduction de Sidi.} Au palais, Baroka flatte Sidi, lui montre un projet de timbre-poste à son effigie, joue de sa vanité et de son désir de modernité. Sidi, vaincue par la ruse du vieux lion, cède à ses avances. À la fin de la partie, Sadiku et Lakunle (l'instituteur amoureux de Sidi) dansent en célébrant la défaite supposée de Baroka, sans savoir que le piège a fonctionné.
\end{enumerate}

\begin{aretenir}[L'essentiel de la deuxième partie]
La deuxième partie met en scène le \cle{conflit entre la ruse traditionnelle} (Baroka, le « lion ») et \cle{l'orgueil moderne} (Sidi, la « perle »). Le stratagème de la fausse impuissance montre que la force du lion n'est pas physique mais intellectuelle : Baroka gagne par l'intelligence et la connaissance de la psychologie humaine.
\end{aretenir}

\begin{savaistu}[Wole Soyinka, un lion des lettres africaines]
Wole Soyinka, dramaturge nigérian né en 1934, est le premier Africain à recevoir le prix Nobel de littérature (1986). \emph{Le Lion et la perle} (1963) est une comédie qui interroge, avec humour, le choc entre tradition et modernité dans l'Afrique de l'indépendance — une question qui concerne aussi le Cameroun de la même époque.
\end{savaistu}

\courssub{La dénotation : le sens premier et objectif}

\textbf{Intuition.} Si je demande à dix élèves : « Qu'est-ce qu'un lion ? », tous répondront à peu près la même chose : un grand félin sauvage à crinière, carnivore, vivant en Afrique. Ce sens partagé, stable, que tout le monde accepte, c'est la dénotation.

\begin{definition}[La dénotation]
La \cle{dénotation} est le sens \textbf{stable, objectif et premier} d'un mot, celui que donne le dictionnaire. Elle est identique pour tous les locuteurs d'une langue, indépendamment des émotions et des contextes.
\end{definition}

\begin{exemplebox}[Dénotations]
\begin{itemize}
\item \emph{Lion} : mammifère carnivore de la famille des félidés, le mâle portant une crinière.
\item \emph{Perle} : concrétion dure et brillante produite par certaines huîtres, utilisée comme bijou.
\item \emph{Village} : agglomération rurale de faible importance, plus petite qu'une ville.
\end{itemize}
\end{exemplebox}

\courssub{La connotation : le sens second, culturel et affectif}

\textbf{Intuition.} Mais quand un élève camerounais entend « mon village », il ne pense pas à une définition de dictionnaire : il pense aux vacances chez grand-mère, au ndolè de Noël, aux veillées, aux ancêtres. Ces idées et émotions \emph{associées} au mot, c'est la connotation.

\begin{definition}[La connotation]
La \cle{connotation} est l'ensemble des idées, des valeurs et des émotions qu'un mot évoque \textbf{en plus} de son sens premier, selon le contexte, la culture et la sensibilité de chacun. Elle est subjective et variable.
\end{definition}

\begin{exemplebox}[Connotations au Cameroun]
\begin{itemize}
\item \emph{Village} : origine, racines, tradition, famille élargie, mais parfois aussi « retard », « simplicité » selon le point de vue.
\item \emph{Lion} : force, courage, royauté, fierté nationale (les Lions indomptables !), mais aussi danger, férocité.
\item \emph{Perle} : beauté rare, pureté, préciosité, objet de désir.
\end{itemize}
\end{exemplebox}

\begin{popfigure}
\begin{center}
\begin{tikzpicture}
% Cercle central : dénotation
\draw[fill=popblueL, draw=popblue, thick] (0,0) circle (1.7);
\node[align=center, text width=3cm] at (0,0) {\textbf{Dénotation}\\ grand félin\\ carnivore\\ à crinière};
% Cercle extérieur : connotations
\draw[draw=poporange, thick, dashed] (0,0) circle (3.6);
\node[text=poporange, align=center] at (0,3.15) {\textbf{Connotations}};
% Rayons et étiquettes
\foreach \angle/\txt in {90/{royauté,\\ chef}, 162/{courage}, 234/{fierté\\ nationale}, 306/{danger,\\ férocité}, 18/{force}}{
  \draw[->, poporange, thick] (\angle:1.7) -- (\angle:3.5);
  \node[align=center, text width=2.2cm] at (\angle:4.35) {\small \txt};
}
\node[align=center] at (0,-4.6) {\emph{Le mot « lion » : un sens central stable, des sens associés variables.}};
\end{tikzpicture}
\end{center}
\legende{Schéma du sens d'un mot : la dénotation occupe le centre (sens du dictionnaire) ; les connotations rayonnent autour selon les contextes culturels.}
\end{popfigure}

\courssub{Exploitation : les connotations du titre \emph{Le Lion et la perle}}

C'est ici que la sémantique devient un outil d'analyse littéraire. Le titre de Soyinka ne désigne pas un animal et un bijou : il fonctionne par \cle{connotation}.

\begin{center}
\begin{tabularx}{\linewidth}{|l|X|X|}
\hline
\rowcolor{popblueL} \textbf{Mot du titre} & \textbf{Dénotation} & \textbf{Connotations exploitées dans la pièce} \\
\hline
Lion & Grand félin carnivore & Baroka : force, autorité du chef, ruse du prédateur, virilité, tradition enracinée \\
\hline
Perle & Concrétion précieuse de l'huître & Sidi : beauté rare, jeunesse, pureté, objet de convoitise que tous veulent posséder \\
\hline
\end{tabularx}
\end{center}

Ainsi, le titre annonce déjà le conflit de la deuxième partie : le lion (Baroka) finira par capturer la perle (Sidi), non par la force brutale, mais par la ruse — car le lion de Soyinka est un \emph{vieux} lion dont la puissance est devenue intelligence.

\begin{attention}[Piège fréquent]
Ne confonds pas connotation et sens propre. Dire « Baroka est un lion » n'est pas une erreur zoologique : c'est une \cle{métaphore} qui repose sur les connotations du mot « lion » (force, royauté). Dans une dissertation, il faut \emph{expliquer} ces connotations au lieu de les subir.
\end{attention}

\courssub{La structure du paragraphe de dissertation}

Un développement de dissertation comporte généralement deux ou trois grandes parties, elles-mêmes divisées en paragraphes. Chaque paragraphe obéit à une architecture rigoureuse.

\begin{definition}[Le paragraphe argumentatif]
Le \cle{paragraphe argumentatif} est une unité de pensée qui développe \textbf{une seule idée directrice}, l'explique, la prouve par une citation et une illustration, puis la conclut par une phrase de synthèse.
\end{definition}

\begin{exemplebox}[Un paragraphe efficace en chiffres]
Un bon paragraphe de dissertation compte de \textbf{8 à 15 lignes} et défend \textbf{une seule idée directrice}. En dessous de 8 lignes, l'idée n'est pas développée ; au-delà de 15 lignes, on mélange probablement plusieurs idées.
\end{exemplebox}

\begin{popfigure}
\begin{center}
\begin{tikzpicture}[node distance=0.55cm,
  etape/.style={draw=popblue, fill=popblueL, rounded corners, thick, text width=8.6cm, align=left, inner sep=6pt},
  fleche/.style={->, thick, poporange}]
\node[etape] (e1) {\textbf{1. Idée principale} : la phrase d'ouverture énonce l'argument du paragraphe.};
\node[etape, below=of e1] (e2) {\textbf{2. Explication} : on développe et on précise l'idée avec ses propres mots.};
\node[etape, below=of e2] (e3) {\textbf{3. Citation} : on appuie l'idée sur le texte de l'œuvre, entre guillemets.};
\node[etape, below=of e3] (e4) {\textbf{4. Illustration / commentaire} : on exploite la citation, on montre ce qu'elle prouve.};
\node[etape, below=of e4] (e5) {\textbf{5. Phrase de synthèse} : on conclut le paragraphe et on prépare la suite.};
\draw[fleche] (e1) -- (e2);
\draw[fleche] (e2) -- (e3);
\draw[fleche] (e3) -- (e4);
\draw[fleche] (e4) -- (e5);
\end{tikzpicture}
\end{center}
\legende{Les cinq étapes de la structure d'un paragraphe argumentatif : idée, explication, citation, illustration, synthèse.}
\end{popfigure}

\courssub{Les techniques d'insertion des citations}

La citation est la preuve que l'on connaît l'œuvre. Mais elle doit être \emph{insérée}, c'est-à-dire amenée et exploitée.

\begin{methode}[Insérer une citation en trois temps]
\begin{enumerate}
\item \textbf{Introduire} la citation avec une formule : \emph{comme le dit Baroka...}, \emph{selon Sadiku...}, \emph{Soyinka écrit...}, ou par une citation intégrée à la phrase : Baroka se présente comme « ... ».
\item \textbf{Citer} le passage exact de l'œuvre, \textbf{entre guillemets}, sans le déformer, en précisant si possible la référence (partie, scène).
\item \textbf{Commenter} la citation : expliquer ce qu'elle signifie et en quoi elle fait avancer l'argumentation.
\end{enumerate}
\end{methode}

\begin{exemplebox}[Les trois formes d'insertion]
\begin{itemize}
\item \textbf{Citation introduite} : Comme le déclare Baroka à Sadiku : « ... », le chef feint la faiblesse pour mieux piéger Sidi.
\item \textbf{Citation intégrée} : Baroka, en avouant à Sadiku sa prétendue « impuissance », tend en réalité un piège à la vanité de Sidi.
\item \textbf{Citation avec référence} : « ... » (deuxième partie, « Midi ») : cette confidence mensongère révèle la ruse du vieux lion.
\end{itemize}
\end{exemplebox}

\begin{attention}[Erreur fréquente : la citation parachutée ou paraphrasée]
Deux fautes tuent un paragraphe :
\begin{itemize}
\item la \textbf{citation parachutée} : poser une citation sans l'introduire ni la commenter. Elle ne prouve rien si on ne l'exploite pas ;
\item la \textbf{paraphrase} : répéter la citation avec d'autres mots au lieu de l'\emph{exploiter}. Commenter, ce n'est pas redire : c'est montrer \emph{ce que la citation prouve} et \emph{comment elle fait avancer l'argument}.
\end{itemize}
Règle d'or : après chaque citation, écris au moins deux phrases de commentaire.
\end{attention}

\courssub{La transition : le pont entre deux paragraphes}

\begin{definition}[La transition]
La \cle{transition} est une phrase (ou deux) placée entre deux paragraphes ou deux parties, qui \textbf{rappelle ce qui précède} (bilan) et \textbf{annonce ce qui suit} (ouverture). Elle assure la cohérence et la fluidité du devoir.
\end{definition}

\begin{methode}[Rédiger une transition en deux formules]
\begin{enumerate}
\item \textbf{Formule de bilan} : résumer l'idée qui vient d'être démontrée. Exemple : « Ainsi, Baroka apparaît d'abord comme le gardien rusé de la tradition. »
\item \textbf{Formule d'annonce} : introduire l'idée du paragraphe suivant. Exemple : « Mais sa victoire sur Sidi ne tient pas seulement à sa ruse : elle repose aussi sur la vanité de la jeune fille. »
\end{enumerate}
La transition complète : « Ainsi, Baroka apparaît d'abord comme le gardien rusé de la tradition ; mais sa victoire sur Sidi ne tient pas seulement à sa ruse : elle repose aussi sur la vanité de la jeune fille. »
\end{methode}

\begin{exoresolu}[Rédiger un paragraphe avec citation et transition]
\textbf{Énoncé.} Sujet : « Dans \emph{Le Lion et la perle}, la ruse est-elle plus forte que la beauté ? » Rédige un paragraphe argumentatif (avec citation insérée) montrant que Baroka triomphe par l'intelligence, puis la transition qui annonce le paragraphe suivant sur la vanité de Sidi.
\tcblower
\textbf{Solution détaillée.}

\emph{Paragraphe :} Baroka, malgré son âge, triomphe de Sidi grâce à son intelligence et non à sa force physique. Conscient qu'il ne peut séduire la jeune fille par la jeunesse ou la beauté, le vieux chef élabore un stratagème fondé sur la connaissance de l'âme humaine : il feint l'impuissance. Comme il le confie à Sadiku dans la confidence mensongère de la deuxième partie, il se présente comme un lion « affaibli », incapable de nuire. Cette fausse faiblesse est en réalité une arme redoutable : en rassurant Sidi, Baroka l'attire dans son palais, terrain où il est maître du jeu. La ruse du lion consiste donc à transformer son apparente défaite en piège ; c'est bien l'intelligence, et non la force, qui fait sa puissance. Ainsi, le stratagème de Baroka prouve que, chez Soyinka, la véritable force du chef est d'ordre intellectuel.

\emph{Transition :} Cependant, la ruse de Baroka n'aurait pas suffi si Sidi n'avait pas été aveuglée par sa propre vanité : c'est donc l'orgueil de la « perle » qu'il convient maintenant d'examiner.

\emph{Commentaire de la copie :} le paragraphe compte environ dix lignes, une seule idée directrice (la victoire par l'intelligence), une citation introduite et commentée, et une phrase de synthèse. La transition fait le bilan (« la ruse de Baroka ») puis annonce la suite (« la vanité de Sidi »).
\end{exoresolu}

\begin{exercice}[À toi de jouer]
\begin{enumerate}
\item Relève dans la deuxième partie deux mots à forte connotation et explique ce qu'ils évoquent dans la culture du village d'Ilujinlè.
\item Rédige un paragraphe argumentatif (8 à 15 lignes) sur l'idée suivante : « Sidi est victime de sa vanité. » Insère une citation de la deuxième partie selon l'une des trois techniques étudiées, puis commente-la.
\item Rédige la transition qui permettrait d'enchaîner ton paragraphe avec un paragraphe suivant consacré à Lakunle, l'instituteur moderniste.
\end{enumerate}
\end{exercice}

\begin{corrige}[Exercice : éléments de correction]
\begin{enumerate}
\item Exemples attendus : « lion » (connote la force et l'autorité du chef Baroka) ; « perle » (connote la beauté rare et précieuse de Sidi) ; « palais » (connote le pouvoir et le territoire du chef) ; « timbre-poste » (connote la modernité et la gloire).
\item Le paragraphe doit contenir : une phrase d'idée, une explication, une citation introduite (ex. « comme le montre l'attitude de Sidi lorsque... »), deux phrases de commentaire et une phrase de synthèse. Vérifier qu'il n'y a qu'une seule idée directrice et que la citation n'est ni parachutée ni paraphrasée.
\item Modèle de transition : « Si la vanité de Sidi l'a conduite au palais de Baroka, elle s'explique aussi par l'influence de la modernité incarnée par Lakunle, l'instituteur dont les idées nouvelles ont troublé le village. »
\end{enumerate}
\end{corrige}

\begin{aretenir}[L'essentiel de la section]
\begin{itemize}
\item La deuxième partie met en scène le stratagème de Baroka (fausse impuissance) qui piège Sidi : la ruse triomphe de la beauté.
\item La \cle{dénotation} est le sens objectif du dictionnaire ; la \cle{connotation} est l'ensemble des idées et émotions associées à un mot selon la culture et le contexte.
\item Le titre \emph{Le Lion et la perle} fonctionne par connotations : Baroka = force et ruse du chef ; Sidi = beauté rare et convoitée.
\item Un paragraphe argumentatif suit cinq étapes : idée, explication, citation, commentaire, synthèse (8 à 15 lignes, une seule idée).
\item Toute citation s'introduit, se met entre guillemets et se commente ; toute transition fait un bilan puis une annonce.
\end{itemize}
\end{aretenir}

\coursec{Lecture méthodique (3), énonciation et conclusion}

Dans la section précédente, nous avons lu l'acte II du \emph{Lion et la perle} et appris à construire un paragraphe argumentatif en insérant des citations et en rédigeant des transitions. Nous avons aussi distingué la \cle{dénotation} de la \cle{connotation}, c'est-à-dire le sens premier d'un mot et les idées qu'il évoque. Nous voici maintenant au terme de la pièce : l'acte III, intitulé « Nuit », nous révèle le choix final de Sidi et le dénouement de l'intrigue. Cette section nous permettra également de découvrir la notion d'\cle{énonciation} et d'apprendre à rédiger la \cle{conclusion} de la dissertation, dernière étape de notre travail d'écriture.

\courssub{Lecture méthodique n° 3 : l'acte III (« Nuit »)}

\textbf{Situation de l'acte.} La nuit est tombée sur le village d'Ilujinle. Sadiku, la première épouse de Baroka, est allée annoncer à Sidi que le Lion a perdu sa virilité : c'est un mensonge calculé, un piège tendu par Baroka lui-même pour attirer la jeune fille dans son palais. Sidi, d'abord moqueuse, se rend chez le vieux chef, prétendument pour partager son repas. Là, Baroka la séduit par la ruse : il lui montre la machine à timbres, symbole du progrès, et lui fait miroiter l'idée que son image pourrait circuler sur tout le territoire.

\textbf{Le choix final de Sidi.} Pendant ce temps, Lakunle, l'instituteur, attend Sidi pour l'épouser « à la moderne », sans dot. Mais Sidi revient du palais en larmes : Baroka l'a possédée. Contre toute attente, elle annonce alors qu'elle épousera... Baroka, le Lion. Lakunle, effondré, se console en regardant une jeune fille qui passe. La pièce se referme sur la danse et le chant : la \cle{tradition} a triomphé, mais une tradition qui a su utiliser les armes de la \cle{modernité} (la machine, les timbres, l'image).

\textbf{Portée de la pièce.} Wole Soyinka ne se contente pas de raconter une histoire d'amour : il met en scène le conflit entre la tradition (Baroka, la dot, la polygamie) et la modernité (Lakunle, l'école, le mariage « civilisé »). Or le dénouement est ambigu : c'est le vieux Lion qui gagne, mais en trichant, et la modernité de Lakunle apparaît comme superficielle et ridicule. Soyinka invite donc le lecteur à réfléchir : faut-il rejeter la tradition au nom d'un progrès mal compris ? C'est précisément le genre de question qu'une dissertation doit savoir exploiter.

\begin{aretenir}[Repères sur l'acte III]
\begin{itemize}
\item Le stratagème de Baroka (fausse impuissance) réussit : Sidi devient sa femme.
\item Lakunle, défenseur de la modernité, est le grand perdant de la pièce.
\item Le dénouement est ambigu : victoire de la tradition, mais par la ruse.
\item Thème central : le choc entre tradition et modernité dans l'Afrique en mutation.
\end{itemize}
\end{aretenir}

\courssub{Communication : l'énonciation}

\textbf{Intuition d'abord.} Quand Sidi dit : « \emph{Je} n'épouserai jamais un homme qui refuse de payer ma dot », elle ne transmet pas seulement une information : elle se met en scène elle-même (\emph{je}), elle s'adresse à quelqu'un (Sadiku, Lakunle), dans une situation précise (une dispute, un moment donné). Le même contenu, dit autrement — « On m'a dit que... », « Il paraît que... » — produirait un effet tout différent. C'est cela, l'énonciation : la manière dont le locuteur s'inscrit dans son propre discours.

\begin{definition}[L'énonciation]
L'\cle{énonciation} est la mise en scène du \cle{locuteur} (celui qui parle ou écrit) et du \cle{destinataire} (celui à qui l'on s'adresse) dans le discours, au sein d'une \cle{situation de communication} déterminée (lieu, moment, circonstances). Elle se repère grâce aux \cle{marques énonciatives} présentes dans l'énoncé.
\end{definition}

\textbf{Les marques énonciatives.} Trois grandes familles de marques permettent d'identifier qui parle, à qui, quand et avec quel degré de certitude :

\begin{itemize}
\item Les \cle{pronoms personnels} : \emph{je}, \emph{nous} désignent le locuteur ; \emph{tu}, \emph{vous} désignent le destinataire ; \emph{il}, \emph{elle}, \emph{ils} renvoient à ce dont on parle (la « non-personne »).
\item Les \cle{temps verbaux} : le présent d'énonciation situe le propos au moment où l'on parle ; le passé composé lie l'événement au présent du locuteur ; l'imparfait et le passé simple créent un récit détaché du moment de parole.
\item Les \cle{modalisateurs} : adverbes et tournures qui expriment l'attitude du locuteur.
\end{itemize}

\begin{definition}[Les modalisateurs]
Les \cle{modalisateurs} sont des mots ou des expressions qui indiquent le degré de certitude, d'engagement ou d'attitude du locuteur face à ce qu'il dit. On distingue les modalisateurs d'\cle{affirmation} (\emph{certes, évidemment, il est clair que, assurément}) et les modalisateurs d'\cle{atténuation} (\emph{sans doute, peut-être, il semble que, probablement}).
\end{definition}

\begin{popfigure}
\begin{center}
\begin{tikzpicture}[scale=1.0, every node/.style={font=\small}]
% Triangle de la communication
\node[draw, fill=popblueL, rounded corners, minimum width=3.2cm, minimum height=1cm] (loc) at (0,3.2) {\textbf{Locuteur} (je / nous)};
\node[draw, fill=poporangeL, rounded corners, minimum width=3.2cm, minimum height=1cm] (dest) at (6,3.2) {\textbf{Destinataire} (tu / vous)};
\node[draw, fill=popgreenL, rounded corners, minimum width=4.4cm, minimum height=1cm] (msg) at (3,0) {\textbf{Message} (l'énoncé)};
\draw[-{Stealth}, thick, popblue] (loc) -- node[above left, align=center]{produit} (msg);
\draw[-{Stealth}, thick, poporange] (msg) -- node[above right, align=center]{reçu par} (dest);
\draw[{Stealth}-{Stealth}, thick, poppurple] (loc) -- node[above]{s'adresse à} (dest);
\node[draw, dashed, fill=popcream, rounded corners, align=center, minimum width=5cm] (sit) at (3,-1.6) {\textbf{Situation d'énonciation}\\ lieu, moment, circonstances};
\draw[dashed, popmuted] (sit) -- (msg);
\end{tikzpicture}
\end{center}
\legende{Le triangle de la communication : tout message est produit par un locuteur, adressé à un destinataire, dans une situation d'énonciation qui en colore le sens.}
\end{popfigure}

\begin{exemplebox}[Repérage des marques énonciatives]
Dans la phrase : « \emph{Je crois que Sidi a sans doute fait le bon choix} », on relève :
\begin{itemize}
\item \emph{je} : pronom de première personne, le locuteur s'engage ;
\item \emph{crois} : présent d'énonciation + verbe d'opinion (faible certitude) ;
\item \emph{sans doute} : modalisateur d'atténuation.
\end{itemize}
Le locuteur donne une opinion prudente, non un fait.
\end{exemplebox}

\courssub{L'énonciation dans la dissertation}

La dissertation est un discours argumenté adressé à un correcteur. L'élève doit donc contrôler ses marques énonciatives avec soin.

\begin{itemize}
\item Le \cle{nous de modestie} : « Nous avons vu que Baroka triomphe par la ruse. » Ce \emph{nous} associe le lecteur au raisonnement et évite le \emph{je} trop personnel.
\item Le \cle{présent de vérité générale} : « La tradition résiste à la modernité. » Il donne au propos une valeur universelle, intemporelle.
\item Les \cle{modalisateurs} : ils dosent l'affirmation. « Il est évident que Sidi exerce un choix » (affirmation forte) ; « On peut penser que Lakunle représente une modernité superficielle » (atténuation prudente).
\end{itemize}

\begin{savaistu}[« Je » ou « nous » ?]
Le \emph{je} est toléré dans la conclusion personnelle, lorsqu'on vous demande explicitement votre point de vue (« Je pense que... »). Mais dans le corps de la dissertation, le \emph{nous} reste la norme : il efface le moi de l'élève derrière un raisonnement partagé avec le lecteur. Au Probatoire et au Baccalauréat technique, l'abus de \emph{je} dans le développement est sanctionné comme une faute de posture énonciative.
\end{savaistu}

\courssub{La conclusion de la dissertation}

\textbf{Intuition.} Une dissertation sans conclusion, c'est un match de football arrêté avant le coup de sifflet final : le correcteur ne sait pas qui gagne. La conclusion est le moment où vous ramassez vos arguments, répondez franchement à la question posée et ouvrez une dernière perspective.

\begin{methode}[Rédiger la conclusion en trois mouvements]
La conclusion compte de 6 à 10 lignes et suit trois mouvements obligatoires :
\begin{enumerate}
\item \textbf{Le \cle{bilan}} : résumez en deux ou trois phrases les idées développées dans le corps du devoir (reprendre les étapes du plan, sans citer les parties).
\item \textbf{La \cle{réponse} à la problématique} : répondez clairement et nettement à la question posée par le sujet (oui, non, ou réponse nuancée, mais assumée).
\item \textbf{L'\cle{ouverture}} : élargissez le propos par une question connexe ou un prolongement vers un autre contexte, sans développer une idée nouvelle.
\end{enumerate}
\end{methode}

\begin{popfigure}
\begin{center}
\begin{tikzpicture}[every node/.style={font=\small}]
% Entonnoir inversé : du bilan vers l'ouverture
\draw[fill=popblueL, draw=popblue, thick] (0,4) -- (6,4) -- (5.2,2.7) -- (0.8,2.7) -- cycle;
\node at (3,3.35) {\textbf{1. Bilan} : rappel des idées développées};
\draw[fill=poporangeL, draw=poporange, thick] (0.8,2.5) -- (5.2,2.5) -- (4.6,1.3) -- (1.4,1.3) -- cycle;
\node at (3,1.9) {\textbf{2. Réponse} à la problématique};
\draw[fill=popgreenL, draw=popgreen, thick] (1.4,1.1) -- (4.6,1.1) -- (5.4,-0.3) -- (0.6,-0.3) -- cycle;
\node at (3,0.4) {\textbf{3. Ouverture} : élargissement};
\draw[-{Stealth}, very thick, popdark] (3,4.4) -- (3,-0.7);
\end{tikzpicture}
\end{center}
\legende{La structure de la conclusion : on resserre d'abord le propos (bilan, réponse précise), puis on l'élargit vers une perspective nouvelle (ouverture).}
\end{popfigure}

\textbf{Les techniques d'ouverture.} Deux procédés sont particulièrement sûrs au niveau de la Première technique :
\begin{itemize}
\item La \cle{question connexe} : poser une question voisine de celle du sujet. Exemple : « Ne faut-il pas alors se demander si la modernité de Lakunle était véritablement un progrès ? »
\item L'\cle{élargissement à un autre contexte} : transposer la réflexion à la société camerounaise contemporaine. Exemple : « Ce débat entre tradition et modernité demeure d'actualité dans nos villes et nos villages camerounais, où le mariage coutumier côtoie le mariage civil. »
\end{itemize}

\begin{attention}[Erreurs fréquentes dans la conclusion]
\begin{itemize}
\item \textbf{Introduire une idée nouvelle} : la conclusion ne doit contenir aucun argument inédit ; elle rassemble, elle n'ajoute pas.
\item \textbf{Terminer sans répondre à la problématique} : une conclusion qui ne dit pas clairement ce que l'on pense du sujet est une conclusion manquée.
\item \textbf{Recopier l'introduction} : le bilan n'est pas une répétition de l'annonce du plan.
\item \textbf{Faire trop court} : moins de 6 lignes, la conclusion paraît bâclée.
\end{itemize}
\end{attention}

\begin{exemplebox}[Conclusion rédigée]
Sujet : « Sidi, victime ou vainqueur dans \emph{Le Lion et la perle} ? »\par
\emph{Au terme de notre réflexion, nous avons montré que Sidi subit d'abord la ruse de Baroka, qui la piège en feignant d'avoir perdu sa virilité ; mais nous avons aussi vu qu'elle choisit librement d'épouser le Lion, refusant le mariage sans dot que lui propose Lakunle. Sidi apparaît donc à la fois victime d'un stratagème et vainqueur de son propre destin : trompée sur les moyens, elle décide en dernier ressort de sa vie et accède au statut envié d'épouse du chef. Ce choix invite à élargir le débat : dans la société camerounaise contemporaine, les jeunes filles ne sont-elles pas encore partagées entre les attentes de la tradition et les promesses de la modernité ?}
\end{exemplebox}

\courssub{Production et amélioration : rédiger et évaluer une conclusion}

\begin{exoresolu}[Rédiger la conclusion d'une dissertation]
Énoncé : à partir du sujet « La tradition triomphe-t-elle de la modernité dans \emph{Le Lion et la perle} ? », rédigez une conclusion de 6 à 10 lignes comportant un bilan, une réponse à la problématique et une ouverture.
\tcblower
Solution détaillée :\par
\textbf{Étape 1 — le bilan.} « Nous avons vu que Baroka, incarnation de la tradition, l'emporte sur Lakunle, le défenseur de la modernité, en séduisant Sidi ; mais nous avons également remarqué que le vieux chef utilise les outils du progrès, la machine à timbres et l'image photographique, pour parvenir à ses fins. »\par
\textbf{Étape 2 — la réponse.} « La tradition triomphe donc en apparence, mais c'est une tradition rusée, capable d'emprunter les armes de la modernité : Soyinka suggère ainsi que les deux mondes ne s'opposent pas autant qu'ils se mélangent. »\par
\textbf{Étape 3 — l'ouverture.} « Ne pourrait-on dès lors se demander si le véritable avenir de nos sociétés africaines ne réside pas dans un équilibre entre l'héritage des anciens et les exigences du monde moderne ? »\par
La conclusion complète compte huit lignes, respecte les trois mouvements et n'introduit aucune idée nouvelle.
\end{exoresolu}

\begin{exercice}[Relecture croisée entre pairs]
Rédigez individuellement la conclusion du sujet : « Lakunle est-il un personnage ridicule ? » puis échangez votre copie avec un camarade. Évaluez sa conclusion à l'aide de la grille suivante (1 point par critère) :
\begin{center}
\begin{tabularx}{\linewidth}{|l|X|c|}
\hline
\rowcolor{popblueL} \textbf{Critère} & \textbf{Question à se poser} & \textbf{Point} \\
\hline
Bilan & Les idées du développement sont-elles résumées ? & /1 \\
\hline
Réponse & La problématique reçoit-elle une réponse claire ? & /1 \\
\hline
Ouverture & Y a-t-il une question connexe ou un élargissement ? & /1 \\
\hline
Énonciation & Le \emph{nous} et le présent de vérité générale sont-ils employés ? & /1 \\
\hline
Longueur & La conclusion compte-t-elle de 6 à 10 lignes ? & /1 \\
\hline
\end{tabularx}
\end{center}
Corrigez ensemble les points faibles repérés avant de remettre la version améliorée au professeur.
\end{exercice}

\begin{corrige}[Exercice — relecture croisée]
Il n'existe pas de conclusion unique : toute conclusion obtenant 5/5 à la grille est recevable. Exemple attendu : « Nous avons montré que Lakunle multiplie les comportements ridicules : il refuse la dot, récite des formules apprises et méprise les coutumes du village. Cependant, son personnage annonce aussi les transformations de la société nigériane. Lakunle est donc ridicule dans ses manières, mais porteur d'un avenir réel. Ce personnage ne rappelle-t-il pas certains de nos jeunes intellectuels camerounais, partagés entre l'école et la tradition familiale ? » Barème : bilan 1 pt ; réponse nuancée et claire 1 pt ; ouverture 1 pt ; énonciation correcte 1 pt ; longueur respectée 1 pt.
\end{corrige}

\begin{aretenir}[L'essentiel de la section]
\begin{itemize}
\item L'acte III montre le triomphe ambigu de Baroka : Sidi choisit la tradition, mais une tradition qui triche avec la modernité.
\item L'énonciation est la mise en scène du locuteur et du destinataire dans le discours ; ses marques sont les pronoms, les temps et les modalisateurs.
\item Dans la dissertation, on privilégie le \emph{nous} de modestie, le présent de vérité générale et des modalisateurs dosés.
\item La conclusion suit trois mouvements : bilan, réponse à la problématique, ouverture — en 6 à 10 lignes, sans idée nouvelle.
\end{itemize}
\end{aretenir}

\coursec{Champ lexical, champ sémantique et texte descriptif}

Après avoir maîtrisé la rédaction de l'introduction, du développement et de la conclusion d'une dissertation, ainsi que les notions d'énonciation lors de la troisième lecture méthodique du \emph{Le Lion et la perle}, nous abordons maintenant deux outils lexicaux indispensables pour enrichir votre écriture et un genre textuel majeur : le \cle{texte descriptif}. Ces notions vous permettront d'éviter les répétitions, de varier votre vocabulaire et de donner de la densité à vos paragraphes argumentatifs comme à vos passages descriptifs.

\courssub{Le champ lexical : définition, repérage et utilisation}

\begin{definition}[Champ lexical]
On appelle \textbf{champ lexical} l'ensemble des mots (noms, verbes, adjectifs, adverbes) qui, dans un texte ou dans la langue, se rapportent à une même \textbf{notion}, à un même \textbf{thème} ou à une même \textbf{réalité conceptuelle}. Ces mots appartiennent à des classes grammaticales différentes mais convergent vers un même centre d'intérêt sémantique.
\end{definition}

\emph{Intuition :} Imaginez un filet de pêche. Le thème est le « poisson » (la notion centrale). Dans le filet, vous trouverez : \emph{truite, saumon, carpe, pêcher, canne, hameçon, eau, rivière, écaille, nageoire}. Tous ces mots forment le champ lexical de la pêche / du poisson. Dans une dissertation, si votre thèse porte sur « la tradition », votre « filet » contiendra : \emph{coutume, ancêtre, rite, danse, héritage, patrimoine, aïeul, transmission, sacré, tabou}.

\begin{exemplebox}[Champ lexical de la tradition dans \emph{Le Lion et la perle}]
Dans la pièce de Wole Soyinka, le village d'Ilujinle est imprégné de tradition. Le champ lexical associé est dense :
\begin{itemize}
    \item \textbf{Noms :} coutume, ancêtre, rite, danse, masque, oracle, fétiche, tabou, héritage, lignée, clan, chefferie, palais, sage, griot, tam-tam, incantation, sacrifice, offrande.
    \item \textbf{Verbes :} vénérer, honorer, transmettre, perpétuer, respecter, célébrer, invoquer, sacrifier, danser, chanter, raconter.
    \item \textbf{Adjectifs :} ancestral, sacré, coutumier, immémorial, immuable, vénérable, mystique, rituel, tribal, communautaire.
\end{itemize}
Ces mots tissent la toile de fond culturelle contre laquelle se détache le conflit modernité / tradition.
\end{exemplebox}

\begin{popfigure}
\begin{tikzpicture}[
    mindmap,
    concept color=popblue,
    grow cyclic,
    level 1/.style={level distance=3.5cm,sibling angle=360/4},
    level 2/.style={level distance=2.5cm,sibling angle=360/5},
    every node/.style={font=\small}
]
\node[concept, text width=2.2cm, align=center] {Coutume\\Tradition}
    child[concept color=poporange] {node[concept, text width=2cm, align=center] {Rites\\Sacrés}
        child {node[concept, text width=1.5cm, align=center] {Sacrifice}}
        child {node[concept, text width=1.5cm, align=center] {Incantation}}
        child {node[concept, text width=1.5cm, align=center] {Offrande}}
        child {node[concept, text width=1.5cm, align=center] {Tabou}}
        child {node[concept, text width=1.5cm, align=center] {Oracle}}
    }
    child[concept color=popgreen] {node[concept, text width=2cm, align=center] {Héritage\\Ancestral}
        child {node[concept, text width=1.5cm, align=center] {Lignée}}
        child {node[concept, text width=1.5cm, align=center] {Aïeul}}
        child {node[concept, text width=1.5cm, align=center] {Patrimoine}}
        child {node[concept, text width=1.5cm, align=center] {Transmission}}
        child {node[concept, text width=1.5cm, align=center] {Mémoire}}
    }
    child[concept color=poppurple] {node[concept, text width=2cm, align=center] {Célébration\\Communautaire}
        child {node[concept, text width=1.5cm, align=center] {Danse}}
        child {node[concept, text width=1.5cm, align=center] {Tam-tam}}
        child {node[concept, text width=1.5cm, align=center] {Chant}}
        child {node[concept, text width=1.5cm, align=center] {Masque}}
        child {node[concept, text width=1.5cm, align=center] {Griot}}
    }
    child[concept color=poppink] {node[concept, text width=2cm, align=center] {Autorité\\Coutumière}
        child {node[concept, text width=1.5cm, align=center] {Chefferie}}
        child {node[concept, text width=1.5cm, align=center] {Palais}}
        child {node[concept, text width=1.5cm, align=center] {Sage}}
        child {node[concept, text width=1.5cm, align=center] {Conseil}}
        child {node[concept, text width=1.5cm, align=center] {Décision}}
    };
\end{tikzpicture}
\legende{Carte mentale : champ lexical de la tradition autour du noyau « coutume » (extrait de \emph{Le Lion et la perle}).}
\end{popfigure}

\begin{methode}[Enrichir un paragraphe grâce au champ lexical]
Pour éviter les répétitions et donner de l'ampleur à votre développement :
\begin{enumerate}
    \item \textbf{Identifiez le mot-clé} de votre paragraphe (ex. : \emph{progrès}).
    \item \textbf{Construisez mentalement (ou sur brouillon) le champ lexical} associé : noms (\emph{évolution, modernisation, innovation, développement, essor}), verbes (\emph{progresser, évoluer, innover, transformer, moderniser}), adjectifs (\emph{progressiste, novateur, moderne, avant-gardiste, dynamique}), adverbes (\emph{progressivement, résolument, radicalement}).
    \item \textbf{Remplacez les occurrences répétées} du mot-clé par des synonymes, des hyperonymes, des périphrases ou des mots du champ lexical.
    \item \textbf{Vérifiez la cohérence} : chaque mot choisi doit coller au contexte précis de la phrase (registre, connotation, syntaxe).
\end{enumerate}
\end{methode}

\begin{exoresolu}[Varier le vocabulaire dans un paragraphe argumentatif]
\textbf{Énoncé :} Améliorez le paragraphe suivant en exploitant le champ lexical de la modernité : « La modernité arrive à Ilujinle. La modernité change les habitudes. Les jeunes veulent la modernité. Baroka refuse la modernité. Lakunle incarne la modernité. »
\tcblower
\textbf{Solution détaillée :}
\begin{enumerate}
    \item \textbf{Mot-clé :} \emph{modernité} (répété 5 fois).
    \item \textbf{Champ lexical de la modernité :} \emph{progrès, civilisation, innovation, renouveau, transformation, occidentalisation, école, chemin de fer, machine, technologie, urbanisation, émancipation, rupture, avant-garde}.
    \item \textbf{Réécriture :} « Le \textbf{progrès} fait irruption à Ilujinle. Cette \textbf{nouvelle donne} bouleverse les habitudes ancestrales. La jeunesse aspire à l'\textbf{émancipation} qu'offre le \textbf{renouveau}. Le Bale Baroka résiste farouchement à cette \textbf{occidentalisation} galopante. Quant à Lakunle, il se fait le héraut de cette \textbf{civilisation} importée, symbole d'une \textbf{modernité} qu'il veut imposer. »
    \item \textbf{Gain :} 5 occurrences $\rightarrow$ 1 seule (\emph{modernité} en finale pour clore), enrichissement sémantique (\emph{émancipation, occidentalisation, civilisation}), nuances argumentatives (Baroka résiste à l'\emph{occidentalisation}, Lakunle porte la \emph{civilisation}).
\end{enumerate}
\end{exoresolu}

\begin{attention}[Piège : confondre champ lexical et famille de mots]
La \textbf{famille de mots} (ou famille morphologique) regroupe des mots formés sur le même \textbf{radical} (\emph{modernité, moderniser, moderne, modernisme, modernement}). Le \textbf{champ lexical} regroupe des mots de \textbf{radicaux différents} unis par le \textbf{sens} (\emph{modernité, progrès, innovation, école, chemin de fer}). Au Bac, on vous demande d'exploiter le champ lexical pour varier le vocabulaire, pas de faire de la morphologie.
\end{attention}

\courssub{Le champ sémantique : des mots aux sens partagés}

\begin{definition}[Champ sémantique]
On appelle \textbf{champ sémantique} l'ensemble des \textbf{sens} (appelés \textbf{sèmes} ou traits sémantiques) communs à plusieurs mots ou expressions. C'est un \textbf{faisceau de significations} partagées. Un même mot peut appartenir à plusieurs champs sémantiques selon le contexte (polysémie). Le champ sémantique dépasse le seul lexique : il inclut les \textbf{périphrases}, les \textbf{métaphores}, les \textbf{métonymies} et les \textbf{expressions figées} qui véhiculent une même idée.
\end{definition}

\emph{Distinction cruciale :}
\begin{itemize}
    \item \textbf{Champ lexical} = ensemble de \textbf{MOTS} (formes graphiques/phoniques) $\rightarrow$ « Quels mots emploie-t-on pour parler de X ? »
    \item \textbf{Champ sémantique} = ensemble de \textbf{SENS} (contenus conceptuels) $\rightarrow$ « Quelles nuances de sens, quelles idées, quelles représentations associe-t-on à X ? »
\end{itemize}

\begin{exemplebox}[Champ sémantique de « la perle » chez Sidi]
Le mot \emph{perle} appartient au champ lexical du \textbf{bijou} (collier, bague, or, argent, joaillier).
Mais dans la bouche de Sidi (\emph{Le Lion et la perle}, p. 12 : « Je suis la perle, la gemme rare... »), il active un \textbf{champ sémantique} tout autre :
\begin{itemize}
    \item \textbf{Sèmes (traits de sens) :} [+ RARE], [+ PRÉCIEUX], [+ PUR], [+ BEAU], [+ INTACT], [+ CONVOITÉ], [+ UNIQUE], [+ FRAGILE].
    \item \textbf{Expressions / Périphrases associées :} « la prunelle de mes yeux », « un trésor inestimable », « une fleur qu'on ne cueille qu'une fois », « l'âme du village ».
\end{itemize}
Sidi ne parle pas de gemmologie ; elle construit son identité et sa valeur marchande symbolique. C'est la \textbf{connotation} (voir section précédente) qui opère ici : le champ sémantique \emph{perle} $\rightarrow$ \emph{beauté rare, virginité, honneur, fierté}.
\end{exemplebox}

\begin{popfigure}
\begin{center}
\begin{tabularx}{\linewidth}{|l|X|X|}
\hline
\rowcolor{popblueL}
\textbf{Critère} & \textbf{Champ lexical} & \textbf{Champ sémantique} \\
\hline
\textbf{Définition} & Ensemble de \textbf{mots} (formes) liés à une même notion. & Ensemble de \textbf{sens} (sèmes) partagés par des mots/expressions. \\
\hline
\textbf{Unité de base} & Le \textbf{mot} (lexème). & Le \textbf{sème} (trait sémantique). \\
\hline
\textbf{Question clé} & « Quels \textbf{mots} pour dire X ? » & « Quels \textbf{sens} pour dire X ? » \\
\hline
\textbf{Exemple : Modernité} & \emph{école, chemin de fer, machine, progrès, ville, usine, livre, écriture, horloge, loi} & Sèmes : [+NOUVEAU], [+RATIONNEL], [+EFFICACE], [+LIBÉRATEUR], [+RUPTURE].\\
& & Périphrases : « le vent du changement », « la marche du monde », « la lumière de la science » \\
\hline
\textbf{Polysémie} & Un mot = une entrée. \emph{Baroque} (style / bijou irrégulier) = 2 champs lexicaux. & Un mot = plusieurs champs sémantiques selon contexte. \emph{Perle} : [bijou] vs [rareté/beauté]. \\
\hline
\textbf{Dans la dissertation} & Varier le \textbf{vocabulaire} (synonymes, hyperonymes, dérivés). & Approfondir l'\textbf{argumentation} (nuances, connotations, images). \\
\hline
\textbf{Relation} & Support du champ sémantique. & Contenu véhiculé par le champ lexical. \\
\hline
\end{tabularx}
\end{center}
\legende{Tableau comparatif : champ lexical (mots) vs champ sémantique (sens) — illustration sur le thème de la modernité.}
\end{popfigure}

\begin{attention}[Erreur fréquente au Probatoire]
Ne confondez pas : \textbf{« Relever le champ lexical »} $\neq$ \textbf{« Analyser le champ sémantique »}.
\begin{itemize}
    \item Si le sujet demande : « Relevez les mots du champ lexical de la nature », vous listez : \emph{arbre, fleuve, oiseau, feuille, vent, ciel, herbe, racine...}
    \item Si le sujet demande : « Quel champ sémantique la nature active-t-elle ici ? », vous répondez : « La nature est investie des sèmes [+MATERNEL], [+PROTECTRICE], [+ÉTERNEL], [+SOURCE DE VIE], opposés aux sèmes [+HOSTILE], [+DÉSTRUCTEUR] du champ sémantique de la machine. »
\end{itemize}
Un même mot (\emph{feu}) appartient au champ lexical de la \emph{cuisine} (four, flamme, cuisson) ET au champ lexical de l'\emph{incendie} (flamme, brasier, cendre, victime). Mais ses champs sémantiques diffèrent : [+OUTIL, +MAÎTRISÉ, +UTILE] vs [+DANGER, +INCONTRÔLABLE, +DÉVASTATEUR].
\end{attention}

\courssub{Le texte descriptif : éléments constitutifs et outils}

\begin{definition}[Texte descriptif]
Le \textbf{texte descriptif} est un texte dont la \textbf{fonction principale} est de \textbf{représenter} un être, un lieu, un objet, une scène ou une atmosphère, en en donnant une \textbf{image mentale} au lecteur. Il fige le temps (souvent à l'imparfait ou au présent de narration) pour déployer l'espace. Il ne raconte pas une action (narration) mais \textbf{peint un tableau}.
\end{definition}

\begin{definition}[Éléments constitutifs du texte descriptif]
\begin{enumerate}
    \item \textbf{Le cadre spatial (ou spatio-temporel) :} lieu, moment, lumière, atmosphère générale (ex. : « Ilujinle, au cœur de l'après-midi torride... »).
    \item \textbf{Le thème descriptif (ou objet décrit) :} l'entité centrale (le village, la place du marché, la danse, le visage de Sidi, la moto de Lakunle).
    \item \textbf{Les notations descriptives (détails signifiants) :} les traits précis, concrets, sensoriels sélectionnés par l'observateur (couleurs, formes, sons, odeurs, textures, mouvements).
    \item \textbf{L'expansion (ou développement) :} l'organisation des notations selon une \textbf{logique spatiale} (du près au loin, du bas en haut, du général au particulier, par plans successifs) ou \textbf{thématique} (visuel $\rightarrow$ auditif $\rightarrow$ olfactif).
    \item \textbf{L'énumération (souvent) :} accumulation de détails pour créer l'abondance, la richesse, le foisonnement (asyndète ou polysyndète).
    \item \textbf{La comparaison / l'image :} métaphore, comparaison, personnification pour transfigurer le réel et marquer la subjectivité du descripteur.
\end{enumerate}
\end{definition}

\begin{methode}[Repérer les éléments descriptifs dans un extrait]
\begin{enumerate}
    \item \textbf{Balisez le cadre} : où ? quand ? quelle lumière ? (repérez les compléments circonstanciels de lieu/temps, les adjectifs d'atmosphère).
    \item \textbf{Identifiez le thème} : quoi ou qui est décrit ? (sujet des phrases descriptives, souvent introduit par \emph{c'est, voici, il y a} ou par une nomination nominale).
    \item \textbf{Surlignez les notations} : adjectifs qualificatifs, noms concrets, verbes d'état (\emph{être, paraître, sembler, avoir}), verbes de perception (\emph{voir, entendre, sentir, toucher, goûter}).
    \item \textbf{Analysez l'organisation (expansion) :} quel ordre ? (plan large $\rightarrow$ détail ; sensoriel visuel $\rightarrow$ auditif $\rightarrow$ olfactif ; statique $\rightarrow$ dynamique).
    \item \textbf{Repérez les figures de style} : comparaisons (\emph{comme, tel que}), métaphores, personnifications, métonymies. Elles révèlent le \textbf{point de vue} (subjectif / émotionnel) du descripteur.
\end{enumerate}
\end{methode}

\begin{exemplebox}[Analyse descriptive : le village d'Ilujinle (incipit de la pièce)]
\emph{« Le théâtre représente la place du village d'Ilujinle. Au fond, l'odàn (arbre sacré) étend ses branches immenses. À droite, le marché. À gauche, la route qui mène à la brousse. Devant, le palais du Bale. Le soleil darde ses rayons sur la place déserte. »} (Didascalies initiales)

\begin{center}
\begin{tabularx}{\linewidth}{|l|X|}
\hline
\textbf{Élément} & \textbf{Repérage dans l'extrait} \\
\hline
Cadre spatial & « place du village d'Ilujinle », « Au fond », « À droite », « À gauche », « Devant » (organisation quadripartite : arbre, marché, route, palais). \\
\hline
Thème descriptif & La place du village (espace central, lieu de vie et de pouvoir). \\
\hline
Notations & « branches immenses » (verticalité, protection), « soleil darde ses rayons » (lumière agressive, chaleur), « place déserte » (attente, silence, potentiel dramatique). \\
\hline
Expansion & Organisation \textbf{spatiale géométrique} (fond, droite, gauche, devant) = plan large, statique, objectif (didascalies). \\
\hline
Figures / Subjectivité & « darde » (verbe animalier, personification du soleil hostile), « odàn (arbre sacré) » (notation culturelle, charge symbolique). \\
\hline
\end{tabularx}
\end{center}
\end{exemplebox}

\begin{exemplebox}[Analyse descriptive : la danse des filles (extrait acte 1)]
\emph{« Les filles entrent en dansant. Elles portent des pots d'eau sur la tête. Leurs corps luisent de sueur. Le rythme du tam-tam s'accélère. Elles tournoient, se penchent, se redressent, les pots immobiles sur leurs têtes. Sidi est au centre, plus belle, plus fière. »}

\begin{itemize}
    \item \textbf{Cadre :} Place du village, plein jour, chaleur (« luisent de sueur »).
    \item \textbf{Thème :} La danse des filles / Sidi au centre.
    \item \textbf{Notations sensorielles :}
        \begin{itemize}
            \item \textbf{Visuel :} « pots d'eau sur la tête », « corps luisent », « tournoient, se penchent, se redressent », « pots immobiles » (contraste mouvement / statique).
            \item \textbf{Auditif :} « rythme du tam-tam s'accélère » (temporalité, crescendo).
            \item \textbf{Tactile / Kinesthésique :} « sueur », « luisent », effort physique implicite.
        \end{itemize}
    \item \textbf{Expansion :} \textbf{Dynamique} (entrée $\rightarrow$ accélération $\rightarrow$ apogée sur Sidi). Du groupe (filles) vers l'individu (Sidi) = \textbf{progression du général au particulier}.
    \item \textbf{Figures :} « corps luisent de sueur » (métaphore de la lumière/beauté), « pots immobiles » (prouesse technique $\rightarrow$ symbolique de l'équilibre traditionnel), Sidi « plus belle, plus fière » (comparatif de supériorité $\rightarrow$ hiérarchisation visuelle).
\end{itemize}
\end{exemplebox}

\begin{aretenir}[Progression thématique dans la description (Complément Bac)]
Une description efficace suit souvent une \textbf{progression thématique} :
\begin{enumerate}
    \item \textbf{Plan d'ensemble (général) :} cadre, atmosphère, silhouette globale. « Le village s'étendait au creux de la vallée. »
    \item \textbf{Plans intermédiaires (découpage spatial) :} zones, bâtiments, groupes. « Au centre, le marché bourdonnait. Autour, les cases rondes... »
    \item \textbf{Plan rapproché (particulier / détail signifiant) :} visages, objets, textures, geste précis. « Sur l'étal, une mangue mûre luisait. Une main ridée la saisit. »
    \item \textbf{Plan symbolique / subjectif (souvent final) :} ce que le détail révèle de l'âme du lieu ou du personnage. « Cette mangue, c'était le soleil du village dans la main d'une grand-mère. »
\end{enumerate}
Cette progression (général $\rightarrow$ particulier $\rightarrow$ symbolique) guide le regard du lecteur et construit le sens. Soyez capable de l'identifier et de la reproduire.
\end{aretenir}

\begin{savaistu}[Description et enjeu colonial dans \emph{Le Lion et la perle}]
La description chez Soyinka n'est jamais neutre. Le village d'Ilujinle est décrit comme un \textbf{espace ordonné, sacré, communautaire} (arbre sacré au centre, palais, marché, route vers la brousse). La modernité (Lakunle, l'école, la moto, le chemin de fer évoqué) apporte une \textbf{description de rupture} : lignes droites, bruit mécanique, lumière artificielle, individualisme. Repérer les champs lexicaux opposés (tradition : cercle, arbre, parole, danse, sueur, communauté / modernité : ligne, machine, écriture, vitesse, huile, individu) dans les passages descriptifs, c'est déjà faire de l'analyse littéraire \textbf{et} de l'argumentation pour votre dissertation.
\end{savaistu}

\courssub{Application intégrée : dissertation, description et champs lexicaux/sémantiques}

\begin{methode}[Mobiliser champs lexicaux, champs sémantiques et description dans la dissertation]
\begin{enumerate}
    \item \textbf{À l'analyse du sujet :} Identifiez les \textbf{mots-clés} du sujet. Construisez pour chacun son \textbf{champ lexical} (pour varier le vocabulaire) et son \textbf{champ sémantique} (pour en explorer les nuances, les connotations, les tensions).
    \item \textbf{Au brouillon (rédaction du développement) :}
        \begin{itemize}
            \item Utilisez le \textbf{champ lexical} pour \textbf{chaîner vos arguments} sans répétition (cohérence lexicale = cohérence thématique).
            \item Mobilisez le \textbf{champ sémantique} pour \textbf{nuancer votre thèse} : distinguez « tradition » (coutume vivante) de « traditionalisme » (repli figé) ; « modernité » (progrès technique) de « westernisation » (aliénation culturelle).
            \item Insérez des \textbf{touches descriptives} (une scène, un décor, un geste) pour \textbf{incarner l'abstrait} : au lieu de dire « la tradition est vivante », décrivez « la danse des filles sous l'odàn, pots immobiles sur la tête, rythme du tam-tam qui s'emballe ».
        \end{itemize}
    \item \textbf{À la relecture :} Traquez les répétitions (remplacez par champ lexical), vérifiez la justesse des sens (champ sémantique), assurez-vous que les passages descriptifs ont un \textbf{cadre}, un \textbf{thème}, des \textbf{notations sensorielles variées}, une \textbf{progression logique} et une \textbf{image forte} en clôture.
\end{enumerate}
\end{methode}

\begin{exoresolu}[Sujet de dissertation type : « La tradition est-elle un frein au progrès ? »]
\textbf{Étape 1 : Analyse lexicologique / sémantique des mots-clés.}

\begin{center}
\begin{tabularx}{\linewidth}{|l|X|X|}
\hline
\textbf{Mot-clé} & \textbf{Champ lexical (mots pour varier)} & \textbf{Champ sémantique (nuances à exploiter)} \\
\hline
Tradition & coutume, héritage, legs, patrimoine, mémoire, ancêtres, rites, transmission, identité, racines, fondements & [+CONTINUITÉ], [+IDENTITAIRE], [+STRUCTURANT], [+VIVANT] vs [+FIGÉ], [+OBSOLÈTE], [+EXCLUANT], [+OPPRESSIF] \\
\hline
Progrès & évolution, développement, innovation, modernisation, essor, avancée, amélioration, émancipation, rupture, conquête & [+LIBÉRATEUR], [+AMÉLIORATIF], [+RATIONNEL], [+UNIVERSEL] vs [+DÉSTRUCTEUR], [+ALIÉNANT], [+INÉGALITAIRE], [+DÉRACINANT] \\
\hline
Frein & obstacle, entrave, barrage, résistance, inertie, poids, chaîne, limite, garde-fou, ancrage & [+NÉGATIF : blocage], [+POSITIF : protection, stabilité, repère, sélection] \\
\hline
\end{tabularx}
\end{center}

\textbf{Étape 2 : Esquisse de paragraphe avec description intégrée (Thèse : La tradition comme garde-fou nécessaire).}

« Si la modernité apporte l'\emph{innovation} technique et l'\emph{émancipation} individuelle, elle charrie aussi son lot de \emph{ruptures} brutales : le \emph{déracinement} des jeunes, la \emph{marchandisation} des relations, la \emph{perte de repères} moraux. \textbf{La tradition, loin d'être un simple frein, agit comme un garde-fou.} \emph{Description :} « Sur la place d'Ilujinle, au crépuscule, les anciens s'assoient en cercle sous l'odàn. Leurs voix graves tissent la trame des récits fondateurs. Le feu crépite, projetant des ombres dansantes sur les visages ridés. Chaque proverbe lancé est une ancre jetée dans le flot du changement. » Ici, le \textbf{cadre} (place, crépuscule, arbre sacré), le \textbf{thème} (cercle des anciens / parole), les \textbf{notations} (visages ridés, feu, ombres, voix graves), la \textbf{progression} (espace $\rightarrow$ parole $\rightarrow$ symbole de l'ancre) et l'\textbf{image finale} (ancre dans le flot) \textbf{incarnent} l'argument : la tradition stabilise sans immobiliser. Le \textbf{champ sémantique} de l'\emph{ancre} ([+STABILITÉ], [+CHOIX], [+MOBILITÉ MAÎTRISÉE]) replace le « frein » sous un jour positif. »

\textbf{Étape 3 : Transition vers paragraphe suivant (antithèse).}
« Pourtant, quand l'\emph{ancrage} se mue en \emph{enfermement}, la tradition devient le \emph{poids} qui empêche l'\emph{essor}. » (Champ lexical : \emph{enfermement, poids, entrave, rigidité, fermeture} ; champ sémantique : [+FIGÉ], [+EXCLUANT], [+VIOLENCE SYMBOLIQUE] — ex. : le sort des femmes stériles, le refus de l'école pour les filles).
\end{exoresolu}

\begin{exercice}[Entraînement : Repérage et réécriture]
\textbf{Extrait (inventé, style Soyinka) :} « Le marché d'Ilujinle bruissait. Les femmes criaient. Les épices sentaient fort. Les tissus étaient colorés. Les enfants couraient. Le soleil tapait. C'était la vie. »
\begin{enumerate}
    \item Identifiez : cadre, thème, notations (classées par sens), expansion (logique), figures.
    \item Relevez le champ lexical du \emph{marché} et de la \emph{vie/animation}.
    \item Proposez une réécriture descriptive enrichie (progression général $\rightarrow$ particulier, notations sensorielles variées, une image forte finale, vocabulaire varié via champs lexicaux).
\end{enumerate}
\end{exercice}

\begin{corrige}[Exercice n°1]
\textbf{1. Analyse de l'extrait original :}
\begin{itemize}
    \item \textbf{Cadre :} Marché d'Ilujinle, plein jour (soleil).
    \item \textbf{Thème :} L'animation du marché / la vie villageoise.
    \item \textbf{Notations :} Auditif (\emph{bruissait, criaient}), Olfactif (\emph{sentaient fort}), Visuel (\emph{colorés, couraient, soleil tapait}), Kinesthésique (\emph{couraient}). Faible variété, vocabulaire générique.
    \item \textbf{Expansion :} Liste plate (parataxe), pas de progression spatiale ni de focalisation.
    \item \textbf{Figures :} Aucune.
\end{itemize}
\textbf{2. Champs lexicaux :}
\begin{itemize}
    \item Marché : étal, marchand, cliente, négocier, prix, mesure, panier, balance, criée, foule, allée, ombre, toit de paille.
    \item Vie / Animation : fourmiller, bourdonner, grouiller, vibrer, pulser, foisonner, bigarrer, chatoyer, clamer, marchander, éclater, jaillir.
\end{itemize}
\textbf{3. Réécriture enrichie :}
« Le marché d'Ilujinle \textbf{bourdonnait} sous le zénit. \textbf{Plan large :} Une mer de toits de paille bigarrés couvrait la place, d'où s'élevait une \textbf{houle de voix} mêlant le \emph{cliquetis des calebasses}, le \emph{chant nasillard des griots} et le \emph{parfum entêtant du piment et du poisson fumé}. \textbf{Plan intermédiaire :} Dans l'allée centrale, \textbf{Mama Nkechi} déployait ses \emph{pagnes teints à l'indigo}, leurs motifs géométriques \emph{chatoyants} sous le soleil cruel. \textbf{Plan rapproché :} Ses doigts experts \emph{pesaient le mil} dans une balance en bois usée, le grain doré \emph{coulant comme de l'or} dans le sac de jute. \textbf{Image finale :} Un enfant, \emph{rire aux dents blanches}, \emph{chapardait une mangue mûre} — \emph{le soleil du village dans sa petite main}. »
\begin{itemize}
    \item \textbf{Progression :} Large (mer de toits, houle de voix) $\rightarrow$ Intermédiaire (Mama Nkechi, pagnes) $\rightarrow$ Rapproché (doigts, grain de mil) $\rightarrow$ Symbolique (mangue = soleil).
    \item \textbf{Notations :} Visuel (bigarré, chatoyant, doré, blanc), Auditif (bourdonnait, houle, cliquetis, chant), Olfactif (parfum entêtant), Tactile (doigts experts, grain coulant, peau de mangue).
    \item \textbf{Champs lexicaux exploités :} Marché (étal $\rightarrow$ allée, balance, calebasse, pagne, mil, jute), Vie (bourdonner, houle, chatoyer, couler, jaillir, rire, chaparder).
    \item \textbf{Figure forte :} Métaphore \emph{« grain doré coulant comme de l'or »} (valeur, richesse), Métaphore finale \emph{« le soleil du village dans sa main »} (vitalité, avenir, préciosité).
\end{itemize}
\end{corrige}

\begin{attention}[Derniers pièges avant l'intégration finale]
\begin{itemize}
    \item \textbf{Accumulation $\neq$ Description :} Une liste de noms sans organisation spatiale, sans point de vue, sans image, c'est un inventaire, pas une description. Structurez (cadre $\rightarrow$ thème $\rightarrow$ plans $\rightarrow$ détail signifiant).
    \item \textbf{Champ lexical $\neq$ Dictionnaire :} N'alignez pas des synonymes au hasard. Chaque mot choisi doit porter une \textbf{nuance argumentative} (connotation) adaptée à votre thèse du moment.
    \item \textbf{Subjectivité assumée :} En dissertation, la description n'est jamais « objective ». Elle sert votre démonstration. Assumez votre \textbf{point de vue} (énonciation) via le choix des notations et des figures.
    \item \textbf{Cohérence globale :} Les champs lexicaux/sémantiques mobilisés dans l'introduction, le développement et la conclusion doivent \textbf{résonner ensemble}. Créez une \textbf{unité lexicale et sémantique} sur tout le devoir.
\end{itemize}
\end{attention}

Cette section vous a doté des outils lexicaux (champ lexical pour varier, champ sémantique pour nuancer) et stylistiques (texte descriptif pour incarner) indispensables pour réussir l'\textbf{intégration finale} : la rédaction d'une dissertation complète, cohérente, riche et persuasive.

\coursec{Intégration : rédiger une dissertation complète}

\courssub{Transition et objectifs de la séance d'intégration}

Après avoir travaillé séparément chaque composante de la dissertation — analyse du sujet, introduction, paragraphes argumentatifs avec citations et transitions, conclusion, sans oublier les outils linguistiques (types de phrases, dénotation/connotation, énonciation, champs lexical et sémantique, texte descriptif) — nous arrivons à l'étape décisive : \textbf{mobiliser l'ensemble de ces savoirs en situation d'examen}. Cette séance d'intégration reproduit les conditions réelles du Probatoire : un sujet inédit, un temps limité (2~heures), une copie à rendre structurée, argumentée, bien écrite.

\begin{aretenir}[Objectifs de la séance d'intégration]
\begin{itemize}
    \item Savoir analyser un sujet de dissertation en 15~minutes chrono (mots-clés, définitions, problématique, plan).
    \item Produire une introduction et une conclusion conformes aux attendus (amorce, sujet posé, sujet divisé, thèse, plan annoncé~; réouverture, bilan, ouverture).
    \item Rédiger un développement de 2 ou 3 parties, chacune découpée en 2 ou 3 paragraphes, avec arguments, exemples, citations insérées, transitions fluides.
    \item Gérer son temps : analyse (15~min) + plan détaillé (10~min) + rédaction (60~min) + relecture (15~min).
    \item S'auto-évaluer grâce à une grille critères et corriger ses points faibles en ateliers de remédiation ciblés.
    \item Transférer la compétence «~argumenter par écrit~» à la vie courante : lettre de motivation, exposé oral, débat structuré.
\end{itemize}
\end{aretenir}

\courssub{Étape 1 — Analyser le sujet proposé (15~minutes)}

L'analyse du sujet n'est pas une formalité : c'est le \textbf{socle de toute la dissertation}. Un sujet mal analysé entraîne un hors-sujet, un plan inadapté, des arguments hors propos. La méthode en trois temps (mots-clés $\rightarrow$ problématique $\rightarrow$ plan) est non négociable.

\begin{methode}[Analyser un sujet de dissertation — 3 temps, 15~minutes]
\begin{enumerate}
    \item \textbf{Repérer et définir les mots-clés (5~min).} Soulignez les termes centraux du sujet. Donnez une définition précise (dictionnaire + contexte du sujet). Identifiez les relations logiques (opposition, cause/effet, concession, finalité). Exemple~: «~La tradition est-elle un frein au progrès~?~» $\rightarrow$ \cle{tradition} (héritage culturel, coutumes), \cle{frein} (obstacle, ralentissement), \cle{progrès} (avancement technique, social, moral). Relation~: opposition apparente.
    \item \textbf{Formuler la problématique (5~min).} Transformez le sujet en question ouverte, débattue, qui appelle une réponse nuancée (thèse + antithèse + synthèse). Évitez les questions fermées (oui/non) ou trop larges. Bonne problématique~: «~Dans quelle mesure la tradition peut-elle à la fois entraver et nourrir le progrès~?~» Mauvaise~: «~La tradition freine-t-elle le progrès~?~» (trop binaire).
    \item \textbf{Construire le plan détaillé (5~min).} Choisissez la structure adaptée à la problématique~: dialectique (thèse/antithèse/synthèse), analytique (causes/conséquences/solutions), thématique (aspects sociaux/culturels/économiques). Pour chaque grande partie (I, II, III), notez 2 à 3 sous-parties (A, B, C) avec l'argument principal, l'exemple/l'œuvre/citation prévu(e), le mot de transition vers le suivant. Le plan doit tenir sur une feuille de brouillon, lisible, complet.
\end{enumerate}
\end{methode}

\begin{attention}[Piège fréquent — Recopier le sujet ou le paraphraser en guise d'introduction]
Ne confondez \textbf{pas} l'analyse du sujet (travail de brouillon) et l'introduction (texte final). L'introduction suit un canevas précis (amorce, sujet posé, sujet divisé, thèse, plan annoncé). Recopier le sujet mot pour mot = 0~point à l'introduction. Paraphraser sans problématiser = introduction «~plate~», note faible. \textbf{Sur le brouillon, vous analysez~; sur la copie, vous rédigez l'introduction.}
\end{attention}

\begin{exemplebox}[Analyse guidée d'un sujet type Probatoire]
\textbf{Sujet~:} «~Le théâtre est-il le miroir de la société~?~»

\textbf{1. Mots-clés~:}
\begin{itemize}
    \item \cle{Théâtre}~: genre littéraire et spectacle vivant, double dimension (texte + représentation).
    \item \cle{Miroir}~: métaphore $\rightarrow$ reflet, imitation, mais aussi déformation, grossissement, mise en lumière.
    \item \cle{Société}~: ensemble d'individus partageant culture, normes, conflits, histoire.
    \item Relation~: correspondance / représentation / critique.
\end{itemize}

\textbf{2. Problématique~:} «~En quoi le théâtre, par sa double nature de texte et de spectacle, offre-t-il une réflexion critique sur la société qui le produit, au-delà du simple reflet~?~»

\textbf{3. Plan dialectique (adapté)~:}
\begin{itemize}
    \item \textbf{I. Le théâtre comme reflet fidèle des réalités sociales} (thèse)
        \begin{itemize}
            \item A. La représentation des mœurs et des conflits de classe (ex.~: \emph{Le Misanthrope}, \emph{Le Lion et la Perle} — Baroka vs Lakunle).
            \item B. La mise en scène des rapports de pouvoir et de genre (ex.~: Sidi, Sadiku, la place de la femme).
            \item C. Le théâtre comme archive vivante de l'histoire collective.
        \end{itemize}
    \item \textbf{II. Le théâtre comme déformation volontaire~: distanciation et critique} (antithèse)
        \begin{itemize}
            \item A. L'exagération comique ou tragique pour mieux dénoncer (satire, caricature).
            \item B. Le jeu sur les codes (métathéâtre, aparté, chœur) pour éveiller la conscience du spectateur.
            \item C. L'exemple de Soyinka~: tradition yoruba revisitée pour questionner la modernité.
        \end{itemize}
    \item \textbf{III. Le théâtre comme laboratoire d'un avenir possible} (synthèse)
        \begin{itemize}
            \item A. La fonction utopique~: imaginer d'autres organisations sociales.
            \item B. La catharsis et l'engagement du spectateur (théâtre-forum, Brecht).
            \item C. Le théâtre comme espace de dialogue interculturel (tradition/oralité vs modernité/écriture).
        \end{itemize}
\end{itemize}
\end{exemplebox}

\courssub{Étape 2 — Produire l'introduction et la conclusion en temps limité}

L'introduction et la conclusion sont les \textbf{vitrines} de votre copie. Le correcteur les lit en premier et en dernier~: elles conditionnent l'impression globale. Elles doivent être rédigées \textbf{après} le développement (sur brouillon) pour garantir la cohérence thèse/plan/conclusion, mais recopiées \textbf{en premier et en dernier} sur la copie finale.

\begin{methode}[Rédiger l'introduction type — 5~minutes sur la copie finale]
\begin{enumerate}
    \item \textbf{Amorce (1~phrase)~:} Accroche large, culturelle, liée au sujet (citation d'auteur, fait de société, référence à l'œuvre au programme). Ex.~: «~Depuis l'Antiquité, le théâtre interroge la cité qui l'accueille.~»
    \item \textbf{Sujet posé (1~phrase)~:} Reprise du sujet en le reformulant, sans le citer textuellement. Ex.~: «~La question de savoir si le théâtre se borne à refléter la société ou s'il la transforme mérite d'être examinée.~»
    \item \textbf{Sujet divisé / Définitions (1--2~phrases)~:} Définition rapide des termes clés dans l'optique de la problématique. Ex.~: «~Par théâtre, entendons l'art dramatique dans sa double dimension textuelle et spectaculaire~; par miroir, non pas une copie passive mais une représentation sélective.~»
    \item \textbf{Problématique (1~phrase, interrogative)~:} La question centrale qui guide tout le devoir. Ex.~: «~Dans quelle mesure le théâtre, au-delà du reflet, agit-il comme un révélateur critique et un laboratoire d'avenir pour la société~?~»
    \item \textbf{Annonce du plan (1~phrase)~:} «~Nous verrons d'abord que le théâtre reflète les réalités sociales (I), puis qu'il les déforme pour mieux les critiquer (II), enfin qu'il propose des horizons nouveaux (III).~»
\end{enumerate}
\end{methode}

\begin{methode}[Rédiger la conclusion type — 5~minutes sur la copie finale]
\begin{enumerate}
    \item \textbf{Bilan synthétique (2--3~phrases)~:} Reprenez la thèse de chaque grande partie en une phrase, montrez l'enchaînement logique. Ex.~: «~Le théâtre reflète la société (I), mais ce reflet est actif~: il déforme pour dénoncer (II), et finit par ouvrir des possibles (III).~»
    \item \textbf{Réponse à la problématique (1~phrase affirmative)~:} Ex.~: «~Ainsi, le théâtre n'est pas un simple miroir~: il est un opérateur de conscience et de transformation sociale.~»
    \item \textbf{Ouverture (1~phrase)~:} Élargissement pertinent (autre genre, autre époque, art différent, actualité). Ex.~: «~Cette fonction critique se retrouve-t-elle dans le cinéma contemporain ou dans les réseaux sociaux, nouveaux «~théâtres~» de la société~?~»
\end{enumerate}
\end{methode}

\begin{propriete}[Règle d'or — Cohérence introduction / développement / conclusion]
La thèse annoncée en introduction doit être \textbf{exactement} celle que le développement défend et que la conclusion résume. Le plan annoncé doit être \textbf{scrupuleusement respecté} (même ordre, mêmes titres de parties). Toute divergence = incohérence = perte de points importante.
\end{propriete}

\courssub{Étape 3 — Produire le développement (60~minutes)}

Le développement est le \textbf{cœur argumentatif} de la dissertation. Chaque grande partie (I, II, III) comporte 2 à 3 paragraphes (A, B, C). Un paragraphe = une idée force + argumentation + illustration (exemple précis, citation insérée) + mini-transition vers le paragraphe suivant. La transition entre grandes parties (I $\rightarrow$ II, II $\rightarrow$ III) est un paragraphe à part entière~: elle fait le bilan de la partie achevée, annonce la limite qui appelle la suivante, et lance la nouvelle problématique partielle.

\begin{methode}[Structure type d'un paragraphe argumentatif — 8 à 10~minutes]
\begin{enumerate}
    \item \textbf{Phrase d'ouverture (thèse du paragraphe)~:} Affirmation claire, liée à la sous-partie du plan. Ex.~: «~Le théâtre expose d'abord les hiérarchies sociales et les conflits de génération.~»
    \item \textbf{Argumentation (2--3~phrases)~:} Explication, analyse, mise en relation. Ex.~: «~Dans \emph{Le Lion et la Perle}, l'opposition Baroka/Lakunle incarne le choc tradition/modernité. Baroka, le Bale, use de ruse et d'autorité coutumière~; Lakunle, instituteur, prône un progrès occidentalisé.~»
    \item \textbf{Illustration par citation insérée (1--2~phrases)~:} Citation courte, intégrée dans la phrase, commentée. Ex.~: «~Lorsque Baroka déclare~: «~\emph{Je ne suis pas un arbre qu'on abat}~», il affirme la résistance de l'autorité traditionnelle face à la modernité brutale.~»
    \item \textbf{Mini-transition (1~phrase)~:} Lien logique vers le paragraphe suivant. Ex.~: «~Mais au-delà des personnages, c'est la place des femmes que la pièce met en lumière.~»
\end{enumerate}
\end{methode}

\begin{exemplebox}[Paragraphe complet — Partie I A~: «~La représentation des mœurs et des conflits de classe~»]
Le théâtre expose d'abord les hiérarchies sociales et les conflits de génération. Dans \emph{Le Lion et la Perle}, l'opposition entre Baroka et Lakunle incarne le choc entre tradition et modernité au village d'Ilujinle. Baroka, le Bale, use de ruse, de proverbes et d'autorité coutumière pour maintenir son pouvoir~; Lakunle, jeune instituteur formé à l'occidentale, prône un progrès technique et moral qui méprise les coutumes. Lorsque Baroka déclare~: «~\emph{Je ne suis pas un arbre qu'on abat}~», il affirme la résistance de l'autorité traditionnelle face à la modernité brutale, tout en révélant sa propre lucidité sur la nécessité d'adaptation. Cette confrontation n'est pas seulement individuelle~: elle met en scène la tension qui traverse toute la société yoruba des années 1960. Mais au-delà des personnages masculins, c'est la place des femmes que la pièce met en lumière.
\end{exemplebox}

\begin{methode}[Rédiger une transition entre grandes parties — 3~minutes]
\begin{enumerate}
    \item \textbf{Bilan de la partie achevée~:} «~En définitive, le théâtre reflète fidèlement les structures sociales et les conflits qui les animent.~»
    \item \textbf{Limite / Pivot~:} «~Pourtant, ce reflet n'est jamais neutre~: le dramaturge sélectionne, grossit, déforme pour provoquer une prise de conscience.~»
    \item \textbf{Annonce de la partie suivante~:} «~C'est cette fonction critique, voire subversive, que nous allons examiner à présent.~»
\end{enumerate}
\end{methode}

\begin{attention}[Erreurs fréquentes dans le développement]
\begin{itemize}
    \item \textbf{Citation «~parachutée~»~:} citation collée sans introduction ni commentaire. \emph{Solution~:} toujours «~Selon X, "..."~; cette image montre que...~»
    \item \textbf{Paragraphe «~liste de courses~»~:} succession d'exemples sans argument liant. \emph{Solution~:} une idée force par paragraphe, les exemples la servent.
    \item \textbf{Transition absente ou mécanique~:} «~Passons maintenant à la deuxième partie.~» \emph{Solution~:} transition logique (mais / toutefois / cependant / par conséquent).
    \item \textbf{Hors-sujet dans l'illustration~:} exemple qui ne prouve pas l'argument. \emph{Solution~:} vérifier sur le brouillon~: cet exemple répond-il à la thèse du paragraphe~?
\end{itemize}
\end{attention}

\courssub{Gestion du temps — Répartition indicative sur 2~heures}

\begin{popfigure}
\centering
\begin{tikzpicture}
\begin{axis}[
    ybar,
    bar width=18pt,
    width=\linewidth,
    height=9cm,
    ymin=0, ymax=70,
    ylabel={Minutes},
    symbolic x coords={Analyse, Plan, Rédaction, Relecture},
    xtick=data,
    nodes near coords,
    nodes near coords align={vertical},
    every node near coord/.append style={font=\small, popdark},
    ymajorgrids=true,
    grid style={popline},
    axis lines*=left,
    enlarge x limits=0.25,
    tick label style={font=\small, popdark},
    label style={font=\small, popdark},
    bar shift=0pt,
]
\addplot[popblue, fill=popblue!80] coordinates {(Analyse,15) (Plan,10) (Rédaction,60) (Relecture,15)};
\end{axis}
\end{tikzpicture}
\legende{Répartition idéale du temps pour une dissertation de 2~heures (120~minutes). La rédaction (60~min) inclut introduction, développement, conclusion. La relecture (15~min) est non négociable.}
\end{popfigure}

\begin{methode}[Gérer son temps à l'examen — Règles non négociables]
\begin{enumerate}
    \item \textbf{Ne jamais rédiger sans plan validé.} Les 25~premières minutes (analyse + plan) sont un investissement~: elles évitent le hors-sujet, les répétitions, les impasses.
    \item \textbf{Rédigez sur brouillon les phrases-clés~:} thèse de chaque paragraphe, citation choisie, transition. Puis recopiez proprement sur la copie en développant.
    \item \textbf{Respectez le quota de temps par partie~:} Intro (5~min), I (15~min), II (15~min), III (15~min), Conclusion (5~min), Relecture (15~min) = 70~min rédaction + relecture. Si vous dépassez sur une partie, accélérez~: mieux vaut une partie un peu courte qu'une partie manquante.
    \item \textbf{Prévoyez 10~minutes minimum de relecture} (orthographe, accords, ponctuation, cohérence des transitions, respect du plan annoncé). C'est souvent la différence entre 12 et 14/20.
    \item \textbf{Si le temps manque~:} sacrifiez l'ouverture de la conclusion, pas la relecture. Une copie sans fautes graves et cohérente vaut mieux qu'une copie longue mais truffée d'erreurs.
\end{enumerate}
\end{methode}

\courssub{Grille d'auto-évaluation — Se corriger soi-même}

Avant de rendre la copie (ou lors de la relecture), cochez chaque critère. Cette grille est calée sur les attendus du Probatoire / Baccalauréat technique.

\begin{center}
\begin{tabularx}{\linewidth}{|l|X|c|}\hline
\textbf{Critère} & \textbf{Indicateurs de réussite} & \textbf{Auto-éval} \\ \hline
Analyse du sujet & Mots-clés définis, problématique pertinente, plan adapté & $\square$ \\ \hline
Introduction & Amorce + sujet posé + divisé + problématique + plan annoncé (5 étapes) & $\square$ \\ \hline
Développement : structure & 3 parties (I, II, III), 2--3 sous-parties chacune, transitions inter-parties & $\square$ \\ \hline
Développement : argumentation & 1 idée force / paragraphe, argument explicité, exemple/citation inséré & $\square$ \\ \hline
Citations & Courtes, intégrées, commentées, sources identifiées (œuvre, auteur) & $\square$ \\ \hline
Transitions & Logiques (mais, toutefois, par conséquent, de plus), pas mécaniques & $\square$ \\ \hline
Conclusion & Bilan + réponse à la problématique + ouverture pertinente & $\square$ \\ \hline
Langue & Orthographe, grammaire, syntaxe, vocabulaire précis, phrases variées & $\square$ \\ \hline
Présentation & Paragraphes visibles (alinea), lisibilité, ratures minimes & $\square$ \\ \hline
Temps & Copie rendue dans les temps, relecture effectuée & $\square$ \\ \hline
\end{tabularx}
\end{center}

\begin{aretenir}[Seuil de réussite — Probatoire / Bac Tech]
\begin{itemize}
    \item \textbf{10/20~:} Sujet traité, plan visible, arguments présents, langue compréhensible malgré fautes.
    \item \textbf{12--14/20~:} Problématique claire, plan dialectique/analytique maîtrisé, citations pertinentes, transitions fluides, peu de fautes.
    \item \textbf{16+/20~:} Problématique fine, synthèse originale, citations analysées (pas seulement citées), style élégant, ouverture pertinente, zéro faute.
\end{itemize}
\end{aretenir}

\courssub{Compte-rendu et remédiation — Ateliers ciblés}

Après la séance d'intégration (ou un devoir maison type examen), le professeur organise une \textbf{correction collective d'une copie anonyme} (choisie pour sa représentativité~: points forts + erreurs types). Cette séance de «~compte-rendu~» suit un protocole précis.

\begin{methode}[Déroulement du compte-rendu collectif — 1 séance]
\begin{enumerate}
    \item \textbf{Distribution de la copie anonyme} (photocopiée ou projetée) + grille d'évaluation vierge.
    \item \textbf{Lecture silencieuse individuelle (5~min)~:} chaque élève note 3 points forts et 3 faiblesses sur la grille.
    \item \textbf{Mise en commun orale (15~min)~:} le professeur guide~: «~Qu'a-t-on aimé dans l'introduction~? Pourquoi cette transition fonctionne-t-elle~? Où le hors-sujet pointe-t-il~?~»
    \item \textbf{Correction guidée par le prof (15~min)~:} sur le tableau, réécriture d'un paragraphe faible, insertion d'une citation manquante, correction d'une transition mécanique.
    \item \textbf{Attribution d'une note collective (5~min)~:} consensus argumenté sur la note /20 selon la grille officielle.
    \item \textbf{Remédiation individuelle (prochaine séance)~:} chaque élève reçoit sa propre copie corrigée + fiche de remédiation personnalisée.
\end{enumerate}
\end{methode}

\begin{experience}[Ateliers de remédiation ciblés — 4 stations, 20~min chacune]
\begin{center}
\begin{tabularx}{\linewidth}{|l|X|}\hline
\textbf{Atelier} & \textbf{Activité ciblée} \\ \hline
A — Introduction & Réécrire 3 introductions «~plates~» en appliquant le canevas 5 étapes. Échange pair-à-pair + validation prof. \\ \hline
B — Citations & Exercice «~citation orpheline $\rightarrow$ citation intégrée~»~: 5 citations données, à insérer dans des phrases avec commentaire. \\ \hline
C — Orthographe / Grammaire & Chasse aux fautes récurrentes (accords participes, homophones, ponctuation) sur sa propre copie. Fiche mémo personnelle. \\ \hline
D — Transitions & Transformer 5 transitions mécaniques («~Ensuite...~», «~Pour finir...~») en transitions logiques (concession, cause, conséquence). \\ \hline
\end{tabularx}
\end{center}
\legende{Organisation en îlots~: les élèves tournent sur les 4 ateliers en 2 séances (2$\times$40~min). Le professeur circule, valide, renforce.}
\end{experience}

\courssub{Application à la vie courante — Transférer la compétence argumentative}

La dissertation n'est pas un exercice scolaire enfermé~: c'est un \textbf{entraînement à la pensée critique et à la communication persuasive}. Toute situation où il faut convaincre, justifier, négocier, expose la même structure~: thèse $\rightarrow$ arguments $\rightarrow$ exemples $\rightarrow$ conclusion.

\begin{exemplebox}[Trois situations de la vie courante = structure de dissertation]
\begin{center}
\begin{tabularx}{\linewidth}{|l|X|X|X|}\hline
\textbf{Élément dissertation} & \textbf{Lettre de motivation (emploi / stage)} & \textbf{Exposé oral (projet technique)} & \textbf{Débat / Négociation} \\ \hline
Amorce / Accroche & «~Actuellement en 1re Tech Maintenance...~» & «~Le défi énergétique de notre région...~» & «~Nous sommes tous concernés par...~» \\ \hline
Sujet posé / Problématique & «~Pourquoi ma candidature correspond à vos besoins~?~» & «~Comment réduire la consommation de 20\%~?~» & «~Faut-il accepter cette proposition~?~» \\ \hline
Plan annoncé / Thèse & «~Je présenterai mes compétences, mon expérience, ma motivation.~» & «~Diagnostic, solutions techniques, planning, budget.~» & «~Avantages, risques, conditions, contre-proposition.~» \\ \hline
Développement (arguments + preuves) & Compétences (diplômes, stages) + exemples concrets (réparation moteur) & Données techniques, chiffrées, schémas, retours d'expérience & Faits, témoignages, calculs, précédents \\ \hline
Transitions & «~Par ailleurs, mon stage m'a permis de...~» & «~Passons maintenant au volet financier.~» & «~Certes, mais avez-vous considéré...~?~» \\ \hline
Conclusion / Appel à l'action & «~Je souhaite vous convaincre en entretien.~» & «~Je propose de lancer l'étude dès la semaine prochaine.~» & «~Si vous acceptez ces conditions, signez ici.~» \\ \hline
\end{tabularx}
\end{center}
\end{exemplebox}

\begin{methode}[Justifier sa démarche — Explicable à l'oral (épreuve orale, entretien)]
\begin{enumerate}
    \item \textbf{Pourquoi ce plan~?} «~J'ai choisi le plan dialectique car le sujet pose une opposition (tradition/progrès) qui appelle une nuance, pas un choix binaire.~»
    \item \textbf{Pourquoi cet exemple / cette citation~?} «~J'ai cité Baroka dans \emph{Le Lion et la Perle} car il incarne la complexité~: il défend la tradition mais utilise la ruse moderne (le timbre-poste). C'est l'exemple parfait de la synthèse.~»
    \item \textbf{Pourquoi cette transition~?} «~J'ai utilisé "toutefois" pour marquer le passage du reflet (partie I) à la déformation critique (partie II)~: c'est une concession qui prépare l'antithèse.~»
    \item \textbf{Comment as-tu géré ton temps~?} «~J'ai passé 25~min sur le brouillon (analyse + plan détaillé), 60~min à recopier en développant, 15~min à relire. J'ai fini à 11h55.~»
\end{enumerate}
\end{methode}

\begin{savaistu}[Le savais-tu~? — Hors-sujet et sanction au Bac Technique]
Au Baccalauréat technique (séries F, G, H, etc.), le \textbf{hors-sujet est sanctionné par une note $\le$ 5/20}, quelle que soit la qualité de la langue, de l'orthographe ou de la culture générale affichée. Le correcteur applique la grille~: «~Hors sujet = 0 à 5 sur 20~». \textbf{Aucune compensation possible.} D'où l'importance vitale des 15~minutes d'analyse et de la vérification \textbf{avant} de commencer la rédaction~: «~Mon plan répond-il \emph{exactement} à la problématique~?~»
\end{savaistu}

\begin{propriete}[Règle d'or — Toute affirmation justifiée, illustrée, sourcée]
Dans une dissertation de niveau Probatoire / Bac Tech, \textbf{aucune affirmation ne doit rester nue}.
\begin{itemize}
    \item \textbf{Argument~:} explication logique, cause/effet, mécanisme.
    \item \textbf{Illustration~:} exemple précis (fait historique, passage d'œuvre, statistique, expérience vécue).
    \item \textbf{Source / Citation~:} référence explicite (auteur, œuvre, acte/scène/vers, page si possible).
\end{itemize}
Schéma~: \textbf{Affirmation $\rightarrow$ Argument $\rightarrow$ Exemple/Citation $\rightarrow$ Commentaire $\rightarrow$ Mini-transition.}
\end{propriete}

\begin{exoresolu}[Sujet complet — Dissertation type 2h (simulation)]
\textbf{Sujet~:} «~La tradition est-elle un frein au progrès~?~» (Durée~: 2h)

\tcblower

\textbf{1. Analyse (15~min) — Brouillon}
Mots-clés~: Tradition (héritage, coutumes, oralité, identité), Frein (obstacle, ralentissement, mais aussi garde-fou), Progrès (technique, scientifique, social, moral). Relation~: tension, dialectique.
Problématique~: «~Dans quelle mesure la tradition, loin d'être un simple obstacle, peut-elle être réinventée comme ressource pour un progrès durable~?~»
Plan dialectique~:
I. Tradition comme frein~: immobilisme, exclusion, résistance au changement (ex.~: excision, castes, refus école).
II. Tradition comme garde-fou et réservoir de sens~: éthique, écologie, cohésion sociale (ex.~: palabre, médecine traditionnelle, valeurs Ubuntu).
III. Tradition réinventée~: modernité endogène, hybridation féconde (ex.~: Soyinka, architecture néo-traditionnelle, start-up africaines).

\textbf{2. Introduction (5~min — copie finale)}
«~"La tradition n'est pas le culte des cendres, mais la transmission du feu", disait Gustav Mahler. Cette parole éclaire le débat qui oppose souvent tradition et progrès comme deux forces ennemies. La tradition, entendue comme l'ensemble des coutumes, savoirs et valeurs transmis par une communauté, est-elle condamnée à freiner le progrès, entendu comme l'amélioration continue des conditions de vie~? Ou bien, paradoxalement, peut-elle en être le ferment~? Nous analyserons d'abord les formes d'immobilisme que la tradition peut générer (I), puis son rôle de régulateur éthique et social (II), enfin les voies d'une modernité enracinée qui réinvente l'héritage (III).~»

\textbf{3. Développement (60~min) — Extraits clés}
\textit{Partie I A~:} «~La tradition freine le progrès lorsqu'elle sacralise l'ordre établi. Dans \emph{Le Lion et la Perle}, le refus de Lakunle de payer la dot illustre le choc~: la coutume yoruba (dot = alliance entre familles) perçue comme "achat de la femme" par le moderne. Baroka exploite cette tradition pour asseoir son pouvoir personnel.~»
\textit{Transition I$\rightarrow$II~:} «~Cependant, réduire la tradition à un frein ignore sa fonction de cohésion et de régulation.~»
\textit{Partie II B~:} «~La palabre, institution traditionnelle de parole réglée, offre un modèle de résolution des conflits que les tribunaux modernes peinent à égaler en acceptabilité sociale. Amadou Hampâté Bâ écrit~: "La tradition est la mémoire de l'avenir."~»
\textit{Transition II$\rightarrow$III~:} «~Mais la tradition ne doit pas être figée~: elle gagne à dialoguer avec la modernité.~»
\textit{Partie III C~:} «~L'architecte Diébédo Francis Kéré (Prix Pritzker 2022) construit des écoles au Burkina en terre crue (technique traditionnelle) avec ventilation passive (ingénierie moderne). Exemple type de "modernité endogène".~»

\textbf{4. Conclusion (5~min)}
«~La tradition n'est pas un frein par nature~: elle devient obstacle quand elle se fige en dogme, mais ressource quand elle est interrogée. Le progrès durable ne s'affranchit pas de l'héritage~; il le réinvente. Le théâtre de Soyinka, l'architecture de Kéré, la médecine intégrée le prouvent~: l'avenir se construit \emph{avec} la mémoire, pas \emph{contre} elle. \textbf{Ouverture~:} Cette dialectique tradition/progrès se retrouve-t-elle dans la gestion des données numériques face aux savoirs oraux~?~»

\textbf{5. Relecture (15~min) — Check-list}
Accords (participes, sujets inversés), homophones (à/a, et/est, ou/où), ponctuation (virgules avant "mais", "car", "donc"), paragraphes visibles (alinea), citations entre guillemets, noms d'œuvres en italique, plan respecté (I, II, III), thèse cohérente.
\end{exoresolu}

\courssub{Exercices d'entraînement — Vers l'autonomie}

\begin{exercice}[Entraînement 1 — Analyse express (10~min chrono)]
Analysez le sujet suivant comme en examen~: mots-clés, définitions, problématique, plan détaillé (I, II, III avec A, B, C).
\textbf{Sujet~:} «~L'école peut-elle former le citoyen sans transmettre la culture nationale~?~»
\end{exercice}

\begin{exercice}[Entraînement 2 — Introduction minute (5~min chrono)]
Rédigez l'introduction complète (5 étapes) pour le sujet de l'exercice 1. Chronométrez-vous.
\end{exercice}

\begin{exercice}[Entraînement 3 — Paragraphe unique (10~min)]
Rédigez le paragraphe I A du plan que vous avez proposé~: «~L'école transmet d'abord un socle commun de connaissances et de valeurs.~» Utilisez une citation (auteur au choix~: Condorcet, Ferry, Freire, ou œuvre au programme). Intégrez-la, commentez-la, faites une mini-transition vers I B.
\end{exercice}

\begin{exercice}[Entraînement 4 — Gestion du temps (simulation 45~min)]
En 45~minutes~: analyse (10~min) + plan (5~min) + introduction + partie I complète (2 paragraphes + transition) + conclusion. Sujet~: «~Le numérique menace-t-il la littérature~?~» Auto-évaluez avec la grille.
\end{exercice}

\begin{corrige}[Exercice 1 — Analyse express (proposition)]
\textbf{Mots-clés~:} École (institution, lieu d'apprentissage, socialisation), Citoyen (sujet de droits/devoirs, acteur de la cité), Culture nationale (héritage littéraire, historique, artistique, linguistique, valeurs partagées). Transmettre = passer, léguer, faire approprié.
\textbf{Problématique~:} «~L'école doit-elle choisir entre formation civique universelle et transmission d'un patrimoine national, ou peut-elle articuler les deux pour former un citoyen enraciné et ouvert~?~»
\textbf{Plan analytique~:}
I. L'école comme creuset du citoyen universel (valeurs républicaines, droits de l'homme, esprit critique) — A. Décloisonnement / universalisme~; B. Éducation à la laïcité, à la démocratie~; C. Risque~: déracinement, uprooting.
II. La culture nationale comme condition de l'appartenance et de l'identité — A. Langue, littérature, histoire partagée~; B. Mémoire collective, rites républicains~; C. Risque~: enfermement identitaire, exclusion.
III. Articulation~: le citoyen "raciné et ouvert" — A. Programme~: littérature nationale + littérature monde~; B. Histoire~: nationale + mondiale~; C. Pédagogie~: comparatisme, pluralisme, éducation interculturelle.
\end{corrige}

\begin{aretenir}[Bilan de l'unité — Ce que vous devez maîtriser pour le Probatoire]
\begin{itemize}
    \item \textbf{Analyser} un sujet en 15~min (mots-clés, définitions, problématique, plan).
    \item \textbf{Rédiger} une introduction canonique (5 étapes) et une conclusion (bilan + réponse + ouverture).
    \item \textbf{Construire} un développement en 3 parties, 6--9 paragraphes, avec arguments, citations insérées et commentées, transitions logiques.
    \item \textbf{Gérer} 2h~: 25~min brouillon (analyse+plan), 60~min rédaction, 15~min relecture, 20~min marge.
    \item \textbf{Relire} efficacement (orthographe, grammaire, cohérence, respect du plan).
    \item \textbf{Transférer} la compétence argumentative à la lettre de motivation, l'exposé technique, le débat citoyen.
    \item \textbf{Justifier} ses choix (plan, exemples, citations, transitions) à l'oral comme à l'écrit.
\end{itemize}
\end{aretenir}

\begin{popfigure}
\centering
\begin{tikzpicture}[
    node distance=1.2cm and 1.5cm,
    every node/.style={font=\small},
    box/.style={draw=popblue, thick, rounded corners, fill=popblue!10, minimum width=3.2cm, minimum height=1.1cm, align=center},
    arrow/.style={-Stealth, thick, popdark},
    labelbox/.style={draw=poporange, thick, rounded corners, fill=poporange!15, minimum width=4cm, minimum height=0.9cm, align=center, font=\small},
]
% Introduction
\node[box, align=center] (intro){\textbf{INTRODUCTION}\\\textit{(5~min)}\\\textcolor{popdark}{Amorce $\rightarrow$ Sujet posé $\rightarrow$ Sujet divisé $\rightarrow$ Problématique $\rightarrow$ Plan annoncé}};
% Développement - Partie I
\node[box, below=of intro, align=center] (I){\textbf{I. PREMIÈRE PARTIE}\\\textit{(15~min)}\\\textcolor{popdark}{Thèse partielle}};
\node[labelbox, left=of I, xshift=-1.2cm, align=center] (IA){A. Paragraphe 1\\Idée + Argument + Citation + Transition};
\node[labelbox, below=of IA, align=center] (IB){B. Paragraphe 2\\Idée + Argument + Citation + Transition};
\node[labelbox, below=of IB, align=center] (IC){C. Paragraphe 3\\(optionnel)};
\node[box, right=of I, xshift=1.2cm, align=center] (trans12){\textbf{TRANSITION I$\rightarrow$II}\\\textit{(3~min)}\\\textcolor{popdark}{Bilan I + Limite + Annonce II}};
% Développement - Partie II
\node[box, below=of I, align=center] (II){\textbf{II. DEUXIÈME PARTIE}\\\textit{(15~min)}\\\textcolor{popdark}{Antithèse / Approfondissement}};
\node[labelbox, left=of II, xshift=-1.2cm, align=center] (IIA){A. Paragraphe 1\\Idée + Argument + Citation + Transition};
\node[labelbox, below=of IIA, align=center] (IIB){B. Paragraphe 2\\Idée + Argument + Citation + Transition};
\node[labelbox, below=of IIB, align=center] (IIC){C. Paragraphe 3\\(optionnel)};
\node[box, right=of II, xshift=1.2cm, align=center] (trans23){\textbf{TRANSITION II$\rightarrow$III}\\\textit{(3~min)}\\\textcolor{popdark}{Bilan II + Limite + Annonce III}};
% Développement - Partie III
\node[box, below=of II, align=center] (III){\textbf{III. TROISIÈME PARTIE}\\\textit{(15~min)}\\\textcolor{popdark}{Synthèse / Dépassement}};
\node[labelbox, left=of III, xshift=-1.2cm, align=center] (IIIA){A. Paragraphe 1\\Idée + Argument + Citation + Transition};
\node[labelbox, below=of IIIA, align=center] (IIIB){B. Paragraphe 2\\Idée + Argument + Citation + Transition};
\node[labelbox, below=of IIIB, align=center] (IIIC){C. Paragraphe 3\\(optionnel)};
% Conclusion
\node[box, below=of III, align=center] (concl){\textbf{CONCLUSION}\\\textit{(5~min)}\\\textcolor{popdark}{Bilan synthétique $\rightarrow$ Réponse à la problématique $\rightarrow$ Ouverture}};
% Flèches
\draw[arrow] (intro) -- (I);
\draw[arrow] (I) -- (trans12);
\draw[arrow] (trans12) -- (II);
\draw[arrow] (II) -- (trans23);
\draw[arrow] (trans23) -- (III);
\draw[arrow] (III) -- (concl);
% Annotations temps
\node[align=center, popmuted, font=\tiny, right=of trans12, xshift=1.5cm] {Temps total\\rédaction~: 60~min};
\node[align=center, popmuted, font=\tiny, left=of IA, xshift=-1.5cm] {1 paragraphe\\$\approx$ 8--10~min};
\end{tikzpicture}
\legende{Architecture complète de la dissertation de Première technique~: introduction (5 étapes), développement en 3 parties (chaque partie = 2 à 3 paragraphes argumentés + transitions inter-parties), conclusion (3 étapes). Temps indicatifs pour une épreuve de 2~heures.}
\end{popfigure}

\begin{exemplebox}[Chiffres-clés — Dissertation de Première technique]
\begin{itemize}
    \item \textbf{Longueur~:} 2 à 3 pages manuscrites (copie double grand carreau).
    \item \textbf{Structure~:} 3 grandes parties (I, II, III) — 6 à 9 paragraphes au total.
    \item \textbf{Citations~:} 4 à 6 minimum (œuvres au programme + culture personnelle), courtes (1 à 2 vers / 1 phrase), intégrées, commentées.
    \item \textbf{Transitions~:} 2 transitions inter-parties (obligatoires) + 5 à 8 mini-transitions intra-parties.
    \item \textbf{Mots de liaison~:} au minimum 15 connecteurs logiques variés (donc, cependant, par ailleurs, en outre, dès lors, c'est pourquoi, en définitive...).
    \item \textbf{Relecture~:} 15~minutes = 1 faute corrigée $\approx$ 0,25 à 0,5 point gagné.
\end{itemize}
\end{exemplebox}

\begin{methode}[Check-list finale — Avant de rendre la copie]
\begin{enumerate}
    \item [$\square$] Sujet analysé sur brouillon (mots-clés, définitions, problématique, plan).
    \item [$\square$] Introduction~: 5 étapes présentes, plan annoncé = plan réalisé.
    \item [$\square$] Développement~: 3 parties visibles (I, II, III), chaque partie a 2--3 paragraphes.
    \item [$\square$] Chaque paragraphe~: 1 idée force + argument + exemple/citation + commentaire + mini-transition.
    \item [$\square$] Citations~: entre guillemets, intégrées, commentées, source indiquée.
    \item [$\square$] Transitions inter-parties~: bilan + limite + annonce (pas «~Passons à...~»).
    \item [$\square$] Conclusion~: bilan + réponse à la problématique + ouverture pertinente.
    \item [$\square$] Relecture orthographe/grammaire/ponctuation (15~min minimum).
    \item [$\square$] Paragraphes matérialisés (alinea), lisibilité, ratures propres.
    \item [$\square$] Nom, numéro, série, date en haut de la copie.
\end{enumerate}
\end{methode}

\newpage

\coursec{Activité d'intégration}

\coursec{Leçon 04 : Rédiger une dissertation}

\courssub{Objectifs de l'activité}
\noindent Cette activité d'intégration vise à mobiliser l'ensemble des savoir-faire travaillés au cours de la leçon « Rédiger une dissertation » pour produire une copie complète, structurée et argumentée, en situation d'examen (Probatoire). Les objectifs spécifiques sont les suivants :
\begin{itemize}
    \item \textbf{Analyser un sujet de dissertation} en identifiant les termes clés, le débat sous-jacent et en formulant une problématique pertinente.
    \item \textbf{Construire un plan dialectique ou thématique} cohérent, articulé en parties et sous-parties équilibrées.
    \item \textbf{Rédiger une introduction} respectant les quatre étapes canoniques (amorce, présentation du sujet, problématique, annonce du plan).
    \item \textbf{Développer des paragraphes argumentatifs} structurés (thèse, arguments, exemples/citations insérés, mini-synthèse) et reliés par des \textbf{transitions} logiques.
    \item \textbf{Rédiger une conclusion} synthétique (réponse à la problématique, reprise des grandes étapes, ouverture).
    \item \textbf{S'auto-évaluer et évaluer un pair} à l'aide d'une grille critériée pour identifier les points forts et les axes de remédiation.
\end{itemize}

\courssub{Situation}
\noindent \textbf{Contexte :} Vous êtes élève en classe de Première Technique (spécialité Secrétariat de Direction / Gestion / Électrotechnique) au Lycée Technique de Bafoussam. Dans le cadre de la préparation à l'épreuve écrite de Français du Probatoire, votre professeur vous propose un entraînement intégral sur deux séances (4 heures) à partir d'un dossier documentaire croisant littérature, sociologie locale et données statistiques.

\medskip
\noindent \textbf{Sujet de dissertation :}
\begin{center}
    \textbf{« La tradition est-elle un obstacle au progrès dans \emph{Le Lion et la Perle} de Wole Soyinka ? »}
\end{center}

\medskip
\noindent \textbf{Dossier documentaire fourni :}

\noindent \textbf{Document 1 : Extraits de \emph{Le Lion et la Perle} (Acte 1 et Acte 3).}
\begin{itemize}
    \item \textit{Extrait A (Lakunle, Acte 1) :} « Je veux vous épouser, Sidi. Mais pas selon la coutume barbare du prix de la dot... Une femme doit être l'égale de l'homme, pas sa propriété. » / « Nous achèterons une machine à coudre, une bicyclette, une radio... Nous serons modernes. »
    \item \textit{Extrait B (Baroka, Acte 1) :} « Le vieux doit couler pour que le nouveau puisse naître... Mais le nouveau ne doit pas oublier le vieux. » / Description de la ruse du Bale pour séduire Sidi (feinte d'impuissance, timbre-poste).
    \item \textit{Extrait C (Sidi, Acte 3) :} « Je suis la perle... Lui, c'est le lion. » / Refus d'épouser Lakunle sans dot : « Tu veux me prendre comme on vole un poulet au marché ? » / Choix final pour Baroka.
\end{itemize}

\noindent \textbf{Document 2 : Photographies commentées de Bafoussam (Région de l'Ouest, 2024).}
\begin{itemize}
    \item \textit{Photo 1 :} La chefferie traditionnelle (toits en chaume, cases des notables, place des fêtes) au pied d'un immeuble administratif moderne (béton, verre, 4 étages).
    \item \textit{Photo 2 :} Un marché de rue : vendeuses en pagnes traditionnels utilisant des terminaux de paiement mobile (Mobile Money) sur smartphones.
    \item \textit{Photo 3 :} Cérémonie du « Nguon » (fête culturelle annuelle) : le Fon (chef traditionnel) en tenue rituelle remet un prix d'excellence à de jeunes ingénieurs formés à l'ENSET de Douala.
\end{itemize}

\noindent \textbf{Document 3 : Données chiffrées Cameroun (Sources : INS, Rapport 2023 ; Banque Mondiale, 2022).}
\begin{itemize}
    \item Taux d'urbanisation : \textbf{58,2\%} (contre 12\% en 1960). Croissance urbaine annuelle : \textbf{3,6\%}.
    \item Taux de scolarisation primaire : \textbf{97\%} (garçons) / \textbf{94\%} (filles). Taux d'achèvement primaire : \textbf{72\%}.
    \item Part de l'économie informelle : \textbf{90\%} des emplois (forte prégnance des réseaux de solidarité traditionnels / tontines).
    \item Indice de développement humain (IDH) : \textbf{0,576} (rang 151/193). Progrès de +15\% depuis 2000.
\end{itemize}

\noindent \textbf{Consigne générale :} À partir de l'étude de la pièce \emph{Le Lion et la Perle} et des documents du dossier, vous traiterez le sujet en rédigeant une dissertation complète. Vous disposerez de 2 heures pour l'analyse, le plan et la rédaction de l'introduction et du développement (Séance 1), et de 2 heures pour la fin du développement, la conclusion et la relecture (Séance 2).

\courssub{Consignes}
\noindent Les questions ci-dessous guident votre travail. Répondez-y dans l'ordre sur votre brouillon avant de rédiger la copie définitive.

\begin{enumerate}
    \item \textbf{Analyse du sujet (30 min) :}
    \begin{enumerate}
        \item Définissez les termes « tradition », « obstacle », « progrès » dans le contexte de la pièce et de la société camerounaise actuelle.
        \item Identifiez les deux thèses opposées (thèse / antithèse) que le sujet invite à confronter.
        \item Formulez la \textbf{problématique} en une question précise qui guide toute la réflexion.
        \item Proposez un \textbf{plan détaillé} (2 ou 3 parties, 2 sous-parties chacune) avec les arguments principaux et les citations précises (Doc 1) ou exemples (Doc 2, 3) à mobiliser.
    \end{enumerate}
    \item \textbf{Rédaction de l'introduction (20 min) :} Rédigez l'introduction complète en respectant les 4 étapes. Soignez l'amorce (accroche sur le dossier ou citation).
    \item \textbf{Rédaction du développement (1h10) :} Rédigez les paragraphes de vos 2 ou 3 parties.
    \begin{itemize}
        \item Chaque paragraphe = 1 idée force + 1 argument + 1 illustration (citation insérée / exemple dossier) + 1 mini-synthèse.
        \item Insérez au minimum \textbf{3 citations exactes} de la pièce (guillemets, référence acte).
        \item Rédigez une \textbf{transition} explicite entre chaque grande partie (rappel thèse précédente + annonce thèse suivante + lien logique).
    \end{itemize}
    \item \textbf{Rédaction de la conclusion (20 min) :} Synthèse, réponse nuancée à la problématique, ouverture (sur le Cameroun d'aujourd'hui, sur l'universel, ou sur une autre œuvre).
    \item \textbf{Relecture et auto-évaluation (20 min) :} Relisez-vous (orthographe, syntaxe, cohérence, respect du sujet). Échangez votre copie avec un voisin. Évaluez-la avec la \textbf{Grille d'évaluation par les pairs} (fournie par le prof).
\end{enumerate}

\courssub{Ressources à mobiliser}
\noindent Rappel des méthodes et outils vus en classe (fiches méthodes) :

\begin{methode}[Analyser un sujet de dissertation]
\begin{enumerate}
    \item \textbf{Cercler les mots-clés} et les relier aux notions du cours (tradition = coutumes, transmission, identité ; progrès = modernité, technique, émancipation, éducation ; obstacle = frein, résistance, mais aussi garde-fou).
    \item \textbf{Repérer la structure} : « Est-ce que X est Y ? » $\rightarrow$ Sujet à débat, plan dialectique (Thèse / Antithèse / Synthèse) ou thématique (Causes / Conséquences / Solutions).
    \item \textbf{Problématiser} : Transformer le sujet en question tensionnée. Ex : « Dans quelle mesure la tradition, loin de freiner le progrès, peut-elle en être le fondement ou le régulateur ? »
    \item \textbf{Planifier} : 2 parties équilibrées (Thèse/Antithèse) + Synthèse (si plan dialectique) OU 3 parties thématiques distinctes. Annoncer le plan par des noms de parties nominaux (ex: « I. La tradition comme frein... II. La tradition comme levier... »).
\end{enumerate}
\end{methode}

\begin{methode}[Rédiger l'introduction (4 temps)]
\begin{enumerate}
    \item \textbf{Amorce :} Citation, fait de société (Doc 2/3), actualité, définition large. \textit{Éviter : « Depuis la nuit des temps... »}.
    \item \textbf{Présentation du sujet / Contexte :} Auteur, œuvre, genre, résumé très bref (2 lignes), ancrage dans le dossier.
    \item \textbf{Problématique :} La question centrale (une seule phrase, point d'interrogation).
    \item \textbf{Annonce du plan :} « Nous verrons d'abord que... Ensuite nous montrerons que... Enfin nous analyserons... »
\end{enumerate}
\end{methode}

\begin{methode}[Construire un paragraphe argumentatif (Structure A.E.I.)]
\begin{enumerate}
    \item \textbf{Argument (Phrase thèse) :} Idée directrice du paragraphe (en tête).
    \item \textbf{Explication / Développement :} Raisonnement, nuances, liens logiques (connecteurs : \emph{en effet, car, par conséquent}).
    \item \textbf{Illustration (Preuve) :} \textbf{Citation insérée} dans la phrase (pas isolée) + référence précise. \textit{Ex: Baroka affirme : « Le vieux doit couler... » (Acte 1).} Ou exemple du dossier (Photo, Chiffre).
    \item \textbf{Mini-synthèse / Transition interne :} Phrase de clôture qui valide l'argument et prépare le suivant.
\end{enumerate}
\end{methode}

\begin{methode}[Rédiger une transition entre parties]
\noindent Structure en 3 temps (phrase unique ou deux courtes) :
\begin{enumerate}
    \item \textbf{Bilan (Rappel) :} « Au terme de cette analyse, il apparaît que la tradition freine le progrès par son conservatisme... »
    \item \textbf{Relance (Limite / Ouverture) :} « Cependant, cette vision unilatérale occulte le rôle dynamique de la coutume... »
    \item \textbf{Annonce (Nouvelle partie) :} « C'est pourquoi nous examinerons maintenant comment la tradition devient le socle d'une modernité assumée. »
\end{enumerate}
\end{methode}

\begin{methode}[Rédiger la conclusion (3 temps)]
\begin{enumerate}
    \item \textbf{Synthèse :} Reprise des grandes idées des parties (sans copier-coller l'annonce du plan).
    \item \textbf{Réponse à la problématique :} Réponse nuancée, définitive. « En définitive, la tradition n'est pas un obstacle absolu mais un partenaire nécessaire... »
    \item \textbf{Ouverture :} Élargissement (autre œuvre : \emph{Les Bouts de bois de Dieu}, actualité camerounaise : décentralisation/chefferies, question philosophique : identité/mondialisation).
\end{enumerate}
\end{methode}

\begin{attention}[Pièges à éviter (Check-list relecture)]
\begin{itemize}
    \item \textbf{Hors-sujet :} Ne pas disséquer sur « La tradition au Cameroun » sans ancrer dans \emph{Le Lion et la Perle}.
    \item \textbf{Citations « orphelines » :} Citation jetée sans introduction ni commentaire. \textit{Règle : 1 citation = 1 phrase hôte + commentaire.}
    \item \textbf{Transitions manquantes ou « Mécaniques » :} « Passons maintenant à la deuxième partie. » $\rightarrow$ Interdit. Faire du lien logique.
    \item \textbf{Conclusion « nouvelle partie » :} Introduire un argument nouveau. Interdit.
    \item \textbf{Langage familier / SMS :} « genre », « grave », « du coup » (oral). Privilégier : \emph{en effet, par conséquent, ainsi, toutefois}.
\end{itemize}
\end{attention}

% --- SCHÉMA RÉCAPITULATIF DE LA STRUCTURE (TikZ) ---
\begin{popfigure}
\begin{tikzpicture}[
    node distance=1.2cm and 2cm,
    box/.style={draw=popblue, thick, rounded corners, fill=popblueL, text width=3.2cm, align=center, minimum height=1.8cm, font=\small},
    arrow/.style={-Stealth, thick, popdark},
    title/.style={font=\bfseries\small, text=popdark}
]
% Nodes
\node[box, title] (sujet) {Sujet \& Dossier};
\node[box, below=of sujet, align=center] (analyse){Analyse\\Mots-clés, Thèses,\\Problématique, Plan};
\node[box, below=of analyse, align=center] (intro){Introduction\\Amorce, Présentation,\\Problématique, Plan};
\node[box, below=of intro, align=center] (dev){Développement\\Partie I (Thèse)\\Partie II (Antithèse)\\Partie III (Synthèse)};
\node[box, below=of dev, align=center] (trans){Transitions\\Bilan + Relance + Annonce};
\node[box, below=of trans, align=center] (concl){Conclusion\\Synthèse, Réponse,\\Ouverture};
\node[box, below=of concl, align=center] (relect){Relecture \& Auto-évaluation\\Grille critériée};
% Arrows
\draw[arrow] (sujet) -- (analyse);
\draw[arrow] (analyse) -- (intro);
\draw[arrow] (intro) -- (dev);
\draw[arrow] (dev) -- (trans);
\draw[arrow] (trans) -- (concl);
\draw[arrow] (concl) -- (relect);
% Side annotations
\node[left=0.5cm of analyse, text=poppurple, font=\small\itshape] {Brouillon (30 min)};
\node[left=0.5cm of intro, text=poppurple, font=\small\itshape] {Rédaction (20 min)};
\node[left=0.5cm of dev, text=poppurple, font=\small\itshape] {Rédaction (1h10)};
\node[left=0.5cm of trans, text=poppurple, font=\small\itshape] {Intégré au dév.};
\node[left=0.5cm of concl, text=poppurple, font=\small\itshape] {Rédaction (20 min)};
\node[left=0.5cm of relect, text=poppurple, font=\small\itshape] {Finalisation (20 min)};
\end{tikzpicture}
\legende{Schéma de la démarche dissertation : de l'analyse du sujet à la relecture finale (4h au total).}
\end{popfigure}

\courssub{Démarche et corrigé commenté}
\begin{exoresolu}[Modélisation de la dissertation complète attendue]
\textbf{Énoncé :} À partir du dossier (pièce, photos Bafoussam, stats Cameroun), traiter le sujet : « La tradition est-elle un obstacle au progrès dans \emph{Le Lion et la Perle} de Wole Soyinka ? » en respectant la méthode (analyse, intro, développement 2-3 parties avec citations/transitions, conclusion).

\tcblower
\textbf{Solution détaillée et commentée}

\medskip
\noindent \textbf{1. ANALYSE DU SUJET (Brouillon structuré)}

\noindent \textbf{Mots-clés \& Définitions contextuelles :}
\begin{itemize}
    \item \textbf{Tradition} : Coutumes Yoruba (dot, polygamie, respect aînés, rites), incarnée par Baroka (Bale), Sadiku. Au Cameroun : chefferies, tontines, rites d'initiation, langues nationales, tenure coutumière.
    \item \textbf{Progrès} : Modernité technique (machine, timbre, école), émancipation individuelle (choix mari, éducation filles), développement infrastructurel. Incarné par Lakunle (instituteur). Au Cameroun : urbanisation (58\%), scolarisation (97\%), Mobile Money, ENSET.
    \item \textbf{Obstacle} : Frein, résistance (Lakunle : « coutume barbare »), mais aussi protection, régulation, fondement identitaire (Baroka : « le nouveau ne doit pas oublier le vieux »).
\end{itemize}

\noindent \textbf{Thèses en présence :}
\begin{itemize}
    \item \textbf{Thèse (Oui) :} La tradition freine l'innovation technique, l'égalité femmes-hommes, l'école laïque. (Lakunle vs Baroka/Sidi). Ex : refus dot = blocage mariage ; rejet école filles.
    \item \textbf{Antithèse (Non) :} La tradition structure la société, absorbe la modernité (Baroka timbre-poste), donne sens au progrès. Sidi choisit le « Lion » (racines) vs « Perle » (modernité vide). Au Cameroun : chefferies moteurs développement local (Nguon), tontines financent PME.
\end{itemize}

\noindent \textbf{Problématique choisie :}
\begin{center}
    \emph{« Dans quelle mesure la confrontation entre Baroka et Lakunle révèle-t-elle que la tradition, loin d'être un simple frein archaïque, est le creuset indispensable d'un progrès authentiquement africain ? »}
\end{center}

\noindent \textbf{Plan détaillé (Dialectique 3 parties) :}
\begin{itemize}
    \item \textbf{I. La tradition comme frein apparent au progrès technique et social (Thèse).}
    \begin{itemize}
        \item A. Le conservatisme des mœurs : obstacle à l'émancipation féminine et à l'école (Lakunle / Sidi / Dot). \textit{Cit. A : « coutume barbare du prix de la dot »}.
        \item B. La résistance au changement technique par méfiance ou intérêt de caste (Baroka feint de refuser le train/route). \textit{Doc 3 : Taux achèvement primaire 72\% (frein culturel ?)}.
    \end{itemize}
    \item \textbf{II. La tradition comme actrice et régulatrice du progrès (Antithèse).}
    \begin{itemize}
        \item A. L'intelligence adaptative : Baroka intègre la modernité (timbre-poste, machine) pour renforcer son pouvoir. \textit{Cit. B : « Le vieux doit couler pour que le nouveau puisse naître »}.
        \item B. Au Cameroun, les chefferies (Bafoussam) pilotent le développement (éducation, santé, foncier). \textit{Doc 2 Photo 3 : Fon remet prix ingénieurs ; Doc 3 : 90\% informel/tontines}.
    \end{itemize}
    \item \textbf{III. Vers une synthèse : Le progrès sans racine est aliénation, la tradition sans souffle est momie (Synthèse).}
    \begin{itemize}
        \item A. Le choix de Sidi : refus de la modernité mimétique (Lakunle « sans dot » = vol) pour la tradition vivante (Baroka). \textit{Cit. C : « Tu veux me prendre comme on vole un poulet » ; « Lui, c'est le lion »}.
        \item B. Modèle camerounais : « Modernité enracinée » (Décentralisation, bilinguisme, promotion cultures nationales). Ouverture sur \emph{Les Bouts de bois de Dieu} (grève cheminots : tradition solidarité = moteur progrès social).
    \end{itemize}
\end{itemize}

\medskip
\noindent \textbf{2. INTRODUCTION (Modèle rédigé)}

\noindent \textit{[Amorce - Doc 2/3]} Dans le Bafoussam de 2024, le toit de chaume de la chefferie traditionnelle voisine harmonieusement avec la façade de verre de la préfecture ; les tontines villageoises alimentent le capital de start-up urbaines via le Mobile Money, tandis que le Fon, gardien des ancêtres, décerne des prix d'excellence à de jeunes ingénieurs formés à l'ENSET. Cette coexistence visible, illustrée par un taux d'urbanisation de 58\% qui n'a pas effacé les solidarités coutumières (90\% d'emplois informels), donne chair à une interrogation universelle.

\noindent \textit{[Présentation]} C'est au cœur de cette tension que Wole Soyinka, prix Nobel de littérature 1986, situe sa comédie \emph{Le Lion et la Perle} (1959). Dans le village yoruba d'Ilujinle, l'instituteur Lakunle, héraut d'une modernité occidentale (école, monogamie, machines), rivalise avec le Bale Baroka, « Lion » rusé et polygame, pour conquérir Sidi, la « Perle », beauté villageoise consciente de sa valeur. Le conflit oppose deux visions du monde : l'importation pure et simple vs l'adaptation endogène.

\noindent \textit{[Problématique]} Dès lors, \textbf{dans quelle mesure la confrontation entre Baroka et Lakunle révèle-t-elle que la tradition, loin d'être un simple frein archaïque, est le creuset indispensable d'un progrès authentiquement africain ?}

\noindent \textit{[Annonce du plan]} Nous analyserons d'abord comment la tradition apparaît comme un obstacle au progrès technique et à l'émancipation individuelle (I). Nous démontrerons ensuite qu'elle en est paradoxalement le moteur et le régulateur, capable d'absorber la modernité (II). Enfin, nous montrerons que le véritable progrès réside dans une synthèse féconde, illustrée par le choix de Sidi et l'expérience camerounaise actuelle (III).

\medskip
\noindent \textbf{3. DÉVELOPPEMENT (Extraits significatifs avec commentaires méthodologiques)}

\noindent \textbf{PARTIE I : La tradition comme frein apparent au progrès technique et social}

\noindent \textbf{Argument A : Le conservatisme des mœurs entrave l'émancipation féminine et l'instruction.}
Lakunle incarne le regard moderne qui perçoit la coutume comme une « barrière ». Pour lui, la dot est une « coutume barbare » (Acte 1) qui transforme la femme en « propriété ». Il refuse de la payer, voulant faire de Sidi une « épouse moderne », « égale de l'homme ». \textbf{Or}, ce refus heurte la logique sociale d'Ilujinle : sans dot, le mariage n'a pas de valeur juridique ni symbolique aux yeux de la communauté. Sidi elle-même rejette cette modernité sans ancrage : « Tu veux me prendre comme on vole un poulet au marché ? » (Acte 3). La tradition, ici, n'est pas seulement un frein ; elle est le \textbf{garant de la dignité sociale} de la femme. Parallèlement, l'opposition à l'école des filles (Baroka détourne le projet d'école pour en faire un outil de séduction) montre une instrumentalisation du pouvoir traditionnel qui peut freiner l'instruction (Doc 3 : taux d'achèvement 72\%, écart filles/garçons).

\noindent \textbf{Argument B : La résistance au changement technique par peur de la perte de pouvoir.}
Baroka se vante d'avoir fait échouer le projet de voie ferrée traversant Ilujinle : « J'ai payé les arpenteurs pour qu'ils détournent le tracé » (Acte 1, évoqué). Il craint que le train n'apporte une modernité qui le dépossède de son autorité (contrôle des flux, des marchés). C'est une vision \textbf{conservatrice du pouvoir} : la tradition protège le chef, pas le peuple. Cela fait écho aux situations où certaines chefferies camerounaises freinent l'aménagement du territoire ou l'expropriation pour utilité publique au nom du foncier coutumier.

\noindent \textbf{[Mini-synthèse Partie I]} En définitive, si la tradition freine le progrès, c'est souvent lorsqu'elle est brandie comme un privilège de caste ou un refus de l'altérité, figeant l'identité dans un passé mythifié.

\medskip
\noindent \textbf{TRANSITION I $\rightarrow$ II}
\begin{center}
    \textit{« Au terme de cette analyse, il apparaît que la tradition peut freiner le progrès lorsqu'elle se fige en conservatisme défenseur d'intérêts particuliers. \textbf{Cependant}, réduire la tradition à ce rôle de blocage serait ignorer sa plasticité historique : à Ilujinle comme à Bafoussam, les détenteurs de l'autorité coutumière sont souvent les premiers artisans de la modernité. \textbf{C'est pourquoi} nous devons examiner comment la tradition devient le levier et le régulateur du progrès. »}
\end{center}

\medskip
\noindent \textbf{PARTIE II : La tradition comme actrice et régulatrice du progrès}

\noindent \textbf{Argument A : L'intelligence adaptative de Baroka : « domestiquer » la modernité.}
Contrairement à Lakunle qui \emph{subït} la modernité (il la mime, il la cite : « machine à coudre, bicyclette, radio » sans les maîtriser vraiment), Baroka la \emph{maîtrise}. Il utilise le timbre-poste (symbole de l'administration coloniale/moderne) pour piéger Sidi (photo dans le magazine). Il déclare : « \textbf{Le vieux doit couler pour que le nouveau puisse naître... Mais le nouveau ne doit pas oublier le vieux} » (Acte 1). Cette maxime est le \textbf{manifeste d'une modernité enracinée}. Baroka veut l'imprimerie, l'électricité, mais à ses conditions. Il transforme l'outil moderne en instrument de pouvoir traditionnel.

\noindent \textbf{Argument B : Le modèle camerounais : Chefferies et développement local (Doc 2 \& 3).}
Les photographies de Bafoussam ne mentent pas. La Photo 3 (Cérémonie du Nguon) est emblématique : le Fon, autorité traditionnelle suprême, préside la remise de prix à des ingénieurs. La tradition n'est pas hors du temps ; elle \textbf{consacre} la réussite moderne. Les données de l'INS confirment cette symbiose : l'économie informelle (90\% emplois), régie par des règles coutumières (tontines, solidarité lignagère), est le véritable poumon de l'économie nationale, bien plus que le secteur formel moderne. La tradition fournit le \textbf{capital social} (confiance, réseaux) sans lequel le progrès technique (Mobile Money, Doc 2 Photo 2) ne serait qu'une coquille vide.

\noindent \textbf{[Mini-synthèse Partie II]} La tradition n'est donc pas l'ennemie du progrès ; elle en est la \textbf{matrice culturelle}. Elle filtre la modernité, en rejette l'aliénant (mimétisme de Lakunle) et en intègre l'utile (timbre, école, hôpital).

\medskip
\noindent \textbf{TRANSITION II $\rightarrow$ III}
\begin{center}
    \textit{« Ainsi, la tradition se révèle être le creuset où le progrès se forge une identité. \textbf{Néanmoins}, cette alliance n'est pas automatique : elle suppose un choix conscient, un arbitrage lucide entre ce qui doit être gardé et ce qui doit être transformé. \textbf{C'est ce moment de vérité} que révèle le dénouement de la pièce à travers le choix définitif de Sidi, et que vit aujourd'hui la jeunesse camerounaise. »}
\end{center}

\medskip
\noindent \textbf{PARTIE III : Synthèse — Le progrès sans racine est aliénation, la tradition sans souffle est momie}

\noindent \textbf{Argument A : Le choix de Sidi : refus de l'aliénation mimétique, acceptation de l'hybridité.}
Sidi a le choix entre deux « modernités » : celle de Lakunle, qui nie ses racines (« femme égale de l'homme » mais sans dot = sans statut), et celle de Baroka, qui assume sa puissance virile et son autorité coutumière tout en possédant les codes du moderne (timbre, voiture, future imprimerie). En choisissant Baroka — « \textbf{Lui, c'est le lion} » (Acte 3) —, Sidi opte pour une \textbf{identité assumée}. Elle refuse d'être une « Perle » décorative, objet de collection pour moderne superficiel. Elle devient l'épouse du « Lion », c'est-à-dire l'actrice d'une tradition vivante, capable d'intégrer le nouveau. C'est la victoire de l'\emph{authenticité} sur le \emph{mimétisme}.

\noindent \textbf{Argument B : Ouverture — Le Cameroun face au défi de la « Modernité enracinée ».}
Le dilemme d'Ilujinle est celui du Cameroun d'aujourd'hui. La décentralisation (lois 2004/2009) institutionnalise le rôle des chefferies dans le développement local. L'enseignement des langues nationales (programme ELAN) et la promotion du patrimoine (Ngondo, Nguon, Nyem-Nyem) visent à faire de la tradition le \textbf{socle cognitif et identitaire} de l'élève-ingénieur de demain. Comme l'écrivait Sembène Ousmane dans \emph{Les Bouts de bois de Dieu}, la grève des cheminots (progrès social) a vaincu parce qu'elle s'est appuyée sur la tradition de solidarité africaine (la « caisse de grève » gérée comme une tontine). Le progrès véritable n'est pas la table rase, mais la \textbf{greffe}.

\medskip
\noindent \textbf{4. CONCLUSION (Modèle rédigé)}

\noindent \textit{[Synthèse]} Au terme de cette étude, il apparaît que la relation entre tradition et progrès dans \emph{Le Lion et la Perle} n'est pas une opposition binaire mais une dialectique féconde. Nous avons vu que la tradition peut freiner le progrès lorsqu'elle se fige en conservatisme (I), mais qu'elle en devient le principal moteur lorsqu'elle fait preuve d'intelligence adaptative, intégrant les techniques sans perdre son âme (II). Le choix final de Sidi pour Baroka, le « Lion », scelle cette alliance : le progrès n'a de sens que s'il est porté par une identité vivante (III).

\noindent \textit{[Réponse à la problématique]} \textbf{Non, la tradition n'est pas un obstacle au progrès dans l'œuvre de Soyinka.} Elle en est la condition de possibilité. L'obstacle réel n'est pas la tradition, mais le \emph{mimétisme} stérile (Lakunle) qui veut greffer des branches modernes sur un tronc coupé. Le progrès authentique, celui que Baroka incarne et que le Cameroun construit à Bafoussam comme à Maroua, est une \textbf{modernité enracinée}.

\noindent \textit{[Ouverture]} Cette leçon dépasse le cadre yoruba. Elle interroge l'avenir de l'Afrique dans la mondialisation : comment « moderniser la tradition et traditionaliser la modernité » (Cheikh Hamidou Kane, \emph{L'Aventure ambiguë}) ? Le défi du Probatoire, comme celui de l'ingénieur technicien, est de devenir ce « Lion » qui manie la perle technique sans perdre sa crinière identitaire.

\medskip
\noindent \textbf{5. GRILLE D'AUTO-ÉVALUATION / ÉVALUATION PAIRE (Extrait)}

\begin{center}
\begin{tabularx}{\linewidth}{|l|X|X|X|}\hline
\rowcolor{popblueL}
\textbf{Critères (Barème indicatif /20)} & \textbf{Non acquis (0-5)} & \textbf{En acquisition (6-11)} & \textbf{Acquis / Expert (12-20)} \\ \hline
\textbf{Analyse \& Problématique} & Sujet recopié, pas de définitions, pb mal formulée. & Termes définis, pb présente mais floue ou hors tension. & Termes précis, pb tensionnée, guide tout le devoir. \\ \hline
\textbf{Introduction (4 temps)} & Éléments manquants, amorce banale, plan absent. & 4 temps présents mais déséquilibrés, plan mécanique. & Amorce percutante (dossier), pb claire, plan annoncé fluide. \\ \hline
\textbf{Développement (Structure A.E.I.)} & Paragraphes sans thèse, citations absentes ou orphelines. & Thèses claires, citations insérées mais peu commentées. & Thèses fortes, citations intégrées \& commentées, ex. Doc 2/3. \\ \hline
\textbf{Transitions} & Absentes ou « Passons à la suite ». & Présentes mais mécaniques (résumé/annonce seul). & Logiques : Bilan + Relance + Annonce (lien causal). \\ \hline
\textbf{Conclusion (3 temps)} & Nouvelle idée, pas de réponse, pas d'ouverture. & Synthèse + Réponse, ouverture banale. & Synthèse nuancée, réponse définitive, ouverture pertinente. \\ \hline
\textbf{Langue \& Présentation} & Nombreuses fautes (syntaxe, orthographe), style oral. & Correct, quelques lourdeurs, style scolaire. & Français standard impeccable, style fluide, connecteurs variés. \\ \hline
\end{tabularx}
\end{center}

\end{exoresolu}

\courssub{Bilan de l'activité}
\noindent Cette activité d'intégration a permis de vérifier l'acquisition effective des compétences visées par la leçon « Rédiger une dissertation » dans un contexte authentique et motivant (Probatoire, ancrage camerounais).

\begin{aretenir}[Bilan pédagogique]
\begin{itemize}
    \item \textbf{Mobilisation des savoirs :} Les élèves ont dû activer simultanément la réception (connaissance fine de \emph{Le Lion et la Perle}, citations exactes), la langue (types de phrases pour varier le rythme, énonciation pour marquer la subjectivité/objectivité, connecteurs logiques), et la production (méthodes intro/développement/conclusion).
    \item \textbf{Transfert contextuel :} L'utilisation du dossier « Bafoussam / Statistiques Cameroun » a forcé le transfert des concepts littéraires (tradition/progrès) vers la réalité sociologique immédiate des élèves, donnant du sens à l'exercice scolaire. La citation de Cheikh Hamidou Kane en ouverture ancre la réflexion dans une francophonie africaine partagée.
    \item \textbf{Métacognition par l'évaluation par les pairs :} La grille critériée a objectivé les attendus. Les élèves, en corrigeant le voisin, sont devenus conscients des critères de réussite (ex: « citation orpheline » vs « citation insérée », « transition mécanique » vs « transition logique »). Cela prépare l'autonomie pour l'examen.
    \item \textbf{Remédiation ciblée :} Le compte-rendu collectif a révélé trois faiblesses récurrentes à retraiter en ateliers :
    \begin{enumerate}
        \item \textbf{Problématique « plate » :} Beaucoup d'élèves posent une question descriptive (« Quels sont les obstacles ? ») au lieu d'une question tensionnée (« Dans quelle mesure... »).
        \item \textbf{Insertion des citations :} Tendance à « coller » la citation entre deux paragraphes sans phrase hôte. Rappel de la règle : \emph{Intro citation + Citation + Commentaire}.
        \item \textbf{Gestion du temps :} Sur 4h, beaucoup ont passé 1h30 sur l'introduction. Nécessité d'entraînements chronométrés (ex: 20 min intro, 1h30 développement, 20 min conclusion, 20 min relecture).
    \end{enumerate}
\end{itemize}
\end{aretenir}

\noindent Cette activité constitue un jalon majeur du portfolio de l'élève pour le Probatoire. Elle valide la capacité à \textbf{produire un texte long, structuré, argumenté et ancré}, compétence transversale essentielle pour la poursuite d'études techniques (rapports de stage, notes de synthèse, mémoires).

\newpage

\coursec{Méthodes et exercices résolus}

\coursec{Rédiger une dissertation}

\courssub{Étape 1 : Lecture méthodique n° 1 du \emph{Lion et la perle} et types de phrase}

\begin{experience}[Lecture méthodique n° 1 : Découverte de l'œuvre et contexte]
\textbf{Objectif} : Situer l'œuvre, identifier le genre, le cadre spatio-temporel et les personnages principaux.\par
\textbf{Support} : \emph{Le Lion et la perle}, Wole Soyinka (traduction française, Présence Africaine).\par
\textbf{Protocole} :
\begin{enumerate}
\item \textbf{Paratexte} : Auteur (Nigérian, prix Nobel 1986), titre (métaphore : Baroka = Lion, Sidi = Perle), date (1959, pré-indépendance), genre (comédie satirique, théâtre).
\item \textbf{Cadre} : Ilujinle (village yoruba fictif), années 50. Conflit Tradition vs Modernité.
\item \textbf{Personnages} : Baroka (Bale, 62 ans, rusé, traditionaliste), Lakunle (instituteur, 23 ans, moderniste caricatural), Sidi (belle villageoise, enjeu du mariage), Sadiku (épouse aînée, manipulée puis manipulatrice).
\item \textbf{Intrigue principale} : Rivalité pour la main de Sidi. Baroka use de ruse (fausse impuissance) ; Lakunle refuse la dot. Triomphe de la tradition adaptée (Baroka).
\end{enumerate}
\textbf{Observations attendues} : Théâtre dialogué, didascalies importantes (mimes, danses), langue imagée (proverbes, chants), ironie dramatique.
\legende{Lecture guidée : Acte 1 (Matin) -- Le mime du voyageur perdu / La proposition de Lakunle.}
\end{experience}

\begin{propriete}[Les quatre types de phrase (Morphosyntaxe)]
En dissertation, le choix du type de phrase structure la pensée.
\begin{center}
\begin{tabularx}{\linewidth}{|l|X|X|}\hline
\textbf{Type} & \textbf{Forme / Marqueurs} & \textbf{Fonction en dissertation} \\ \hline
\cle{Déclarative} (Assertive) & Sujet + Verbe (indicatif). Point final. & \textbf{Majoritaire}. Affirme les arguments (A), explications (E), interprétations (I). Ex : \emph{La tradition freine le progrès.} \\ \hline
\cle{Interrogative} & Inversion, \emph{est-ce que}, mot interrogatif. Point d'interrogation. & \textbf{Problématique} (intro), \textbf{questions rhétoriques} (transition, relance). Ex : \emph{La tradition est-elle un obstacle ?} \\ \hline
\cle{Impérative} & Verbe à l'impératif (2e pers.), pas de sujet. Point/Exclamation. & \textbf{Annonce du plan} (\emph{Voyons, Observons}), \textbf{chute} (appel). Rare dans le corps. \\ \hline
\cle{Exclamative} & \emph{Que/Quelle/Comme} + indicatif/subjonctif. Point d'exclamation. & \textbf{Amorce} (indignation, admiration), \textbf{chute} (force affective). Ex : \emph{Quel gâchis que cette modernité sans âme !} \\ \hline
\end{tabularx}
\end{center}
\textbf{Attention} : Au Probatoire, l'abus de phrases interrogatives/exclamatives dans le développement affaiblit l'argumentation (ton polémique vs analytique). Privilégier la déclarative pour l'autorité du raisonnement.
\end{propriete}

\courssub{Étape 2 : Analyser un sujet et rédiger l'introduction}

\begin{methode}[Analyser un sujet de dissertation : la méthode en 4 temps]
\textbf{Étape 1 : Délimiter le sujet} — Souligner les \cle{mots-clés} (notions centrales) et les \cle{mots-outils} (verbes opérateurs : \emph{discuter, analyser, dans quelle mesure, comment, pourquoi}). Définir chaque terme clé (dictionnaire, sens commun, sens littéraire/contextuel).\par
\textbf{Étape 2 : Problématiser} — Transformer le sujet en \cle{problématique} (question ouverte, débattue, qui appelle une réponse nuancée). Éviter les questions fermées (oui/non) ou les reformulations plates.\par
\textbf{Étape 3 : Trouver les arguments} — Faire un \emph{remue-méninges} (brainstorming) : idées pour, idées contre, nuances, exemples littéraires (\emph{Le Lion et la perle}, autres œuvres au programme ou culture personnelle), faits de société, actualité camerounaise. Trier, regrouper, hiérarchiser (argument fort / faible / illustration).\par
\textbf{Étape 4 : Construire le plan} — Choisir une structure adaptée au sujet :
\begin{itemize}
\item \textbf{Dialectique} (Thèse / Antithèse / Synthèse) : sujet opposant deux concepts (\emph{obstacle / progrès}).
\item \textbf{Thématique} (2-3 aspects distincts) : sujet portant sur une notion complexe (\emph{le théâtre, lieu de débat}).
\item \textbf{Analytique} (Causes / Conséquences / Solutions) : sujet processuel.
\end{itemize}
Rédiger les \cle{titres de parties} (phrase nominale affirmative, majuscule initiale, pas de point final) et les \cle{sous-parties} (1 argument = 1 paragraphe).
\end{methode}

\begin{exoresolu}{Analyser le sujet : « La tradition est-elle un obstacle au progrès ? »}
\textbf{Étape 1 — Délimitation} :\par
\cle{Tradition} : ensemble de coutumes, croyances, valeurs, rites transmis de génération en génération (ex. : coutumes ilujinle, dot, chefferie, masque dans \emph{Le Lion et la perle}).\par
\cle{Progrès} : évolution vers un mieux (technique, social, moral), modernité (ex. : école, poste, route, chemin de fer, imprimerie dans la pièce).\par
\cle{Obstacle} : ce qui freine, empêche, ralentit, bloque.\par
\cle{Verbe opérateur} : « Est-ce que » / « Dans quelle mesure » $\rightarrow$ réponse nuancée attendue (pas de binaire).\par
\textbf{Étape 2 — Problématique} : \emph{La tradition freine-t-elle systématiquement le progrès ou peut-elle, au contraire, en constituer le fondement vivant et le gardien critique ?}\par
\textbf{Étape 3 — Arguments} :\par
\textit{Thèse (Oui, obstacle)} : Baroka refuse la route/le chemin de fer pour garder son pouvoir (scène du géomètre) ; Lakunle veut imposer la modernité à marche forcée (refus de la dot, mimétisme) ; Sidi hésite, perd ses repères.\par
\textit{Antithèse (Non, socle/levier)} : La tradition assure la cohésion sociale (fêtes, rites, mime du voyageur) ; Sadiku use de ruse traditionnelle pour piéger Baroka ; le progrès sans racine mène à l'aliénation (Lakunle caricature « à l'anglaise »).\par
\textit{Synthèse} : Dialogue nécessaire. Le progrès doit s'enraciner ; la tradition doit s'adapter (ex. : Baroka imprime les timbres à l'effigie de Sidi).\par
\textbf{Étape 4 — Plan dialectique} :\par
\textbf{I. La tradition comme frein au progrès} (Poids du passé, conservatisme des puissants, rejet de l'innovation technique).\par
\textbf{II. La tradition comme socle et levier du progrès} (Identité collective, sagesse ancestrale, capacité d'adaptation).\par
\textbf{III. Vers une réconciliation : le progrès enraciné} (Modernité critique, tradition vivante, synthèse chez Soyinka).
\tcblower
\textbf{Analyse de la démarche} : L'analyse respecte les 4 temps. La problématique évite le piège du « Oui/Non » en intégrant « systématiquement » et « ou peut-elle en être le fondement ». Le plan dialectique est le plus pertinent car le sujet oppose deux concepts. Les arguments s'appuient précisément sur \emph{Le Lion et la perle} (œuvre au programme) et non sur des généralités. Les titres de parties sont des phrases nominales affirmatives.
\end{exoresolu}

\begin{methode}[Rédiger l'introduction : l'« entonnoir » en 4 temps]
\textbf{1. L'amorce (accroche)} — 1 phrase large : citation d'auteur, fait d'actualité, définition générale, référence à l'œuvre au programme. \textbf{Interdit} : « Depuis la nuit des temps... », « Je vais vous parler de... », « Dans cette dissertation... ».\par
\textbf{2. La présentation du sujet} — Reformuler le sujet en le situant (contexte, enjeux). Nommer l'œuvre de référence (\emph{Le Lion et la perle} de Wole Soyinka) et les personnages clés si pertinent.\par
\textbf{3. La problématique} — La question centrale (issue de l'analyse). Elle doit être \cle{unique}, \cle{ouverte}, \cle{débattue}.\par
\textbf{4. L'annonce du plan} — « Nous verrons d'abord que... Puis nous montrerons que... Enfin nous démontrerons que... » (utiliser \cle{connecteurs logiques} : \emph{premier lieu, ensuite, finalement}). \textbf{Pas de numérotation I, II, III dans le texte courant}.
\end{methode}

\begin{exoresolu}{Rédiger l'introduction pour : « La tradition est-elle un obstacle au progrès ? »}
\textbf{Proposition d'introduction} :\par
\emph{« La tradition n'est pas le culte des cendres, mais la conservation du feu », écrivait Gustav Mahler. Cette maxime éclaire le débat qui oppose, dans l'Afrique des années 50, les tenants de la modernité aux gardiens des coutumes. Dans \emph{Le Lion et la perle}, Wole Soyinka met en scène, à travers le trio Baroka, Lakunle et Sidi, le choc entre le village d'Ilujinle et les influences occidentales. \textbf{La tradition freine-t-elle systématiquement le progrès ou peut-elle, au contraire, en constituer le fondement vivant et le gardien critique ?} Nous analyserons d'abord comment la tradition peut apparaître comme un frein au changement, avant d'étudier son rôle de socle identitaire, pour enfin montrer la nécessité d'un progrès enraciné. »}\par
\tcblower
\textbf{Analyse de la rédaction} :\par
1. \textbf{Amorce} : Citation de Mahler (pertinente, courte), lien immédiat avec le sujet (feu/cendres = vie/mort de la tradition).\par
2. \textbf{Présentation} : Contexte (Afrique années 50), œuvre citée (titre en italique, auteur), personnages clés nommés, enjeu (choc).\par
3. \textbf{Problématique} : Reprise exacte de celle trouvée à l'étape 2, en gras ici pour repérage.\par
4. \textbf{Annonce} : Trois segments clairs, connecteurs (\emph{d'abord, avant d'étudier, pour enfin montrer}), pas de « I, II, III », vocabulaire varié (\emph{frein, socle, nécessité, enraciné}).
\end{exoresolu}

\courssub{Étape 3 : Lecture méthodique n° 2, Dénotation/Connotation et paragraphe argumentatif}

\begin{experience}[Lecture méthodique n° 2 : Le conflit des langages (Acte 1 \& 2)]
\textbf{Focus} : Opposition lexicale et stylistique entre Baroka et Lakunle.\par
\textbf{Extrait clé} : Scène de la poste (Baroka vs Géomètre) / Proposition de Lakunle à Sidi.\par
\textbf{Activité} : Relever les champs lexicaux opposés (Nature/Tradition vs Technique/Occident). Noter l'ironie de Soyinka.
\end{experience}

\begin{definition}[Dénotation et Connotation (Sémantique/Lexicologie)]
\begin{itemize}
\item \cle{Dénotation} : Sens premier, littéral, objectif, partagé par tous (dictionnaire). Ex : \emph{Lion} = grand félin.
\item \cle{Connotation} : Sens secondaire, subjectif, affectif, culturel, suggéré par le contexte. Ex : \cle{Lion} $\rightarrow$ force, royauté, danger, Baroka (dans la pièce).
\end{itemize}
\textbf{En dissertation} : Maîtriser la connotation permet de \textbf{nuancer} (éviter le manichéisme), d'\textbf{analyser} les citations (pourquoi ce mot ?), d'\textbf{argumenter} par le style (le choix lexical \emph{est} un argument).
\end{definition}

\begin{methode}[Construire un paragraphe argumentatif : la structure « A.E.C.I.T. »]
\textbf{A — Argument (Phrase-thèse)} : Une phrase affirmative qui résume l'idée directrice du paragraphe (ex. : « La tradition devient un obstacle au progrès lorsqu'elle est instrumentalisée par le pouvoir pour conserver ses privilèges. »).\par
\textbf{E — Explication} : Développer le raisonnement (cause, conséquence, mécanisme psychologique/social). 2-3 phrases. Contextualiser.\par
\textbf{C — Citation / Exemple précis} — \textbf{Obligatoire} au Probatoire. Insérer une \cle{citation courte} (guillemets français « ... ») ou une référence précise (scène, acte, didascalie). \emph{Ne jamais laisser une citation « nue »} (sans introduction ni commentaire).\par
\textbf{I — Interprétation (Commentaire)} — Expliquer \emph{pourquoi} cet exemple prouve l'argument. Analyser les mots (champ lexical, connotation), le contexte, l'intention de l'auteur, l'effet sur le lecteur/spectateur. C'est le cœur de la note (valeur 50\% du paragraphe).\par
\textbf{T — Transition (Mini-conclusion + Ouverture)} — Résumer l'apport du paragraphe ET annoncer le suivant (mots de liaison + reprise lexicale). Ex : « \emph{Toutefois}, si le pouvoir use de la tradition pour se figer, le peuple en fait un levier... »
\end{methode}

\begin{exoresolu}{Rédiger le paragraphe : « Le conservatisme de Baroka freine le développement d'Ilujinle » (Partie I, §1)}
\textbf{Paragraphe rédigé} :\par
\textbf{[A]} La tradition devient un obstacle au progrès lorsqu'elle est instrumentalisée par le pouvoir pour conserver ses privilèges. \textbf{[E]} Baroka, le Bale, incarne ce conservatisme intéressé : il perçoit la modernité (route, chemin de fer, école) comme une menace directe pour son autorité absolue et ses prérogatives coutumières. \textbf{[C]} Ainsi, face au géomètre venu tracer le passage de la voie ferrée, il affirme avec morgue : « \emph{Je ne veux pas de cette route qui amènera le bruit et la fumée dans mon village} » (Acte 2, scène de la poste). \textbf{[I]} Le lexique négatif (\emph{bruit, fumée}) et le possessif « \emph{mon village} » révèlent un égoïsme patrimonial : Baroka confond l'intérêt collectif et sa propre toute-puissance. En refusant la voie ferrée, il prive Ilujinle d'un développement économique et d'une ouverture sur le monde. \textbf{[T]} \emph{Toutefois, cette vision figée de la tradition masque une autre réalité : celle d'un peuple qui, à travers ses rites, préserve une identité sans laquelle le progrès serait vide.}
\tcblower
\textbf{Analyse A.E.C.I.T.} :\par
- \textbf{A} : Phrase-thèse claire, nuancée (\emph{lorsqu'elle est instrumentalisée}).\par
- \textbf{E} : Contextualisation psychologique et sociale (pouvoir, autorité, menace).\par
- \textbf{C} : Citation exacte (traduction standard), intégrée dans la phrase, source localisée (Acte 2, scène poste).\par
- \textbf{I} : Analyse lexicale (\emph{bruit, fumée, mon village}), interprétation psychologique (\emph{égoïsme patrimonial}), conséquence concrète (\emph{prive Ilujinle}).\par
- \textbf{T} : Pivot dialectique (\emph{Toutefois}), lien logique vers la partie II (le peuple / l'identité).
\end{exoresolu}

\begin{methode}[Rédiger les transitions : les 3 registres de liens logiques]
La transition assure la \cle{cohérence} du développement. Elle se place \textbf{à la fin du paragraphe} (mini-conclusion + ouverture) OU \textbf{au début du suivant} (rappel + annonce). Trois registres :\par
\textbf{1. L'addition / L'approfondissement} (même thèse) : \emph{Par ailleurs, De surcroît, En outre, D'ailleurs, Qui plus est} $\rightarrow$ nouvel argument parallèle.\par
\textbf{2. L'opposition / La concession} (Thèse $\rightarrow$ Antithèse / Nuance) : \emph{Cependant, Néanmoins, Certes... mais, En revanche, Au contraire, Si... toutefois} $\rightarrow$ pivot argumentatif.\par
\textbf{3. La conséquence / La synthèse} (Antithèse $\rightarrow$ Synthèse / Fin de partie) : \emph{Ainsi, Par conséquent, C'est pourquoi, En définitive, Au total} $\rightarrow$ conclusion logique.\par
\textbf{Règle d'or} : La transition doit \textbf{reprendre un mot-clé} de l'idée précédente et \textbf{en introduire un nouveau}.
\end{methode}

\begin{exoresolu}{Corriger et améliorer des transitions maladroites}
\textbf{Exercice} : Améliorer ces transitions entre deux paragraphes.\par
\textit{Cas 1 :} P1 finit sur « ... Baroka refuse l'école pour ses fils. » / P2 commence par « \textbf{Aussi}, la tradition empêche l'émancipation des femmes. »\par
\textit{Cas 2 :} P1 finit sur « ... Sadiku subit la polygamie. » / P2 commence par « \textbf{Donc}, Lakunle représente la modernité. »\par
\tcblower
\textbf{Corrections argumentées} :\par
\textbf{Cas 1} : « \emph{Aussi} » marque une conséquence logique. Or le deuxième argument (femmes) n'est pas la conséquence du premier (école/fils), c'est un \textbf{nouvel argument parallèle} (addition).\par
$\rightarrow$ \textbf{Meilleur} : « \emph{Par ailleurs}, la tradition entrave aussi l'émancipation des femmes... » ou en fin de P1 : « \emph{De surcroît}, ce conservatisme pèse lourdement sur le sort des femmes... »\par
\textbf{Cas 2} : « \emph{Donc} » marque une déduction. Or, Lakunle n'est pas la conséquence de la polygamie ; c'est un \textbf{personnage opposé} (contraste / antithèse).\par
$\rightarrow$ \textbf{Meilleur} : « \emph{À l'opposé de cette soumission}, Lakunle incarne une modernité conquérante... » ou « \emph{Néanmoins}, face à cette tradition oppressive, se dresse la figure de Lakunle... »
\end{exoresolu}

\courssub{Étape 4 : Lecture méthodique n° 3, Énonciation et conclusion}

\begin{experience}[Lecture méthodique n° 3 : La résolution et l'ironie finale (Acte 3)]
\textbf{Focus} : Le triomphe de Baroka (ruse), le choix de Sidi, l'adaptation (timbres).\par
\textbf{Énonciation} : Repérer les indices de subjectivité (modalisateurs, temps, personnes, déictiques) dans les discours de victoire de Baroka vs le désarroi de Lakunle. Qui parle ? À qui ? Pourquoi ? Comment ?
\end{experience}

\begin{definition}[L'énonciation (Communication)]
L'énonciation est l'acte par lequel un \cle{locuteur} produit un énoncé adressé à un \cle{allocutaire}, dans un \cle{contexte} (spatio-temporel).
\begin{itemize}
\item \textbf{Indices du locuteur} : \emph{Je, me, moi, nous} (personnes), \emph{ici, maintenant, demain} (déictiques spatio-temporels), \emph{ceci, cela} (démonstratifs).
\item \textbf{Modalisation} (subjectivité) : Verbes d'opinion (\emph{je pense, je crois}), adverbes (\emph{certes, sans doute, hélas}), modaux (\emph{devoir, pouvoir, falloir}), conditionnel, subjonctif, phrases exclamatives/interrogatives.
\item \textbf{En dissertation} : Le « Je » est \textbf{proscrit} (sauf citation). On utilise le \cle{« Nous » de modestie} (\emph{Nous verrons, Nous analyserons}) ou l'\cle{impersonnel} (\emph{Il apparaît que, Force est de constater}). L'analyse de l'énonciation dans les citations (ex. : le « Je » arrogant de Baroka vs le « Je » mimétique de Lakunle) nourrit l'étape I (Interprétation).
\end{itemize}
\end{definition}

\begin{methode}[Rédiger la conclusion : la symétrie inversée en 3 temps]
\textbf{1. Le bilan (Récapitulatif synthétique)} — 2-3 phrases. Reprendre la problématique et y répondre \textbf{en résumant la trajectoire de la pensée} (Thèse $\rightarrow$ Antithèse $\rightarrow$ Synthèse). \textbf{Interdit} : liste « I. J'ai dit que... II. J'ai montré que... ». Utiliser : « Au terme de cette étude, il apparaît que... », « En définitive... »\par
\textbf{2. L'ouverture (Élargissement)} — 1-2 phrases. Dépasser le sujet sans hors-sujet. Pistes : autre œuvre (autre pièce de Soyinka : \emph{La Danse des forêts}, autre auteur africain : Achebe, Césaire), autre époque, autre genre (roman, poème, essai), actualité camerounaise (chefferie, décentralisation), question philosophique durable. \textbf{Interdit} : « Pour conclure, je dirais que... », question rhétorique bête (« Et vous, qu'en pensez-vous ? »).\par
\textbf{3. La chute (Dernière phrase forte)} — Une phrase courte, mémorable, souvent imagée ou aphoristique, qui claque la porte du devoir.
\end{methode}

\begin{exoresolu}{Rédiger la conclusion pour le sujet sur tradition/progrès}
\textbf{Proposition de conclusion} :\par
\textbf{[Bilan]} Au terme de cette analyse, il apparaît que la tradition n'est pas intrinsèquement ennemie du progrès : elle en devient l'obstacle seulement quand les puissants la figent pour se perpétuer (Baroka), mais elle en constitue le socle vital quand le peuple la fait vivre pour s'adapter (Sidi, le mime de la danse du voyageur). \textbf{[Ouverture]} Cette tension rappelle le dilemme universel des sociétés en mutation, que l'on retrouve chez Chinua Achebe dans \emph{Le Monde s'effondre}, où la rigidité d'Okonkwo le perd face au changement, ou encore dans l'actualité camerounaise où la chefferie traditionnelle négocie sa place face à l'administration moderne. \textbf{[Chute]} \emph{Le progrès sans mémoire est un arbre sans racines ; la tradition sans avenir est un musée sans visiteurs.}
\tcblower
\textbf{Analyse} :\par
- \textbf{Bilan} : Réponse nuancée à la problématique (« pas intrinsèquement ennemie »), synthèse des 3 parties (Baroka / Peuple / Synthèse) sans numérotation.\par
- \textbf{Ouverture} : Pertinente (Achebe, même aire culturelle, même thème), concise, lien avec le Cameroun (contexte élève).\par
- \textbf{Chute} : Métaphore équilibrée (arbre/musée), rythme binaire, image forte.
\end{exoresolu}

\courssub{Étape 5 : Champ lexical, champ sémantique et texte descriptif}

\begin{definition}[Champ lexical vs Champ sémantique]
\begin{itemize}
\item \cle{Champ lexical} : Ensemble de mots (différents signifiants) liés par le \textbf{sens} et rattachés à une même \cle{notion} (thème). Ex : Champ lexical de la \emph{tradition} : \emph{coutume, rite, ancêtre, héritage, tabou, chefferie, dot, masque, fête}.
\item \cle{Champ sémantique} : Ensemble des \textbf{sens} (acceptions) d'un \textbf{même mot} (polysémie). Ex : \cle{Progrès} = 1. Technique (route, train) 2. Social (éducation, droits) 3. Moral (sagesse, humanité) 4. Économique (développement).
\end{itemize}
\textbf{Utilité en dissertation} : 1. Analyser une citation (champ lexical dominant = idéologie). 2. Varier le vocabulaire (synonymes, hyperonymes, antonymes) pour éviter les répétitions. 3. Préciser sa problématique (quel sens de « progrès » ?).
\end{definition}

\begin{methode}[Rédiger un paragraphe descriptif (argument d'illustration / autorité)]
La description n'est pas ornementale : elle \emph{prouve} par l'image. Structure \cle{« Point de vue + Détails significatifs + Verbes d'état/perception + Figures »}.\par
\textbf{1. Point de vue} : Qui regarde ? (Narrateur, personnage, regard interne/externe). \textbf{2. Détails significatifs} : Sélectionner 3-4 détails qui \emph{prouvent} l'argument (pas l'inventaire). \textbf{3. Verbes} : Verbes d'état (\emph{être, paraître, sembler}), de perception (\emph{voir, entendre, sentir}), d'action (peu). \textbf{4. Figures} : Comparaison, métaphore, personification pour l'effet de sens.
\end{methode}

\begin{exoresolu}{Analyser le champ lexical de la « modernité » chez Lakunle (Extrait Acte 1) + Paragraphe descriptif}
\textbf{Extrait} : Lakunle : « \emph{Je veux une épouse qui marche à mes côtés, pas derrière moi. Une femme éduquée, moderne, qui boit du thé à l'anglaise...} »\par
\textbf{1. Champ lexical de la modernité (Lakunle)} :\par
\emph{épouse, côtés, éduquée, moderne, thé, anglaise, machine, livre, école, progrès, civilisation, développement, confort, modernité}.\par
\textbf{2. Champ sémantique de « Progrès » chez Lakunle} : Réduit à \emph{occidentalisation / mimétisme} (connotation péjorative).\par
\textbf{3. Connotation} : \textbf{Péjorative / Caricaturale}. Lakunle associe modernité $\rightarrow$ imitation servile de l'Occident (\emph{thé à l'anglaise}), rupture brutale (\emph{pas derrière moi}), mépris du corps/terroir.\par
\textbf{4. Paragraphe argumentatif intégré (Partie I ou II)} :\par
Lakunle incarne une modernité \emph{par procuration}, vidée de son sens profond. Son vocabulaire trahit une aliénation culturelle : le champ lexical de l'\emph{anglicité} (\emph{thé, anglaise, machine}) supplante celui de la \emph{réalité yoruba}. \emph{« Je veux une épouse qui marche à mes côtés »}, proclame-t-il, mais il nie la dot, fondement du contrat social traditionnel. Cette citation révèle une \emph{connotation} négative : le « progrès » lakunlien est rejet de l'Autre, non construction \emph{avec} l'Autre. Soyinka, par l'ironie, dénonce ce mimétisme stérile qui, loin de faire progresser Ilujinle, ne fait que le déraciner.
\tcblower
\textbf{Validation} : Le paragraphe utilise l'analyse lexicale (champ lexical/sémantique, connotation) comme outil d'interprétation (étape I). La citation est intégrée et commentée.
\end{exoresolu}

\begin{exoresolu}{Écrire un paragraphe descriptif : La vitalité de la tradition (Mime du voyageur, Acte 1)}
\textbf{Consigne} : Décrire la scène pour montrer que la tradition est vivante, créatrice, collective (Argument Partie II).\par
\textbf{Proposition} :\par
\textbf{[Point de vue]} Regard interne : le spectateur est immergé dans la place du village. \textbf{[Détails]} La foule en cercle, le tambour qui « parle », les danseurs qui miment la voiture qui cale, le « blanc » perdu (casque pith, moustache, carte), les rires francs. \textbf{[Verbes]} La place \emph{bruisse}, le tam-tam \emph{bat}, les corps \emph{ondulent}, la foule \emph{éclate} de rire. \textbf{[Figures]} Métaphore : « le tambour \emph{parle} » (langage codé) ; Personification : la voiture \emph{tousse}, \emph{crache}, \emph{meurt} (technique fragile). \textbf{[Interprétation]} Ce n'est pas un folklore mort : la tradition \emph{digère} la modernité (la voiture) par le rire et le mime, elle l'apprivoise. Elle est intelligence collective face à l'inconnu.
\end{exoresolu}

\courssub{Étape 6 : Intégration — Rédiger une dissertation complète}

\begin{methode}[Rédiger une dissertation complète : la check-list « Bac Ready » (Auto-évaluation)]
Avant de rendre la copie, vérifier \textbf{10 points} :\par
1. \textbf{Sujet} : Recopié exact ? Compris ? Mots-clés soulignés ?\par
2. \textbf{Introduction} : Amorce $\rightarrow$ Sujet $\rightarrow$ Problématique $\rightarrow$ Plan (4 étapes distinctes, saut de ligne) ?\par
3. \textbf{Développement} : 2 ou 3 parties équilibrées ? 2-3 paragraphes par partie ? Titres de parties (marge ou texte) ?\par
4. \textbf{Paragraphes} : Structure A.E.C.I.T. respectée ? \textbf{Citation intégrée + Commentée} à chaque paragraphe ?\par
5. \textbf{Transitions} : Liens logiques explicites (mots de liaison + reprise lexicale) entre \emph{tous} les paragraphes ?\par
6. \textbf{Conclusion} : Bilan (réponse prob.) $\rightarrow$ Ouverture (pertinente) $\rightarrow$ Chute (forte) ? Saut de ligne avant ?\par
7. \textbf{Langue} : Phrases complètes, syntaxe soignée, vocabulaire précis (champs lexicaux), \textbf{zéro SMS / argot / familier}. Orthographe soignée (accords, homophones).\par
8. \textbf{Présentation} : Saut de ligne après intro, entre parties, avant conclusion. Paragraphes indentés (alinéa visible). Marge gauche respectée.\par
9. \textbf{Œuvre} : \emph{Le Lion et la perle} cité au moins 3-4 fois (différentes scènes/actes) ? Titres en italique.\par
10. \textbf{Temps} : Gestion du temps (30 min analyse/plan, 1h30 rédaction, 15 min relecture).
\end{methode}

\begin{exoresolu}{Sujet d'intégration : « Le théâtre est-il le lieu privilégié du débat d'idées ? » (Plan + Intro + 1 paragraphe type)}
\textbf{Analyse rapide} : \cle{Théâtre} = genre dialogué, public, représentation, incarnation. \cle{Débat d'idées} = confrontation thèses, visée didactique. \cle{Privilégié} = unique ? supérieur aux autres genres (roman, essai, poésie) ?\par
\textbf{Problématique} : Le théâtre, par sa nature dialoguée et spectaculaire, offre-t-il une efficacité supérieure aux autres formes littéraires pour mettre en scène et résoudre le conflit des idées ?\par
\textbf{Plan dialectique} :\par
\textbf{I. La spécificité théâtrale : un débat incarné et dynamique} (Dialogue/Conflit immédiat, Didactique par l'exemple/le corps, Public juge).\par
\textbf{II. Les limites du théâtre face à l'essai ou au roman} (Contraintes scéniques / temps, Manichéisme des personnages, Primat de l'émotion sur la raison).\par
\textbf{III. La complémentarité des genres : le théâtre comme mise à l'épreuve publique de la pensée} (Essayiste $\rightarrow$ Théâtre $\rightarrow$ Public / Roman $\rightarrow$ Intériorité).\par
\textbf{Introduction (rédigée)} :\par
\emph{« Le théâtre est une école de larmes et de rire », disait Musset. Mais est-il aussi une école de pensée ? De l'Antiquité grecque à la scène contemporaine, la pièce a toujours mis en présence des voix discordantes. Wole Soyinka, dans \emph{Le Lion et la perle}, transforme le village d'Ilujinle en agora où s'affrontent tradition et modernité. \textbf{Le théâtre, par sa nature dialoguée et spectaculaire, offre-t-il une efficacité supérieure aux autres formes littéraires pour mettre en scène et résoudre le conflit des idées ?} Nous verrons d'abord que le théâtre incarne le débat comme nulle part ailleurs, avant d'en analyser les contraintes intrinsèques, pour enfin montrer comment il s'articule aux autres genres pour former le citoyen. »}\par
\textbf{Paragraphe type (Partie I, argument 1 : L'incarnation)} :\par
\textbf{[A]} Le théâtre possède une force unique : il donne chair et voix aux idées abstraites. \textbf{[E]} Là où l'essai expose des concepts, la pièce les \emph{met en corps} : le conflit d'idées devient conflit de personnages, de regards, de gestes, sous les yeux du spectateur. \textbf{[C]} Dans \emph{Le Lion et la perle}, l'opposition Tradition/Progrès n'est pas un cours magistral : elle est le duel verbal et physique entre Baroka et Lakunle, notamment lors de la scène de la poste (Baroka qui détourne le géomètre) ou du mime de la voiture (Lakunle qui mime le conducteur). \textbf{[I]} Le spectateur \emph{voit} la ruse de Baroka (intelligence traditionnelle) et \emph{voit} la maladresse de Lakunle (modernité hors-sol). L'idée « la tradition a une intelligence pratique » devient une \emph{évidence sensible}, bien plus persuasive qu'un traité de sociologie. \textbf{[T]} \emph{Cependant, cette force de l'incarnation a son revers : la contrainte du temps scénique oblige à simplifier, voire à caricaturer, la complexité des idées.}
\tcblower
\textbf{Validation check-list} : Sujet recopié ? Oui. Problématique ouverte ? Oui. Plan dialectique ? Oui. Intro 4 temps ? Oui. Paragraphe A.E.C.I.T. ? Oui (Citation scène précise, Interprétation « évidence sensible », Transition vers limite). Ouverture en conclusion (à faire) : Roman (intériorité) / Essai (abstraction).
\end{exoresolu}

\begin{attention}[Les 5 erreurs qui coûtent des points au Probatoire]
\textbf{1. Introduction « catalogue » ou « vie de l'auteur »} — Pas de biographie de Soyinka, pas de résumé de la pièce. L'amorce doit \textbf{accrocher le sujet}, pas l'auteur. \textbf{Sanction} : -2 à -4 pts sur 20 (hors-sujet apparent).\par
\textbf{2. Problématique fermée ou absente} — « La tradition est-elle un obstacle ? » (réponse oui/non) au lieu de « Dans quelle mesure / Comment... ». Ou pire : pas de problématique explicite. \textbf{Sanction} : Plan incohérent, note plafond 8/20.\par
\textbf{3. Paragraphes sans citation / Citation « orpheline »} — Un paragraphe d'argumentation littéraire \textbf{doit} contenir une référence textuelle précise (citation ou renvoi scène/acte). Laisser une citation sans commentaire (pas d'étape I) vaut 0 point pour l'argument. \textbf{Sanction} : -1 pt par paragraphe concerné.\par
\textbf{4. Transitions inexistantes ou purement mécaniques} — « II. » écrit dans la marge sans phrase de liaison. Ou « Donc » utilisé à tort (voir méthode 4). \textbf{Sanction} : Perte de cohérence, note style/structure baissée.\par
\textbf{5. Conclusion « nouvelle idée » ou « question bête »} — Introduire un argument nouveau dans la conclusion (hors structure). Finir par « Qu'en pensez-vous ? » ou « C'est tout. ». \textbf{Sanction} : Dernière impression négative, perte de 1-2 pts sur l'ensemble.
\end{attention}

\begin{aretenir}[Synthèse des savoirs-clés de la leçon]
\begin{itemize}
\item \textbf{Dissertation} : Exercice codifié (Intro 4 temps / Dev A.E.C.I.T. / Conclu 3 temps). Exige rigueur, citations, transitions.
\item \textbf{Le Lion et la perle} : Œuvre pivot. Connaître 5-6 citations clés (Baroka/Lakunle/Sidi), 3 scènes majeures (Mime, Poste, Final), thèmes (Tradition/Modernité, Pouvoir/Ruse, Femme/Regard).
\item \textbf{Outils langue} : Types de phrase (Déclarative roi), Dénotation/Connotation (nuance), Énonciation (analyse citations), Champs lexicaux/sémantiques (vocabulaire précis, analyse style).
\item \textbf{Description} : Outil argumentatif (Point de vue + Détails + Verbes + Figures), pas ornement.
\item \textbf{Probatoire} : 3h. 30 min brouillon (analyse/plan), 1h30 rédaction, 15 min relecture. Propreté, alinéas, italique titres.
\end{itemize}
\end{aretenir}

\newpage

\coursec{Exercices}

\courssub{Je vérifie mes connaissances}

\begin{exercice}[QCM : les types de phrase]
Pour chaque phrase extraite ou inspirée du \emph{Lion et la perle}, choisis la bonne réponse.
\begin{enumerate}
\item « Lakunle, écoute-moi ! » est une phrase : a) déclarative ; b) impérative ; c) interrogative.
\item « Quel beau matin sur Ilujinle ! » est une phrase : a) exclamative ; b) impérative ; c) déclarative.
\item « Sidi porte l'eau sur la tête. » est une phrase : a) interrogative ; b) exclamative ; c) déclarative.
\item « Pourquoi refuses-tu le magazine ? » est une phrase : a) déclarative ; b) interrogative ; c) impérative.
\end{enumerate}
\end{exercice}

\begin{exercice}[Vrai ou faux, en justifiant]
Réponds par vrai ou faux, puis justifie en une phrase.
\begin{enumerate}
\item L'introduction d'une dissertation se compose de trois parties : l'amorce, la problématique et l'annonce du plan.
\item Une citation peut être insérée dans un paragraphe sans verbe introducteur.
\item La dénotation d'un mot est son sens chargé d'émotions et d'associations personnelles.
\item Une transition sert à relier deux parties du développement en rappelant la partie précédente et en annonçant la suivante.
\end{enumerate}
\end{exercice}

\begin{exercice}[Définitions]
Définis en deux ou trois phrases chacun des termes suivants : \cle{problématique}, \cle{champ lexical}, \cle{énonciation}, \cle{connotation}.
\end{exercice}

\begin{exercice}[Repérage rapide]
On donne les mots suivants : \emph{école, progrès, magazine, route, civilisation, photographie, modernité}.
\begin{enumerate}
\item Quel est le champ lexical commun à ces mots ?
\item Cite deux autres mots qui pourraient enrichir ce champ lexical.
\item Parmi ces mots, lequel possède une connotation péjorative dans la bouche de Baroka ? Explique.
\end{enumerate}
\end{exercice}

\courssub{Je m'entraîne}

\begin{exercice}[Analyser un sujet]
On donne le sujet : « Lakunle, un moderniste ridicule ou un visionnaire incompris ? »
\begin{enumerate}
\item Dégage les mots-clés du sujet.
\item Indique le thème général de la réflexion.
\item Formule la problématique sous la forme d'une question.
\item Propose un plan en deux parties (axes seulement).
\end{enumerate}
\end{exercice}

\begin{exercice}[Insérer des citations]
Complète le paragraphe suivant en insérant correctement les trois citations données, avec des verbes introducteurs variés (affirmer, déclarer, s'écrier, prétendre, ajouter...).

Citations : a) « Je ne veux pas d'un mari qui paie ma dot » ; b) « Le progrès viendra par l'école » ; c) « Le vieux lion a encore ses griffes ».

Paragraphe à compléter : Le conflit entre tradition et modernité traverse toute la pièce. Sidi refuse les arguments de Lakunle : \ldots\ldots\ldots\ldots. De son côté, le jeune instituteur croit à l'instruction : \ldots\ldots\ldots\ldots. Quant à Baroka, il incarne la ruse de la tradition : \ldots\ldots\ldots\ldots.
\end{exercice}

\begin{exercice}[Dénotation et connotation]
Dans l'extrait suivant, cinq mots sont soulignés. Pour chacun, donne le sens dénoté, puis le sens connoté dans le contexte de la pièce.

« Lakunle, ce \underline{fou} de l'\underline{école}, court après Sidi comme un \underline{enfant}. Il rêve d'une Ilujinle \underline{moderne}, débarrassée de ses vieilles \underline{chaînes}. »
\end{exercice}

\begin{exercice}[Rédiger une transition]
On rédige une dissertation sur le sujet : « La tradition est-elle un obstacle au bonheur des personnages du \emph{Lion et la perle} ? » La première partie montre les avantages de la tradition (solidarité, sagesse, identité) ; la deuxième partie en expose les limites (poids de la dot, autorité de Baroka, refus du changement). Rédige la transition entre ces deux parties en trois à cinq phrases.
\end{exercice}

\courssub{Je me prépare au Bac}

\begin{exercice}[Sujet de dissertation : introduction complète]
Sujet : « Le choix de Sidi à la fin du \emph{Lion et la perle} est-il une victoire de la tradition ? »
\begin{enumerate}
\item Analyse le sujet : mots-clés, thème, problématique.
\item Rédige une amorce (accroche) de deux ou trois phrases, en lien avec la situation de communication.
\item Présente l'œuvre et l'auteur (Wole Soyinka, dramaturge nigérian, prix Nobel de littérature 1986).
\item Formule la problématique.
\item Annonce clairement un plan en deux parties.
\item Rassemble le tout en une introduction complète et rédigée (10 à 15 lignes).
\end{enumerate}
\end{exercice}

\begin{exercice}[Langue : analyse d'un extrait]
On donne l'extrait suivant, adapté de la pièce :

« Lakunle : Sidi, ma belle, écoute-moi ! Ne porte plus cette calebasse sur ta tête. Je te le dis : l'avenir est à l'école, au progrès, aux machines. Comme tu es obstinée ! Veux-tu enfin me comprendre ? »

\begin{enumerate}
\item Relève dans cet extrait un exemple de chaque type de phrase obligatoire (déclarative, impérative, interrogative, exclamative) et justifie chaque classement.
\item Relève deux marques d'énonciation (pronoms personnels, interpellation, temps verbaux, modalisateurs) et explique ce qu'elles révèlent de la situation de communication.
\item Distingue le sens dénoté et le sens connoté des mots \emph{calebasse}, \emph{école} et \emph{obstinée}.
\item Construis le champ lexical de la modernité présent dans cet extrait, puis propose deux mots supplémentaires qui pourraient l'enrichir.
\end{enumerate}
\end{exercice}

\begin{exercice}[Dissertation complète et remédiation]
\textbf{Partie A.} Rédige une dissertation complète (introduction, développement en deux parties avec transitions, conclusion) sur le sujet : « Faut-il rejeter la tradition pour entrer dans la modernité ? » Tu t'appuieras sur le \emph{Lion et la perle} et sur tes connaissances personnelles. Chaque paragraphe du développement devra contenir une idée, un exemple et au moins une citation correctement insérée.

\textbf{Partie B.} Voici l'introduction défaillante d'un élève :

« Dans la vie, il y a la tradition et la modernité. Beaucoup de gens en parlent. Soyinka aussi en parle dans sa pièce. Nous allons voir cela. »

\begin{enumerate}
\item Releve trois défauts de cette introduction.
\item Réécris-la en respectant les étapes de l'introduction (amorce, présentation de l'œuvre, problématique, annonce du plan).
\end{enumerate}
\end{exercice}

\newpage

\coursec{Corrigés des exercices}

\coursec{Rédiger une dissertation : \emph{Le Lion et la perle}}

\courssub{Je vérifie mes connaissances}

\begin{corrige}[Exercice 1 — QCM : les types de phrase]
\begin{enumerate}
\item \textbf{b) impérative}. Le verbe \emph{écoute} est conjugué à l'impératif présent, sans sujet exprimé ; la phrase exprime un ordre, renforcé par le point d'exclamation.
\item \textbf{a) exclamative}. La phrase exprime un sentiment (l'admiration) ; elle est introduite par le déterminant exclamatif \emph{quel} et se termine par un point d'exclamation.
\item \textbf{c) déclarative}. La phrase affirme un fait, sans marque d'interrogation ni d'exclamation ; elle se termine par un point simple.
\item \textbf{b) interrogative}. La phrase pose une question : mot interrogatif \emph{pourquoi}, inversion du sujet (\emph{refuses-tu}) et point d'interrogation.
\end{enumerate}
\textbf{Point de vigilance :} ne pas confondre le \emph{point d'exclamation} (ponctuation) avec la phrase exclamative (type de phrase) : une phrase impérative peut aussi porter un point d'exclamation, comme à la question 1.
\end{corrige}

\begin{corrige}[Exercice 2 — Vrai ou faux, en justifiant]
\begin{enumerate}
\item \textbf{Vrai.} L'introduction d'une dissertation comprend bien trois étapes obligatoires : l'\cle{amorce} (ou accroche), la \cle{problématique} précédée de la présentation de l'œuvre, et l'\cle{annonce du plan}.
\item \textbf{Faux.} Une citation doit toujours être introduite par un \cle{verbe introducteur} (affirmer, déclarer, s'écrier...) qui l'intègre à la phrase et à l'argumentation ; une citation « jetée » sans introduction perd sa valeur argumentative.
\item \textbf{Faux.} Le sens chargé d'émotions et d'associations est la \cle{connotation} ; la \cle{dénotation} est le sens premier, objectif et stable du mot, celui que donne le dictionnaire.
\item \textbf{Vrai.} La \cle{transition} assure la liaison entre deux parties : elle \emph{bilan} la partie qui se termine et \emph{annonce} celle qui commence, ce qui donne de la cohérence au devoir.
\end{enumerate}
\end{corrige}

\begin{corrige}[Exercice 3 — Définitions]
\textbf{Problématique :} question centrale que pose le sujet de dissertation et à laquelle tout le devoir devra répondre. Elle se formule généralement sous la forme d'une phrase interrogative et oriente le plan.

\textbf{Champ lexical :} ensemble de mots qui se rapportent à une même notion ou à un même thème et qui apparaissent ensemble dans un texte (par exemple \emph{école, cahier, maître, leçon} forment le champ lexical de l'instruction). Il contribue à l'unité de sens du texte.

\textbf{Énonciation :} mise en scène de la communication dans un énoncé : qui parle (\cle{énonciateur}), à qui (\cle{destinataire}), où et quand. On la repère grâce aux pronoms personnels, aux temps verbaux, aux interpellations et aux modalisateurs.

\textbf{Connotation :} ensemble des valeurs affectives, culturelles ou imagées qu'un mot évoque en plus de son sens premier. Elle dépend du contexte : ainsi \emph{chaînes} connote l'oppression dans la bouche de Lakunle.
\end{corrige}

\begin{corrige}[Exercice 4 — Repérage rapide]
\begin{enumerate}
\item Ces mots appartiennent au \cle{champ lexical de la modernité} (ou du progrès) : ils renvoient tous aux transformations venues de l'extérieur du village d'Ilujinle.
\item Deux mots possibles pour enrichir ce champ : \emph{machine}, \emph{voiture} (on accepte aussi \emph{train}, \emph{téléphone}, \emph{électricité}, \emph{uniforme}...).
\item C'est \emph{civilisation} (et, dans une moindre mesure, \emph{progrès}) qui prend une connotation péjorative dans la bouche de Baroka : le Baale y voit une menace venue de l'Occident, un prétexte pour détruire les coutumes d'Ilujinle. Le mot, positif pour Lakunle, devient pour Baroka synonyme d'arrogance et de trahison des valeurs du village.
\end{enumerate}
\end{corrige}

\courssub{Je m'entraîne}

\begin{corrige}[Exercice 5 — Analyser un sujet]
\begin{enumerate}
\item \textbf{Mots-clés :} \emph{Lakunle} (personnage), \emph{moderniste} (thèse favorable), \emph{ridicule} (jugement négatif), \emph{visionnaire} (jugement positif), \emph{incompris} (rapport aux autres personnages).
\item \textbf{Thème général :} le conflit entre tradition et modernité, incarné par le personnage de Lakunle, et le regard que la pièce porte sur le « progressiste ».
\item \textbf{Problématique :} Lakunle est-il seulement un personnage comique que la pièce tourne en dérision, ou porte-t-il, malgré ses maladresses, un message valable sur l'avenir de la société africaine ?
\item \textbf{Plan en deux parties (dialectique) :}
\begin{itemize}
\item I. Lakunle, un moderniste ridicule (son snobisme, son rejet maladroit des coutumes, le comique de ses situations).
\item II. Lakunle, un visionnaire incompris (la justesse de certaines de ses idées : l'école, l'égalité, le refus de la dot ; l'aveuglement du village).
\end{itemize}
\end{enumerate}
\textbf{Point de vigilance :} le sujet est formulé comme une alternative (« ou ») ; le plan dialectique en deux parties est donc le plus adapté. La problématique doit reprendre les deux termes de l'alternative sans choisir d'avance.
\end{corrige}

\begin{corrige}[Exercice 6 — Insérer des citations]
Le conflit entre tradition et modernité traverse toute la pièce. Sidi refuse les arguments de Lakunle : elle \textbf{affirme} avec fierté : « Je ne veux pas d'un mari qui paie ma dot ». De son côté, le jeune instituteur croit à l'instruction : il \textbf{prétend} que « le progrès viendra par l'école ». Quant à Baroka, il incarne la ruse de la tradition : la rumeur \textbf{rapporte} que « le vieux lion a encore ses griffes ».

\textbf{Points de vigilance :}
\begin{itemize}
\item Varier les verbes introducteurs (\emph{affirmer}, \emph{prétendre}, \emph{rapporter}) et les accorder au sens : \emph{prétendre} marque ici la distance du narrateur envers Lakunle.
\item Respecter la ponctuation : deux-points avant la citation, guillemets, majuscule à la citation si elle forme une phrase complète.
\item Ne jamais laisser une citation seule : elle doit être introduite \emph{et} exploitée (commentée) ensuite.
\end{itemize}
\end{corrige}

\begin{corrige}[Exercice 7 — Dénotation et connotation]
\begin{center}
\begin{tabularx}{\linewidth}{|l|X|X|}
\hline
\rowcolor{popblueL} \textbf{Mot} & \textbf{Sens dénoté} & \textbf{Sens connoté (contexte)} \\
\hline
\emph{fou} & Personne atteinte de troubles mentaux. & Excentricité, comportement jugé absurde : Lakunle est moqué par le village. \\
\hline
\emph{école} & Établissement où l'on enseigne. & Modernité occidentale, rupture avec les coutumes ; source de conflit. \\
\hline
\emph{enfant} & Jeune être humain. & Immaturité, naïveté : Lakunle manque de sagesse et d'expérience. \\
\hline
\emph{moderne} & Qui appartient au temps présent, récent. & Séduisant pour Lakunle, suspect et menaçant pour les tenants de la tradition. \\
\hline
\emph{chaînes} & Anneaux de métal pour attacher. & Servitude, oppression : les coutumes vécues comme une prison par Lakunle. \\
\hline
\end{tabularx}
\end{center}
\textbf{Point de vigilance :} la connotation dépend toujours du contexte et du locuteur ; \emph{moderne} n'a pas la même valeur dans la bouche de Lakunle (éloge) et dans celle de Baroka (rejet).
\end{corrige}

\begin{corrige}[Exercice 8 — Rédiger une transition]
\textbf{Modèle de transition :}

Ainsi, la tradition apparaît d'abord comme une richesse pour Ilujinle : elle assure la solidarité entre les habitants, transmet la sagesse des anciens et donne à chacun une identité. Cependant, cette force protectrice peut se transformer en poids étouffant. Ne faut-il pas alors s'interroger sur ses limites ? C'est ce que nous allons montrer dans une seconde partie, en examinant le poids de la dot, l'autorité abusive de Baroka et le refus de tout changement.

\textbf{Méthode vérifiée :} la transition contient bien (1) un bilan de la partie I (\emph{solidarité, sagesse, identité}), (2) une charnière (\emph{cependant}) et une question-relais, (3) l'annonce de la partie II (\emph{dot, autorité de Baroka, refus du changement}).
\end{corrige}

\courssub{Je me prépare au Bac}

\begin{corrige}[Exercice 9 — Sujet de dissertation : introduction complète]
\textbf{1. Analyse du sujet.} Mots-clés : \emph{choix de Sidi}, \emph{fin de la pièce}, \emph{victoire}, \emph{tradition}. Thème : le sens du dénouement du \emph{Lion et la perle} et le rapport des personnages à la coutume. Problématique : en épousant Baroka plutôt que Lakunle, Sidi fait-elle triompher la tradition, ou son choix obéit-il à d'autres motifs ?

\textbf{2. Amorce.} Dans les sociétés africaines en pleine mutation, le mariage est souvent le lieu où s'affrontent coutumes ancestrales et aspirations nouvelles. Le choix d'une épouse devient alors un choix de société.

\textbf{3. Présentation de l'œuvre.} C'est tout l'enjeu du \emph{Lion et la perle}, pièce du dramaturge nigérian Wole Soyinka, prix Nobel de littérature en 1986, qui met en scène le village d'Ilujinle partagé entre le Baale Baroka et l'instituteur Lakunle.

\textbf{4. Problématique.} Le choix final de Sidi, qui épouse Baroka, signifie-t-il la victoire de la tradition sur la modernité ?

\textbf{5. Annonce du plan.} Nous verrons d'abord que ce choix consacre le triomphe de la tradition et de la ruse du « vieux lion » ; nous montrerons ensuite qu'il relève aussi d'une décision personnelle et ambiguë, qui ne constitue pas une victoire totale de la coutume.

\textbf{6. Introduction rédigée (modèle).}

Dans les sociétés africaines en pleine mutation, le mariage est souvent le lieu où s'affrontent coutumes ancestrales et aspirations nouvelles : choisir un époux, c'est parfois choisir un monde. C'est tout l'enjeu du \emph{Lion et la perle}, pièce du dramaturge nigérian Wole Soyinka, prix Nobel de littérature en 1986. L'œuvre met en scène le village d'Ilujinle, déchiré entre Baroka, le Baale gardien des coutumes, et Lakunle, le jeune instituteur chantre du progrès, tous deux prétendants de la belle Sidi. Or, au terme de la pièce, Sidi choisit d'épouser Baroka. Ce dénouement inattendu interroge : le choix de Sidi est-il une victoire de la tradition ? Nous verrons dans un premier temps que ce choix consacre le triomphe de la coutume et de la ruse du « vieux lion » ; nous nuancerons ensuite ce constat en montrant que Sidi agit aussi par décision personnelle, ce qui interdit de parler de victoire totale de la tradition.

\textbf{Barème indicatif (sur 5 points) :} amorce en lien avec le sujet (1 pt) ; présentation exacte de l'œuvre et de l'auteur (1 pt) ; problématique claire sous forme interrogative (1 pt) ; annonce d'un plan en deux parties (1 pt) ; correction de la langue et qualité de rédaction (1 pt).

\textbf{Points de vigilance :} ne pas raconter toute la pièce dans la présentation ; la problématique doit reprendre les termes du sujet ; l'annonce du plan ne doit pas détailler les arguments, seulement les axes.
\end{corrige}

\begin{corrige}[Exercice 10 — Langue : analyse d'un extrait]
\textbf{1. Les quatre types de phrase :}
\begin{itemize}
\item \textbf{Déclarative :} « l'avenir est à l'école, au progrès, aux machines » — elle affirme un fait, ponctuée par un point.
\item \textbf{Impérative :} « écoute-moi ! » (ou « Ne porte plus cette calebasse sur ta tête ») — verbe à l'impératif, sans sujet, exprimant un ordre.
\item \textbf{Interrogative :} « Veux-tu enfin me comprendre ? » — inversion du sujet et point d'interrogation.
\item \textbf{Exclamative :} « Comme tu es obstinée ! » — introduite par \emph{comme} exclamatif, elle exprime le reproche et l'agacement.
\end{itemize}

\textbf{2. Marques d'énonciation (deux attendues) :}
\begin{itemize}
\item Les pronoms de première et deuxième personnes (\emph{je}, \emph{te}, \emph{tu}) : situation de dialogue direct entre Lakunle (énonciateur) et Sidi (destinataire).
\item L'interpellation « Sidi, ma belle » : apostrophe au destinataire, qui révèle la volonté de séduction et de persuasion de Lakunle.
\item (On accepte aussi : le présent de l'indicatif, temps du dialogue ; les impératifs et l'exclamation, modalisateurs qui trahissent l'autoritarisme et l'énervement du locuteur.)
\end{itemize}

\textbf{3. Dénotation et connotation :}
\begin{itemize}
\item \emph{Calebasse} : dénotation = récipient végétal pour porter l'eau ; connotation = poids de la tradition, tâche jugée rétrograde par Lakunle.
\item \emph{École} : dénotation = lieu d'enseignement ; connotation = salut, avenir, modernité (valeur très positive pour Lakunle).
\item \emph{Obstinée} : dénotation = qui persiste dans son refus ; connotation péjorative = entêtement irrationnel, reproche adressé à Sidi.
\end{itemize}

\textbf{4. Champ lexical de la modernité :} \emph{école, progrès, machines, avenir}. Deux mots pour l'enrichir : \emph{route}, \emph{civilisation} (ou \emph{magazine}, \emph{photographie}).

\textbf{Barème indicatif (sur 8 points) :} 4 types de phrase relevés et justifiés (4 $\times$ 0,5 = 2 pts) ; deux marques d'énonciation expliquées (2 pts) ; dénotation/connotation des trois mots (2 pts) ; champ lexical constitué et enrichi (2 pts).

\textbf{Points de vigilance :} justifier chaque classement par un indice grammatical (mode, ponctuation, inversion) et non par le seul « sens » ; distinguer clairement dénotation (dictionnaire) et connotation (contexte de la pièce).
\end{corrige}

\begin{corrige}[Exercice 11 — Dissertation complète et remédiation]
\textbf{Partie A — Éléments de correction (plan détaillé et passages rédigés attendus).}

\textbf{Introduction :} amorce sur la confrontation tradition/modernité dans l'Afrique contemporaine ; présentation du \emph{Lion et la perle} de Wole Soyinka ; problématique : faut-il nécessairement rejeter la tradition pour accéder à la modernité ? ; annonce du plan.

\textbf{I. La tradition peut apparaître comme un obstacle à la modernité.}
\begin{itemize}
\item Paragraphe 1 : le poids des coutumes. Exemple : la dot, que Lakunle refuse de payer ; citation : Sidi affirme « Je ne veux pas d'un mari qui paie ma dot », ce qui montre que la coutume enferme les individus.
\item Paragraphe 2 : l'autorité traditionnelle et la ruse. Exemple : Baroka fait circuler la rumeur de son impuissance pour piéger Sidi ; la tradition y devient instrument de domination.
\end{itemize}

\textbf{Transition :} bilan (la tradition freine) + charnière (\emph{cependant, tout rejeter serait une erreur}) + annonce de la partie II.

\textbf{II. La modernité ne peut s'épanouir qu'en s'appuyant sur la tradition.}
\begin{itemize}
\item Paragraphe 1 : les valeurs positives de la tradition (solidarité villageoise, sagesse, identité culturelle) ; Soyinka lui-même, moderniste, écrit en puisant dans la culture yoruba.
\item Paragraphe 2 : les dérives d'une modernité superficielle ; exemple : Lakunle, ridicule dans son rejet total des coutumes, montre qu'une modernité copiée de l'extérieur reste un vernis. Connaissance personnelle admise : écoles, routes et téléphones au Cameroun transforment la vie sans abolir les familles ni les chefferies.
\end{itemize}

\textbf{Conclusion :} réponse à la problématique (non, il ne faut pas rejeter la tradition, mais la dépasser en la conjuguant avec le progrès) ; ouverture possible : la question de l'éducation des jeunes filles, ou le rôle des écrivains africains dans ce dialogue des cultures.

\textbf{Barème indicatif (sur 20 points) :} introduction complète en trois étapes (4 pts) ; deux parties équilibrées avec transition (4 pts) ; paragraphes construits (idée + exemple + citation insérée avec verbe introducteur) (6 pts) ; conclusion avec réponse et ouverture (3 pts) ; correction de la langue (3 pts).

\textbf{Partie B — Remédiation.}

\textbf{1. Trois défauts de l'introduction de l'élève :}
\begin{itemize}
\item Amorce vide et banale (« Beaucoup de gens en parlent »), sans lien précis avec le sujet.
\item Présentation de l'œuvre inexistante : ni le titre (\emph{Le Lion et la perle}), ni la nationalité de l'auteur, ni le thème de la pièce ne sont donnés.
\item Absence de problématique et d'annonce du plan : « Nous allons voir cela » ne formule ni question ni axes.
\end{itemize}

\textbf{2. Réécriture (modèle).}

Toute société qui se modernise se trouve confrontée à un dilemme : faut-il abandonner ses coutumes pour entrer dans le monde nouveau ? C'est la question que pose le \emph{Lion et la perle}, pièce du dramaturge nigérian Wole Soyinka, prix Nobel de littérature en 1986, qui oppose dans le village d'Ilujinle le Baale Baroka, gardien de la tradition, à Lakunle, jeune instituteur passionné de progrès. Faut-il alors rejeter la tradition pour entrer dans la modernité ? Nous montrerons d'abord que la tradition peut freiner l'évolution des individus, puis qu'une modernité réussie doit au contraire s'enraciner dans les valeurs héritées.

\textbf{Points de vigilance :} une introduction sans problématique est sanctionnée à l'examen ; chaque citation du développement doit être introduite par un verbe et commentée ; la conclusion doit répondre explicitement à la question posée, sans introduire d'argument nouveau.
\end{corrige}

\newpage

\coursec{Fiche bilan}

\begin{popfigure}
\begin{tikzpicture}[mindmap, grow cyclic, every node/.style={concept, font=\small}, text=white, root concept/.append style={concept color=popdark, font=\bfseries}, level 1/.append style={level distance=3.4cm, sibling angle=60}, level 1 concept/.append style={font=\footnotesize}]
\node[root concept]{Rédiger une dissertation}
  child[concept color=popblue] { node {Analyser le sujet : mots-clés, type de sujet, problématique, plan} }
  child[concept color=poporange] { node {Introduction : amorce, sujet, problématique, annonce du plan} }
  child[concept color=popgreen] { node {Développement : 3 parties, paragraphes argumentés, citations, transitions} }
  child[concept color=poppurple] { node {Conclusion : bilan, réponse à la problématique, ouverture} }
  child[concept color=poppink] { node {Langue : types de phrases, énonciation, texte descriptif} }
  child[concept color=popgold] { node {Lexique : dénotation, connotation, champ lexical et sémantique} };
\end{tikzpicture}
\legende{Carte mentale de la leçon : les six compétences clés pour réussir la dissertation au Probatoire technique.}
\end{popfigure}

\begin{bilan}
\textbf{1. Définitions clés}
\begin{itemize}
  \item \cle{Dissertation} : exercice écrit qui répond de manière argumentée et organisée à une question (problématique) posée par un sujet.
  \item \cle{Problématique} : question centrale, issue de l'analyse du sujet, à laquelle toute la dissertation répond.
  \item \cle{Amorce} (ou accroche) : première phrase de l'introduction, générale, qui met le lecteur dans le sujet.
  \item \cle{Transition} : phrase ou paragraphe qui fait le bilan d'une partie et annonce la suivante.
  \item \cle{Dénotation} : sens premier, objectif et stable d'un mot (celui du dictionnaire).
  \item \cle{Connotation} : sens second, subjectif, lié aux idées et aux émotions qu'évoque un mot.
  \item \cle{Champ lexical} : ensemble de mots qui se rapportent à une même notion (ex. champ lexical de la guerre : combat, arme, soldat).
  \item \cle{Champ sémantique} : ensemble des sens (sèmes) qu'un mot peut prendre selon les contextes.
  \item \cle{Énonciation} : situation dans laquelle on produit un énoncé (qui parle ? à qui ? quand ? où ?).
\end{itemize}

\textbf{2. Les quatre types de phrases obligatoires}
\begin{center}
\begin{tabularx}{\linewidth}{|l|X|X|}
\hline
\rowcolor{popblueL} \textbf{Type} & \textbf{Fonction} & \textbf{Marques} \\
\hline
Déclarative & informer, constater & point final \\
\hline
Interrogative & poser une question & point d'interrogation, inversion, mot interrogatif \\
\hline
Impérative & ordonner, conseiller & verbe à l'impératif, point d'exclamation possible \\
\hline
Exclamative & exprimer un sentiment & point d'exclamation, « comme », « que » \\
\hline
\end{tabularx}
\end{center}

\textbf{3. Méthodes express}
\begin{itemize}
  \item \textbf{Analyser le sujet} : (1) lire plusieurs fois ; (2) souligner les mots-clés ; (3) définir ces mots ; (4) repérer le type de sujet (question directe, citation, affirmation à discuter) ; (5) formuler la problématique ; (6) choisir le plan (thématique ou thèse/antithèse/synthèse).
  \item \textbf{Introduction en 4 étapes} : amorce $\rightarrow$ présentation du sujet $\rightarrow$ problématique $\rightarrow$ annonce du plan.
  \item \textbf{Paragraphe argumentatif} : idée (argument) $\rightarrow$ explication $\rightarrow$ exemple ou citation introduite (« Comme le dit... ») $\rightarrow$ mini-conclusion.
  \item \textbf{Conclusion en 3 étapes} : bilan des idées $\rightarrow$ réponse claire à la problématique $\rightarrow$ ouverture.
  \item \textbf{Texte descriptif} : cadre spatio-temporel, expansion du nom, champs lexicaux, temps simples (imparfait, présent), rythme et sonorités.
\end{itemize}
\end{bilan}

\begin{aretenir}[Checklist avant le Bac]
Avant de rendre ma copie, je vérifie :
\begin{itemize}
  \item $\square$ J'ai souligné et défini les mots-clés du sujet au brouillon.
  \item $\square$ Ma problématique est une vraie question, liée au sujet.
  \item $\square$ Mon introduction contient les 4 étapes (amorce, sujet, problématique, plan).
  \item $\square$ Mon développement comporte 3 parties annoncées dans l'introduction.
  \item $\square$ Chaque paragraphe contient une idée, une explication et un exemple ou une citation.
  \item $\square$ Mes citations sont introduites et entre guillemets.
  \item $\square$ J'ai rédigé une transition entre chaque partie.
  \item $\square$ Ma conclusion fait le bilan, répond à la problématique et s'ouvre sur une autre question.
  \item $\square$ Je n'ai pas recopié le sujet ni écrit « je » de façon excessive.
  \item $\square$ J'ai varié les types de phrases et employé un vocabulaire précis (champs lexicaux).
  \item $\square$ Je me suis relu : orthographe, accords, ponctuation.
\end{itemize}
\end{aretenir}

\findecours

\end{document}
